% Libretto dei canti della parrocchia (A5, solo testo)
% Compilazione: pdflatex -> songidx -> pdflatex  (ricetta "libretto canti" in VS Code)
% Per la stampa a libretto su A4 compilare poi stampa.tex
\documentclass[11pt]{article}
\usepackage[T1]{fontenc}
\usepackage[italian]{babel}
\usepackage{tgpagella}
\usepackage{tgheros}            % sans-serif usato da songs per numeri e indice
\usepackage[a5paper,top=1.4cm,bottom=1.6cm,left=1.5cm,right=1.5cm]{geometry}
\usepackage[lyric]{songs}
\usepackage[hidelinks]{hyperref}
% I canti da ricontrollare sono elencati in DA_VERIFICARE.md

% ---- Dati del frontespizio (da personalizzare) -------------------------
\newcommand{\parrocchia}{Parrocchia di \dots}
\newcommand{\titololibretto}{Canti per la liturgia}
\newcommand{\edizione}{Edizione 2026}

% ---- Impostazioni del pacchetto songs ----------------------------------
\songcolumns{1}                 % una colonna: in A5 due sarebbero troppo strette
\songpos{0}                     % i canti possono spezzarsi tra due pagine (niente buchi)
\newindex{titoli}{titoli}       % indice alfabetico dei titoli (titoli.sxd -> titoli.sbx)
\indexsongsas{titoli}{\thepage} % nell'indice: pagina d'inizio del canto, non il suo numero
\renewcommand{\lyricfont}{\normalfont\large}
\renewcommand{\stitlefont}{\rmfamily\bfseries\Large}

% ---- Indice alfabetico: personalizzazioni --------------------------------
% (copia di \SB@maketitleindex del pacchetto songs, con due modifiche)
% 1) i titoli che iniziano con un numero vanno sotto "#" invece che sotto la cifra;
% 2) i canti omonimi vanno su righe separate (una riga per canto) invece che
%    su una riga sola con le pagine separate da virgola.
% 3) \soloindice{ (testo)} nel titolo di un canto: il testo compare solo
%    nell'indice alfabetico (serve a distinguere i canti omonimi), non nel titolo.
\newif\ifnellindice
\DeclareRobustCommand\soloindice[1]{\ifnellindice#1\fi}
\pdfstringdefDisableCommands{\def\soloindice#1{}}
\makeatletter
\newcommand\idxlettera[1]{%
  \ifnum0\ifnum`#1>47 \ifnum`#1<58 1\fi\fi>0 \#\else#1\fi}
\def\SB@unarigaperpagina#1#2\\#3\SB@fine{%
  \SB@ellipspread{\idxtitlefont\relax\ignorespaces#1\unskip}%
                 {{\idxrefsfont\relax#2}}%
  \ifx\relax#3\relax\else\SB@unarigaperpagina{#1}#3\SB@fine\fi}
\renewcommand\SB@maketitleindex{%
  \ifnum\idxheadwidth>\z@%
    \renewenvironment{SB@lgidx}[1]{
      \hbox{\SB@colorbox\idxbgcolor{\vbox{%
        \hbox to\idxheadwidth{{\idxheadfont\relax\idxlettera{##1}}\hfil}%
      }}}%
      \nobreak\vskip3\p@\@plus2\p@\@minus2\p@\nointerlineskip%
    }{\penalty-50\vskip5\p@\@plus5\p@\@minus4\p@}%
  \else%
    \renewenvironment{SB@lgidx}[1]{}{}%
  \fi%
  \renewenvironment{SB@smidx}[1]{}{}%
  \renewcommand\idxentry[2]{\SB@unarigaperpagina{##1}##2\\\SB@fine}%
  \renewcommand\idxaltentry[2]{%
    \SB@ellipspread{\idxlyricfont\relax\ignorespaces##1\unskip}%
                   {{\idxrefsfont\relax##2}}%
  }%
  \SB@displayindex%
}
\makeatother

% Intestazione di sezione per momento liturgico (va anche nell'indice iniziale)
\newcommand{\momento}[1]{%
  \par\bigskip
  \phantomsection\addcontentsline{toc}{section}{#1}%
  {\centering\scshape\LARGE #1\par}%
  \smallskip\hrule\medskip}

\pagestyle{plain}

\begin{document}

% ---- Frontespizio ------------------------------------------------------
\begin{titlepage}
  \centering
  \vspace*{3cm}
  {\scshape\Large \parrocchia\par}
  \vspace{2cm}
  {\bfseries\Huge \titololibretto\par}
  \vfill
  {\large \edizione\par}
\end{titlepage}

% ---- Indice per momento liturgico --------------------------------------
\renewcommand{\contentsname}{Indice}
\tableofcontents
\clearpage

% ---- Canti (un file per momento liturgico) -----------------------------
% Le sezioni ancora vuote hanno il modello commentato dentro il file.
% !TEX root = ../libretto.tex
% Sezione: Ingresso
\momento{Ingresso}
\begin{songs}{titoli}

% Fonte: LibrettoAccordi.pdf (Canti Coro), A4.
% Nell'elenco: Ingresso, Ordinario.
\beginsong{Acclamate al Signore}[by={M. Frisina}]
\beginchorus
Acclamate al Signore,
voi tutti della terra,
e servitelo con gioia,
andate a lui con esultanza.
Acclamate voi tutti
al Signore.
\endchorus
\beginverse
Riconoscete che il Signore, che il Signore è Dio,
Egli ci ha fatti, siamo suoi, suo popolo e gregge del suo pascolo.
\endverse
\beginverse
Entrate nelle sue porte con degli inni di grazia,
suoi atri nella lode, benedite, lodate il suo nome.
\endverse
\beginverse
Poiché buono è il Signore, eterna la sua misericordia,
la sua fedeltà si estende sopra ogni generazione.
\endverse
\endsong

% Fonte: LibrettoAccordi.pdf (Canti Coro), J42.
% Nell'elenco: Ingresso.
\beginsong{Beati voi}[by={F. Buttazzo}]
\beginchorus
Beati voi! Beati voi! Beati voi! Beati voi!
\endchorus
\beginverse
Se sarete poveri nel cuore, beati voi:
sarà vostro il Regno di Dio Padre.
Se sarete voi che piangerete, beati voi,
perché un giorno vi consolerò.
\endverse
\beginverse
Se sarete miti verso tutti, beati voi:
erediterete tutto il mondo.
Quando avrete fame di giustizia, beati voi,
perché un giorno Io vi sazierò.
\endverse
\beginverse
Se sarete misericordiosi, beati voi:
la misericordia troverete.
Se sarete puri dentro il cuore, beati voi,
perché voi vedrete il Padre mio.
\endverse
\beginverse
Se lavorerete per la pace, beati voi:
chiameranno voi «figli di Dio».
Se per causa mia voi soffrirete, beati voi,
sarà grande in voi la santità.
\endverse
\endsong

% Fonte: LibrettoAccordi.pdf (Canti Coro), A9.
% Nell'elenco: Ingresso.
\beginsong{Cantate al Signore}[by={Gen Rosso}]
\beginchorus
Alleluja! Alleluja! Alleluja! Alleluja! \rep{2}
\endchorus
\beginverse
Cantate al Signore un cantico nuovo, tutta la terra canti al Signore,
per tutta la terra si è accesa una luce: uomini nuovi cantano in cor
un cantico nuovo di gioia infinita, un canto d'amore a Dio fra noi.
Alleluja!
\endverse
\beginverse
Un coro di voci s'innalza al Signore, re della vita, luce del mondo;
discende dal cielo un fuoco d'amore: il Paradiso canta con noi
un cantico nuovo di gioia infinita, un canto d'amore a Dio fra noi.
Alleluja!
\endverse
\endsong

% Fonte: testo fornito dall'utente (tolti gli accordi). Non presente nell'elenco.
% Nel foglio il ritornello è scritto due volte di seguito; qui una volta sola. Finale: ultimo verso ripetuto.
\beginsong{Cantiamo al Signore con gioia}[by={F. Volpintesta}]
\beginchorus
Cantiamo al Signore con gioia,
cantiamo per amor:
un inno di lode la vita sarà,
un canto d'amore per lui.
\endchorus
\beginverse
In ogni momento, un dono,
la vita per sempre sarà,
e nasce la gioia che viene da te:
la gioia che mai finirà.
\endverse
\beginverse
La vita sarà solo amore
se tu resti ancora con noi:
la tua presenza la luce darà,
la luce che viene da te.
\endverse
\beginchorus
Cantiamo al Signore con gioia,
cantiamo per amor:
un inno di lode la vita sarà,
un canto d'amore per lui. \rep{2}
\endchorus
\endsong

% Fonte: LibrettoAccordi.pdf (Canti Coro), A2.
% Nell'elenco: Ingresso, Ordinario.
\beginsong{Cantiamo con gioia}[by={RnS}]
\beginchorus
Cantiamo con gioia al nostro Dio,
gridiamo con forza al nostro Re!
Le mani battiamo, insieme danziamo, cantiamo per la felicità!
\endchorus
\beginverse
Rialza chiunque è caduto, rafforza chi è debole tra noi.
Chi è triste ritrova in Lui la gioia.
Su canta insieme a noi e grida più che puoi
che Gesù Cristo è il Signor!
\endverse
\beginverse
Eterna è la sua misericordia, fedele per sempre è il suo amor.
Con forza ha sconfitto il nemico.
Su canta insieme a noi e grida più che puoi
che Gesù Cristo è il Signor!
\endverse
\endsong

% Testo fornito dall'utente (fonte: genverde.it/en/?p=34261).
% Versione italiana di "You Call Us by Name".
% Struttura cantata: Rit. - Strofa 1 - Rit. - Strofa 2 - Rit. + finale.
% Nell'elenco: Ingresso, Ordinario.
\beginsong{Chiamati per nome}[by={Gen Verde (J. Bai, N. L. Uelmen, S. G. McAllister)}]
\beginchorus
Veniamo da te,
chiamati per nome.
Che festa, Signore, tu cammini con noi.
Ci parli di te,
per noi spezzi il pane,
ti riconosciamo e il cuore arde, sei tu.
E noi tuo popolo
siamo qui.
\endchorus
\beginverse
Siamo come terra ed argilla
e la tua parola ci plasmerà;
brace pronta per la scintilla
e il tuo spirito soffierà,
c'infiammerà.
\endverse
\beginverse
Siamo come semi nel solco,
come vigna che il suo frutto darà;
grano del Signore risorto,
la tua messe che fiorirà d'eternità.
\endverse
\beginverse*
\textit{Finale (dopo l'ultimo ritornello):}
E noi tuo popolo
siamo qui, siamo qui.
\endverse
\endsong

% Fonte: LibrettoAccordi.pdf (Canti Coro), A19.
% Nell'elenco: Ingresso, Matrimonio (Baggio-Buttazzo). CD: Dall'amore di Dio (2010).
\beginsong{Dall'amore di Dio}[by={F. Buttazzo (CD Dall'amore di Dio)}]
\beginverse
Nasce dall'amore di Dio il miracolo di questa vita.
Vive per l'amore di Dio ogni angolo dell'universo.
Ma nel cuore d'ogni uomo risplende
il dono più divino: la vita sua più vera.
\endverse
\beginchorus
Creati come te, amati noi da te,
voliamo più su, più in alto di ogni cosa.
Creati come te, chiamati noi da te,
doniamo la vita con te.
Creati come te, amati noi da te,
voliamo più su, più in alto di ogni cosa.
Creati come te, chiamati noi da te,
amiamo per l'eternità.
\endchorus
\beginverse
Nasce dentro il cuore di Dio il sentiero della nostra vita.
Viene dall'amore di Dio la bellezza della nostra storia.
Noi, chiamati a donarci la vita,
seguiamo in due una strada, vediamo in lui la meta.
\endverse
\endsong

% Fonte: LibrettoAccordi.pdf (Canti Coro), A5.
% Nell'elenco: Ingresso, Ordinario.
\beginsong{Dall'aurora al tramonto}[by={M. Balduzzi (CD Verbum Panis)}]
\beginchorus
Dall'aurora io cerco te, fino al tramonto ti chiamo
ha sete solo di te l'anima mia come terra deserta.
\endchorus
\beginverse
Non mi fermerò un solo istante sempre canterò la tua lode
perché sei il mio Dio, il mio riparo
mi proteggerai all'ombra delle tue ali.
\endverse
\beginverse
Non mi fermerò un solo istante io racconterò le tue opere
perché sei il mio Dio, unico bene
nulla mai potrà la notte contro di me.
\endverse
\endsong

% Fonte: LibrettoAccordi.pdf (Canti Coro), A1.
% Nell'elenco: Ingresso, Ordinario.
\beginsong{È bello lodarti}[by={Gen Verde}]
\beginchorus
È bello cantare il tuo amore, è bello lodare il tuo nome.
È bello cantare il tuo amore, è bello lodarti Signore,
è bello cantare a te. \rep{2}
\endchorus
\beginverse
Tu che sei l'Amore infinito, che neppure il cielo può contenere,
ti sei fatto uomo, tu sei venuto qui
ad abitare in mezzo a noi, allora
\endverse
\beginverse
Tu, che conti tutte le stelle e le chiami una ad una per nome,
da mille sentieri ci hai radunati qui,
ci hai chiamati figli tuoi, allora
\endverse
\endsong

% Fonte: LibrettoAccordi.pdf (Canti Coro), O7.
% Nell'elenco: Ingresso.
\beginsong{È giunta l'ora}[by={L. Scaglianti}]
\beginverse
È giunta l'ora, Padre, per me: ai miei amici ho detto che
questa è la vita: conoscere te e il Figlio tuo, Cristo Gesù.
\endverse
\beginverse
Erano tuoi, li hai dati a me ed ora sanno che torno a te.
Hanno creduto: conservali tu nel tuo amore, nell'unità.
\endverse
\beginverse
Tu mi hai mandato ai figli tuoi: la tua parola è verità.
E il loro cuore sia pieno di gioia: la gioia vera viene da te.
\endverse
\beginverse
Io sono in loro e tu in me: che sian perfetti nell'unità
e il mondo creda che tu mi hai mandato, li hai amati come ami me.
\endverse
\endsong

% Fonte: LibrettoAccordi.pdf (Canti Coro), K12.
% Nell'elenco: Ingresso, Ordinario.
\beginsong{È la gioia che fa cantare}[by={RnS}]
\beginchorus
È la gioia che fa cantare, celebrando il Signore
il suo spirito oggi canta in me.
È la gioia che fa cantare, celebrando il Signore
il suo spirito oggi canta in me.
\endchorus
\beginverse
Io canto la gloria tua, perchè hai vinto la morte
Mia potente salvezza, mia forza sei tu.
\endverse
\beginverse
Tu raduni il tuo popolo e sconfiggi le tenebre.
Il tuo esercito siamo noi: «Vittoria di Dio».
\endverse
\endsong

% Fonte: LibrettoAccordi.pdf (Canti Coro), J28.
% Nell'elenco: Ingresso, Ordinario.
\beginsong{Eccomi\soloindice{ (Gen Verde)}}[by={Gen Verde}]
\beginchorus
Eccomi Gesù anche oggi davanti a te,
tutto rinnovato, proprio come Tu mi vuoi,
io sarò la risposta al tuo perché
un frutto degno del tuo abbandono, eccomi. \rep{2}
\endchorus
\beginverse
Un patto ormai ci stringe tutti qui di fronte a te
per dichiararti il nostro amore esclusivo
e corrispondere ai doni del tuo amore
perché sei stato abbandonato, siamo morti e risorti con te.
\endverse
\beginverse
Schierati in prima fila nella divina avventura,
per far conoscere in questo mondo l'amore,
un solo cuore, un solo coro che
ripete una preghiera con la vita dice a Te.
\endverse
\endsong

% Fonte: "99 - Raccolte e libretti accordi/frisina-marco-non-temere-fascicolo-completo_compress.pdf" (pag. 15). Salmo 39.
% Nell'elenco: Ingresso, Ordinario. CD: Non temere, Signore è il suo nome.
\beginsong{Eccomi\soloindice{ (M. Frisina)}}[by={M. Frisina (CD Non temere, Signore è il suo nome)}]
\beginchorus
Eccomi, eccomi,
Signore, io vengo.
Eccomi, eccomi!
Si compia in me la tua volontà.
\endchorus
\beginverse
Nel mio Signore ho sperato
e su di me s'è chinato,
ha dato ascolto al mio grido,
m'ha liberato dalla morte.
\endverse
\beginverse
I miei piedi ha reso saldi,
sicuri ha reso i miei passi.
Ha messo sulla mia bocca
un nuovo canto di lode.
\endverse
\beginverse
Il sacrificio non gradisci,
ma m'hai aperto l'orecchio,
non hai voluto olocausti,
allora ho detto: «Io vengo!».
\endverse
\beginverse
Sul tuo libro di me è scritto:
si compia il tuo volere.
Questo, mio Dio, desidero,
la tua legge è nel mio cuore.
\endverse
\beginverse
La tua giustizia ho proclamato,
non tengo chiuse le labbra.
Non rifiutarmi, Signore,
la tua misericordia.
\endverse
\endsong

% Fonte: LibrettoAccordi.pdf (Canti Coro), A8.
% Nell'elenco: Ingresso, Ordinario.
\beginsong{Giorno di festa}[by={C. Bertoni, F. Righetto}]
\beginchorus
È giorno di festa torniamo all'Amore
è giorno di gioia, è Dio fra noi.
Nella tua casa ogni attimo è pace, incontrare Te,
è dono e vita per me.
\endchorus
\beginverse
Come un fiume torna al mare, come un fiore guarda al cielo,
i nostri occhi, le nostre mani,
si alzano a Te, ritrovano Te, che sei Amore.
\endverse
\beginverse
Riconoscerti presente, Tu parola, Tu in quel pane.
Ardente fuoco, brucia nei cuori
che cercano Te, ritrovano Te, che sei Amore.
\endverse
\endsong

% Fonte: LibrettoAccordi.pdf (Canti Coro), D29.
% Nell'elenco: Ingresso, Ordinario.
\beginsong{I cieli narrano}[by={M. Frisina}]
\beginchorus
I cieli narrano la gloria di Dio
e il firmamento annunzia l'opera Sua.
Alleluia, alleluia, alleluia, alleluia.
\endchorus
\beginverse
Il giorno al giorno ne affida il messaggio,
la notte alla notte ne trasmette notizia,
non è linguaggio, non sono parole
di cui non si oda il suono.
\endverse
\beginverse
Là pose una tenda per il sole che sorge,
è come uno sposo dalla stanza nuziale,
esulta come un prode che corre
con gioia la sua strada.
\endverse
\beginverse
Lui sorge dall'ultimo estremo del cielo
e la sua corsa l'altro estremo raggiunge.
Nessuna delle creature potrà
mai sottrarsi al suo calore.
\endverse
\beginverse
La legge di Dio rinfranca l'anima mia,
la testimonianza del Signore è verace,
gioisce il cuore ai suoi giusti precetti
che danno la luce agli occhi.
\endverse
\endsong

% Fonte: LibrettoAccordi.pdf (Canti Coro), A10.
% Nell'elenco: Ingresso.
\beginsong{Lodate il Signore dai cieli}[by={M. Frisina}]
\beginchorus
Lodate il signore dai cieli, nell'alto dei cieli lodatelo.
Lodatelo, voi suoi angeli, lodatelo, voi sue schiere.
\endchorus
\beginverse
Lodatelo, sole e luna, lodatelo, fulgide stelle.
Lodatelo, cieli dei cieli, voi acque al di sopra dei cieli.
\endverse
\beginverse
Lodate il Signore, creature, voi mostri marini e abissi,
il fuoco, la neve, la nebbia, il vento che a lui obbedisce.
\endverse
\beginverse
Voi monti e voi tutte colline, voi alberi e tutti voi cedri,
voi fiere e tutte le bestie, voi rettili e uccelli alati.
\endverse
\beginverse
I re della terra e i popoli, i giudici e i governanti,
i giovani e le fanciulle, i vecchi insieme ai bambini.
\endverse
\beginverse
Perché il suo nome è sublime, la gloria sua splende sul mondo.
È lode per i suoi fedeli, il popolo che egli ama.
\endverse
\endsong

% Fonte: testo fornito dall'utente. Non presente nell'elenco.
\beginsong{Open the eyes of my heart}[by={M. W. Smith}]
\beginverse*
Open the eyes of my heart, Lord,
open the eyes of my heart:
I want to see you,
I want to see you.
\endverse
\beginverse*
To see you high and lifted up,
shinin' in the light of your glory.
Pour out your power and love
as we sing holy, holy, holy.
\endverse
\beginverse*
Holy, holy, holy,
we cry holy, holy, holy,
you are holy, holy, holy.
I want to see you.
\endverse
\endsong

% Fonte: LibrettoAccordi.pdf (Canti Coro), J33.
% Nell'elenco: senza tipologia (nel libretto accordi tra i canti di comunione). L'utente lo mette in Ingresso.
\beginsong{Ora è tempo di gioia}[by={Gen Rosso}]
\beginverse
L'eco torna d'antiche valli, la sua voce non porta più
ricordo di sommesse lacrime di esili in terre lontane.
\endverse
\beginchorus
Ora è tempo di gioia, non ve ne accorgete?
Ecco faccio una cosa nuova,
nel deserto una strada aprirò.
\endchorus
\beginverse
Come l'onda che sulla sabbia copre le orme e poi passa e va,
così nel tempo si cancellano le ombre scure del lungo inverno.
\endverse
\beginverse
Fra i sentieri dei boschi il vento con i rami ricomporrà
nuove armonie che trasformano i lamenti in canti di festa.
\endverse
\endsong

% Fonte: LibrettoAccordi.pdf (Canti Coro), J40 (li' tra i canti di comunione).
% Nell'elenco: Ingresso, Ordinario. CD: Venti dello Spirito (1988-2007).
\beginsong{Popoli tutti, acclamate}[by={RnS (CD Venti dello Spirito)}]
\beginverse
Mio Dio, Signore, nulla è pari a Te.
Ora e per sempre, voglio lodare il tuo grande amor per noi.
Mia roccia Tu sei, pace e conforto mi dai,
con tutto il cuore e le mie forze, sempre io ti adorerò.
\endverse
\beginchorus
Popoli tutti, acclamate al Signore,
gloria e potenza cantiamo al Re,
mari e monti si prostrino a Te, al tuo nome, o Signore.
Canto di gioia per quello che fai,
per sempre, Signore, con Te resterò,
non c'è promessa, non c'è fedeltà che in Te.
\endchorus
\endsong

% Fonte: LibrettoAccordi.pdf (Canti Coro), A21.
% Nell'elenco: Ingresso, Ordinario.
\beginsong{Quale gioia è star con te}[by={RnS}]
\beginverse
Ogni volta che ti cerco, ogni volta che t'invoco,
sempre mi accogli, Signor.
Grandi sono i tuoi prodigi, tu sei buono verso tutti,
santo tu regni tra noi.
\endverse
\beginchorus
Quale gioia è star con te, Gesù, vivo e vicino,
bello è dar lode a Te, Tu sei il Signor.
Quale dono è aver creduto in Te, che non mi abbandoni,
io per sempre abiterò la tua casa, mio Re.
\endchorus
\beginverse
Hai guarito il mio dolore, hai cambiato questo cuore,
oggi rinasco, Signor.
Grandi sono i tuoi prodigi, tu sei buono verso tutti,
santo tu regni tra noi.
\endverse
\beginverse
Hai salvato la mia vita, hai aperto la mia bocca,
canto per te, mio Signor.
Grandi sono i tuoi prodigi, tu sei buono verso tutti,
santo tu regni tra noi.
\endverse
\beginverse*
Ti loderò, ti adorerò, ti canterò,
noi ti canteremo.
\endverse
\endsong

% Fonte: LibrettoAccordi.pdf (Canti Coro), J19; c'e' anche
% "1 Canti di Ingresso/Una goccia nel mare (M. Balduzzi).pdf".
% Nell'elenco: Ingresso, Ordinario. CD: Xaipe.
\beginsong{Una goccia del mare}[by={M. Balduzzi (CD Xaipe)}]
\beginverse
Tu che per nome chiami tutte le stelle,
tu che raccogli ogni nostro segreto sospiro,
non mi abbandonare,
non mi lasciare nell'ombra del cuore,
e da te io verrò.
\endverse
\beginchorus
Una goccia del mare ti chiedo,
che disseti come rugiada,
dammi solo una goccia del mare
dell'infinita sapienza tua.
\endchorus
\beginverse
Tu che hai creato il calice dell'universo,
tu che per gli uomini hai preparato una casa,
non mi abbandonare,
non mi lasciare nell'ombra del cuore
e da te io verrò.
\endverse
\endsong

% Fonte: testo fornito dall'utente. Non presente nell'elenco. Autore non indicato.
\beginsong{Vida en abundancia}
\beginverse
Los lirios del campo
y las aves del cielo
no se preocupan
porque están en mis manos.
Tené confianza en mí,
acá estoy junto a vos.
Amá lo que sos y tus circunstancias,
estoy con vos, con tu cruz en mi espalda.
Todo terminará bien,
yo hago nuevas todas las cosas.
\endverse
\beginchorus
Yo vengo a traerte vida,
vida en abundancia, en abundancia.
Yo soy el camino, la verdad y la vida,
vida en abundancia, en abundancia.
\endchorus
\beginverse
No hice al hombre para que esté solo:
caminen juntos, como hermanos,
sopórtense mutuamente,
ámense unos a otros.
La felicidad de la vida eterna
empieza conmigo en la tierra.
Sentite vivo:
la fiesta del reino comienza acá.
\endverse
\endsong

% Fonte: LibrettoAccordi.pdf (Canti Coro), D8.
% Nell'elenco: Ingresso, Matrimonio.
\beginsong{Vieni dal Libano}[by={K. Arguello}]
\beginchorus
Vieni dal Libano, mia sposa, vieni dal Libano, vieni!
Avrai per corona le vette dei monti, le alte cime dell'Hermon.
Tu m'hai ferito, ferito il cuore, o sorella mia sposa.
Vieni dal Libano, mia sposa, vieni dal Libano, vieni.
\endchorus
\beginverse*
Cercai l'amore dell'anima mia, lo cercai senza trovarlo.
Trovai l'amore dell'anima mia, l'ho abbracciato e non lo lascerò mai!
\endverse
\beginverse
Alzati in fretta, mia diletta, vieni, colomba, vieni.
L'inverno ormai è già passato, il canto della tortora si ode.
I fiori son tornati sulla terra, il grande sole è venuto.
Alzati in fretta, mia diletta, vieni, colomba, vieni.
\endverse
\beginverse
Come sigillo sul tuo cuore, come sigillo sul braccio;
che l'amore è forte come la morte e le acque non lo spegneranno.
Dare per esso tutti i beni della casa sarebbe disprezzarlo.
Come sigillo sul tuo cuore, come sigillo sul braccio.
\endverse
\endsong

\end{songs}

% !TEX root = ../libretto.tex
% Sezione: Signore, pietà
\momento{Signore, pietà}
\begin{songs}{titoli}

% Fonte: LibrettoAccordi.pdf (Canti Coro), B5.
% Nell'elenco: Kyrie, Ordinario.
\beginsong{Kyrie}[by={M. Balduzzi}]
\beginverse
Kyrie, Kyrie eleison! Christe, Christe eleison!
Kyrie, Kyrie, Kyrie, Kyrie eleison!
\endverse
\beginverse
Kyrie, Kyrie eleison! Christe, Christe eleison!
Kyrie, Kyrie, Kyrie, Kyrie eleison!
\endverse
\endsong

% Fonte: LibrettoAccordi.pdf (Canti Coro), B6.
% Nell'elenco: Kyrie, Ordinario.
\beginsong{Kyrie eleison}[by={C. Bevilacqua}]
\beginverse
Kyrie eleison. Kyrie eleison.
Kyrie eleison, eleison.
\endverse
\beginverse
Christe eleison. Christe eleison.
Christe eleison, eleison.
\endverse
\beginverse
Kyrie eleison. Kyrie eleison.
Kyrie eleison, eleison.
\endverse
\endsong

% Fonte: LibrettoAccordi.pdf (Canti Coro), B2.
% Nell'elenco: Kyrie, Ordinario.
\beginsong{Signore, pietà\soloindice{ (D. Ricci)}}[by={D. Ricci}]
\beginverse
Signore, che a Pietro pentito hai offerto il tuo perdono,
abbi pietà di noi.
Signore, pietà! Signore, pietà! Signore, pietà! \rep{2}
\endverse
\beginverse
Cristo, che al buon ladrone hai promesso il paradiso,
abbi pietà di noi.
Cristo, pietà! Cristo, pietà! Cristo, pietà! \rep{2}
\endverse
\beginverse
Signore, che accogli ogni uomo che si affida alla tua misericordia,
abbi pietà di noi.
Signore, pietà! Signore, pietà! Signore, pietà! \rep{2}
\endverse
\endsong

% Fonte: LibrettoAccordi.pdf (Canti Coro), B1.
% Nell'elenco: Kyrie, Ordinario. CD: Vita nuova con te (1995).
\beginsong{Signore, pietà\soloindice{ (F. Buttazzo)}}[by={F. Buttazzo (CD Vita nuova con te)}]
\beginverse
Signore, che sei venuto a perdonare,
abbi pietà di noi,
abbi pietà di noi.
Signore, pietà, Signore, pietà.
\endverse
\beginverse
Cristo, che fai festa per chi ritorna a te,
abbi pietà di noi,
abbi pietà di noi.
Cristo, pietà, Cristo, pietà.
\endverse
\beginverse
Signore, che perdoni molto a chi molto ama,
abbi pietà di noi,
abbi pietà di noi.
Signore, pietà, Signore, pietà.
\endverse
\endsong

% Fonte: LibrettoAccordi.pdf (Canti Coro), B4.
% Nell'elenco: Kyrie, Ordinario.
\beginsong{Signore, pietà\soloindice{ (G. Cento)}}[by={G. Cento}]
\beginverse
Signore, Signore, pietà di noi!
Signore, pietà di noi! Signore, pietà, pietà di noi!
\endverse
\beginverse
Cristo, Cristo, pietà di noi!
Cristo, pietà di noi! Cristo, pietà, pietà di noi!
\endverse
\beginverse
Signore, Signore, pietà di noi!
Signore, pietà di noi! Signore, pietà, pietà di noi!
Pietà di noi!
\endverse
\endsong

\end{songs}

% !TEX root = ../libretto.tex
% Sezione: Gloria
\momento{Gloria}
\begin{songs}{titoli}

% Fonte: LibrettoAccordi.pdf (Canti Coro), C3.
% Nell'elenco: Gloria (senza autore).
\beginsong{Gloria a Dio}
\beginverse
Gloria a Dio nell'alto dei cieli,
e sia gloria, sia gloria,
pace in terra agli uomini di buona volontà.
\endverse
\beginverse
Noi ti lodiamo, ti benediciamo,
ti adoriamo, ti glorifichiamo,
ti rendiamo grazie per la tua gloria immensa,
Signore Dio, Re del cielo, Dio Padre Onnipotente.
\endverse
\beginchorus
E sia gloria, gloria, gloria.
\endchorus
\beginverse
Signore, Figlio Unigenito,
Gesù Cristo, Signore Dio,
Agnello di Dio, Figlio del Padre:
tu, che togli i peccati del mondo, abbi pietà di noi;
tu, che togli i peccati del mondo, accogli la nostra supplica;
tu, che siedi alla destra del Padre, abbi pietà di noi.
\endverse
\beginverse
Perché tu solo il Santo,
tu solo il Signore,
tu solo l'Altissimo, Gesù Cristo,
con lo Spirito Santo, nella gloria di Dio Padre. Amen!
\endverse
\beginverse*
Amen! Amen! Amen!
\endverse
\endsong

% Fonte: LibrettoAccordi.pdf (Canti Coro), C5.
% Nell'elenco: Gloria.
\beginsong{Gloria a Dio\soloindice{ (D. Ricci)}}[by={D. Ricci}]
\beginverse*
Gloria a Dio nell'alto dei cieli,
pace in terra agli uomini di buona volontà.
Noi ti lodiamo, ti benediciamo,
noi ti adoriamo, ti glorifichiamo.
Ti rendiamo grazie per la tua gloria immensa.
Signore Dio, Re del cielo,
Dio Padre Onnipotente, Signore.
Figlio Unigenito, Gesù Cristo, Signore Dio,
Agnello di Dio, Figlio del Padre.
Tu, che togli i peccati del mondo,
abbi pietà di noi.
Tu, che togli i peccati del mondo,
accogli la nostra supplica.
Tu, che siedi alla destra del Padre,
abbi pietà di noi.
Perché tu solo il Santo, tu solo il Signore,
tu solo l'Altissimo Gesù Cristo,
con lo Spirito Santo, nella gloria di Dio Padre.
Amen.
\endverse
\endsong

% Fonte: LibrettoAccordi.pdf (Canti Coro), C6.
% Nell'elenco: Gloria.
\beginsong{Gloria a Dio\soloindice{ (D. Ricci)}}[by={D. Ricci}]
\beginchorus
Gloria a Dio nell'alto dei cieli,
e pace in terra agli uomini di buona volontà.
Gloria a Dio nell'alto dei cieli!
Gloria. \rep{2}
\endchorus
\beginverse
Noi ti lodiamo e ti benediciamo,
noi ti adoriamo e ti glorifichiamo
e ti rendiamo noi grazie per la tua gloria immensa.
Signore Dio, Re del cielo, Dio Padre onnipotente.
\endverse
\beginverse
Signore, Figlio unigenito, Gesù Cristo,
Signore Dio, Agnello di Dio, Figlio del Padre.
Tu che togli i peccati del mondo, abbi pietà di noi.
Tu che togli i peccati del mondo, accogli la nostra supplica.
Tu che siedi alla destra di Dio Padre, abbi pietà di noi.
\endverse
\beginverse
Perché tu solo il Santo, tu solo il Signore,
tu solo l'Altissimo, Gesù Cristo,
con lo Spirito Santo, nella gloria di Dio Padre,
con lo Spirito Santo, nella gloria.
\endverse
\beginverse*
Gloria a Dio nell'alto dei cieli.
Amen.
\endverse
\endsong

% Fonte: LibrettoAccordi.pdf (Canti Coro), C7.
% Nell'elenco: Gloria.
\beginsong{Gloria a Dio\soloindice{ (E. Selleri)}}[by={E. Selleri}]
\beginchorus
Gloria, gloria, gloria a Dio nell'alto dei cieli
e pace in terra agli uomini di buona volontà.
\endchorus
\beginverse*
Noi ti lodiamo, ti benediciamo, ti adoriamo, ti glorifichiamo,
ti rendiamo grazie per la tua immensa gloria,
Signore Dio, Re del cielo, Dio Padre Onnipotente.
Signore, Figlio Unigenito, Gesù Cristo, Signore Dio,
Agnello di Dio, Figlio del Padre:
tu, che togli i peccati del mondo, abbi pietà di noi;
tu, che togli i peccati del mondo, accogli la nostra supplica;
tu, che siedi alla destra del Padre, abbi pietà di noi.
Perché tu solo il Santo, tu solo il Signore,
tu solo l'Altissimo, Gesù Cristo,
con lo Spirito Santo, nella gloria di Dio Padre.
Amen! Amen!
\endverse
\endsong

% Fonte: LibrettoAccordi.pdf (Canti Coro), C11.
% Nell'elenco: Gloria.
\beginsong{Gloria a Dio\soloindice{ (K. Arguello)}}[by={K. Arguello, G. Amadei}]
\beginchorus
Gloria a Dio nell'alto dei cieli
e pace in terra agli uomini amati dal Signore!
Ti lodiamo, ti benediciamo,
ti adoriamo, ti rendiamo grazie per la tua gloria immensa!
\endchorus
\beginverse
Signore Dio, Re del cielo, Dio Padre onnipotente,
Signore, Figlio unigenito, Gesù Cristo.
\endverse
\beginverse
Signore Dio, Agnello di Dio, Figlio del Padre:
tu che togli i peccati del mondo abbi pietà di noi.
\endverse
\beginverse
Tu che togli i peccati del mondo accogli la nostra supplica.
Tu che siedi alla destra del Padre abbi pietà di noi.
\endverse
\beginverse
Perché tu solo il Santo, tu solo il Signore,
tu solo l'Altissimo, Gesù Cristo.
Con lo Spirito Santo nella gloria di Dio Padre.
\endverse
\beginverse*
Amen! Amen!
\endverse
\endsong

% Fonte: LibrettoAccordi.pdf (Canti Coro), C8.
% Nell'elenco: Gloria.
\beginsong{Gloria a Dio nell'alto dei cieli}[by={F. Buttazzo}]
\beginchorus
Gloria a Dio nell'alto dei cieli e pace in terra
agli uomini di buona volontà.
\endchorus
\beginverse
Noi ti lodiamo, ti benediciamo,
ti adoriamo, ti glorifichiamo,
ti rendiamo noi grazie per la tua gloria immensa,
Signore Dio, Re del cielo, Dio Padre onnipotente.
\endverse
\beginverse
Signore, Figlio unigenito,
Gesù Cristo, Signore Dio,
Agnello di Dio, Figlio del Padre,
tu che togli i peccati del mondo,
abbi pietà di noi;
tu che togli i peccati del mondo,
accogli la nostra supplica;
tu che siedi alla destra del Padre,
abbi pietà di noi.
\endverse
\beginverse
Perché tu solo il Santo, tu solo il Signore,
tu solo l'Altissimo, Gesù Cristo,
con lo Spirito Santo, nella gloria di Dio Padre,
Amen. Amen.
\endverse
\endsong

% Fonte: LibrettoAccordi.pdf (Canti Coro), C10.
% Nell'elenco: Gloria. Nel finale una voce canta il ritornello e l'altra la 3a strofa.
\beginsong{Gloria\soloindice{ (Balduzzi, Giorno senza tramonto)}}[by={M. Balduzzi (CD Giorno senza tramonto)}]
\beginchorus
Gloria, gloria, gloria a Dio nei cieli e pace in terra agli uomini.
Gloria, gloria, gloria a Dio nei cieli e pace a tutti gli uomini.
\endchorus
\beginverse
Noi ti lodiamo, benediciamo, noi ti adoriamo, Signore Dio,
noi ti lodiamo, benediciamo, glorifichiamo, Signore Dio.
Rendiamo grazie per la tua gloria, immensa gloria, Signore Dio,
o re del cielo, Signore padre onnipotente, Signore Dio.
\endverse
\beginverse
Signore, figlio, unico figlio, Signore Dio, Gesù Cristo
Signore Dio, agnello di Dio, abbi pietà, abbi pietà di noi.
Tu, che dal mondo togli i peccati, accogli la nostra supplica,
Tu, che alla destra del padre siedi, abbi pietà, abbi pietà di noi.
\endverse
\beginverse
Tu solo il Santo, solo il Signore,
perché tu solo l'Altissimo,
tu nella gloria di Dio Padre,
tu con lo Spirito Santo.
Tu solo il Santo, solo il Signore,
perché tu solo l'Altissimo,
tu con lo Spirito nella gloria
di Dio Padre. Amen.
\endverse
\endsong

% Fonte: LibrettoAccordi.pdf (Canti Coro), C4. Botta e risposta Uomini/Donne/Tutti unito.
% Nell'elenco: Gloria.
\beginsong{Gloria\soloindice{ (Balduzzi, Verbum Panis)}}[by={M. Balduzzi (CD Verbum Panis)}]
\beginchorus
Gloria in excelsis Deo! Gloria in excelsis Deo! \rep{2}
\endchorus
\beginverse
Et in terra pax hominibus bonae voluntatis,
laudamus te, benedicimus te,
adoramus te, glorificamus te,
gratias agimus tibi
propter magnam gloriam tuam.
\endverse
\beginverse
Domine Deus, Domine Deus,
Rex caelestis, Rex caelestis,
Deus Pater omnipotens.
Domine Fili, Domine Fili,
Unigenite, Unigenite,
Iesu Christe.
Domine Deus, Agnus Dei, Filius Patris.
\endverse
\beginverse
Qui tollis peccata mundi, miserere nobis,
qui tollis peccata mundi,
suscipe deprecationem nostram,
qui sedes ad dexteram Patris,
miserere nobis.
\endverse
\beginverse
Quoniam tu solus, quoniam tu solus,
solus Sanctus, solus Sanctus,
tu solus Dominus.
Quoniam tu solus, quoniam tu solus,
Altissimus, Altissimus,
Iesu Christe.
Cum Sancto Spiritu, in gloria Dei Patris.
\endverse
\beginverse*
Amen!
\endverse
\endsong

% Fonte: LibrettoAccordi.pdf (Canti Coro), C2.
% Nell'elenco: Gloria ("Gloria a Dio", Buttazzo).
\beginsong{Gloria\soloindice{ (F. Buttazzo)}}[by={F. Buttazzo}]
\beginchorus
Gloria a Dio nell'alto dei cieli,
pace in terra agli uomini. \rep{2}
\endchorus
\beginverse
Ti lodiamo, ti benediciamo,
ti adoriamo, ti glorifichiamo.
Ti rendiamo grazie
per la tua gloria immensa.
Signore Dio, Re del cielo,
Dio Padre onnipotente.
Gesù Cristo, Agnello di Dio,
tu Figlio del Padre.
\endverse
\beginverse
Tu, che togli i peccati del mondo,
la nostra supplica ascolta, Signore.
Tu che siedi alla destra del Padre, abbi pietà di noi.
\endverse
\beginverse
Tu solo il Santo, tu solo il Signore, tu l'Altissimo Gesù Cristo,
con lo Spirito Santo,
nella gloria del Padre.
\endverse
\endsong

% Fonte: LibrettoAccordi.pdf (Canti Coro), C1. Le parti a eco sono unite.
% Nell'elenco: Gloria, Ordinario.
\beginsong{Gloria nell'alto dei cieli\soloindice{ (D. Ricci)}}[by={D. Ricci}]
\beginchorus
Gloria nell'alto dei cieli, gloria nell'alto dei cieli,
gloria! Gloria a Dio!
Pace agli uomini in terra, pace agli uomini in terra,
pace agli uomini di buona volontà.
\endchorus
\beginverse
Noi ti lodiamo, ti benediciamo, ti adoriamo, ti glorifichiamo.
Ti rendiamo grazie per la tua gloria immensa, gloria immensa.
Signore Dio, Re del Cielo, Dio Padre Onnipotente. \rep{2}
\endverse
\beginverse
Signore Figlio Unigenito Gesù Cristo,
Signore Dio, Agnello di Dio, Figlio del Padre.
Tu che togli i peccati del mondo, abbi pietà di noi,
tu che togli i peccati del mondo, accogli la nostra supplica,
tu che siedi alla destra di Dio Padre, abbi pietà di noi.
\endverse
\beginverse
Perché tu solo il Santo, tu solo il Signore,
tu solo l'Altissimo Gesù Cristo,
con lo Spirito Santo nella Gloria di Dio Padre,
Amen.
\endverse
\endsong

% Fonte: LibrettoAccordi.pdf (Canti Coro), C9.
% Nell'elenco: Gloria. CD: Cerco il tuo volto (2000).
\beginsong{Gloria nell'alto dei cieli\soloindice{ (Gen Verde)}}[by={Gen Verde (CD Cerco il tuo volto)}]
\beginchorus
Gloria, gloria a Dio,
gloria, gloria nell'alto dei cieli.
Pace in terra agli uomini
di buona volontà.
\endchorus
\beginverse
Noi ti lodiamo, ti benediciamo, ti adoriamo, glorifichiamo te,
ti rendiamo grazie per la tua immensa gloria.
Signore Dio, gloria!
Re del cielo, gloria!
Dio Padre, Dio onnipotente, gloria!
\endverse
\beginverse
Signore, Figlio unigenito, Gesù Cristo,
Signore, Agnello di Dio, Figlio del Padre.
Tu che togli i peccati del mondo, abbi pietà di noi;
tu che togli i peccati del mondo, accogli la nostra supplica;
tu che siedi alla destra del Padre, abbi pietà di noi!
\endverse
\beginverse
Perché tu solo il Santo, il Signore, tu solo l'Altissimo, Cristo Gesù,
con lo Spirito Santo nella gloria del Padre.
\endverse
\endsong

\end{songs}

% !TEX root = ../libretto.tex
% Sezione: Salmi
\momento{Salmi}
\begin{songs}{titoli}

% Fonte: LibrettoAccordi.pdf (Canti Coro), O23.
% Nell'elenco: Salmo, Pasqua. Messo tra i Salmi (regola dell'utente: i salmi stanno nei Salmi).
\beginsong{Alleluia, rendete grazie (Epistola)}[by={Gen Verde (CD Il mistero pasquale)}]
\beginchorus
Alleluia, alleluia, alleluia,
alleluia, alleluia. \rep{2}
\endchorus
\beginverse
Rendete grazie a Dio, Egli è buono,
eterno e fedele è il suo amore.
Sì, è così: lo dica Israele,
dica che il suo amore è per sempre.
\endverse
\beginverse
La destra del Signore si è innalzata
a compiere grandiose meraviglie.
Non morirò, ma resterò in vita
e annuncerò i prodigi del Signore.
\endverse
\beginverse
La pietra che avevano scartato
è divenuta pietra angolare.
Questo prodigio ha fatto il Signore:
una meraviglia ai nostri occhi.
\endverse
\endsong

% Fonte: LibrettoAccordi.pdf (Canti Coro), D13.
% Nell'elenco: Salmo, Anno A.
\beginsong{Ascoltate oggi}[by={P. Spoladore}]
\beginchorus
Ascoltate oggi la voce del Signore:
non indurite il vostro cuore.
\endchorus
\beginverse
Venite, cantiamo al Signore,
acclamiamo alla roccia della nostra salvezza.
Accostiamoci a lui per rendergli grazie,
a lui acclamiamo con canti di gioia.
\endverse
\beginverse
Entrate: prostrati, adoriamo,
in ginocchio davanti al Signore che ci ha fatti.
È lui il nostro Dio e noi il popolo del suo pascolo,
il gregge che egli conduce.
\endverse
\beginverse
Se ascoltaste oggi la sua voce!
«Non indurite il cuore come a Meriba,
come nel giorno di Massa nel deserto,
dove mi tentarono i vostri padri:
mi misero alla prova pur avendo visto le mie opere».
\endverse
\endsong

% Fonte: LibrettoAccordi.pdf (Canti Coro), D18.
% Nell'elenco: Salmo, Anno B.
\beginsong{Camminerò}[by={P. Spoladore}]
\beginchorus
Camminerò alla presenza del Signore
nella terra dei viventi.
\endchorus
\beginverse
Ho creduto anche quando dicevo:
«Sono troppo infelice».
Agli occhi del Signore è preziosa
la morte dei suoi fedeli.
\endverse
\beginverse
Ti prego, Signore, perché sono tuo servo;
io sono tuo servo, figlio della tua schiava:
tu hai spezzato le mie catene.
A te offrirò un sacrificio di ringraziamento
e invocherò il nome del Signore.
\endverse
\beginverse
Adempirò i miei voti al Signore
davanti a tutto il suo popolo,
negli atri della casa del Signore,
in mezzo a te, Gerusalemme.
\endverse
\endsong

% Fonte: LibrettoAccordi.pdf (Canti Coro), D21.
% Nell'elenco: Salmo, Anno B.
\beginsong{Crea in me, o Dio, un cuore puro}[by={P. Spoladore}]
\beginchorus
Crea in me, o Dio, un cuore puro. \rep{2}
\endchorus
\beginverse
Pietà di me, o Dio, nel tuo amore; nella tua grande misericordia
cancella la mia iniquità.
Lavami tutto dalla mia colpa,
dal mio peccato rendimi puro.
\endverse
\beginverse
Crea in me, o Dio, un cuore puro,
rinnova in me uno spirito saldo.
Non scacciarmi dalla tua presenza
e non privarmi del tuo santo spirito.
\endverse
\beginverse
Rendimi la gioia della tua salvezza,
sostienimi con uno spirito generoso.
Insegnerò ai ribelli le tue vie
e i peccatori a te ritorneranno.
\endverse
\endsong

% Fonte: LibrettoAccordi.pdf (Canti Coro), D1.
% Nell'elenco: Salmo, Pasqua. Messo tra i Salmi (regola dell'utente: i salmi stanno nei Salmi).
\beginsong{Del tuo Spirito, Signore}[by={Gen Verde}]
\beginchorus
Del tuo Spirito, Signore,
è piena la terra, è piena la terra. \rep{2}
\endchorus
\beginverse
Benedici il Signore, anima mia, Signore, Dio, tu sei grande!
Sono immense, splendenti tutte le tue opere e tutte le creature.
\endverse
\beginverse
Se tu togli il tuo soffio muore ogni cosa, e si dissolve nella terra.
Il tuo Spirito scende: tutto si ricrea e tutto si rinnova.
\endverse
\beginverse
La tua gloria, Signore, resti per sempre. Gioisci, Dio, del creato.
Questo semplice canto salga a te Signore, sei tu la nostra gioia.
\endverse
\endsong

% Fonte: LibrettoAccordi.pdf (Canti Coro), D16.
% Nell'elenco: Salmo, Pasqua. Messo tra i Salmi (regola dell'utente: i salmi stanno nei Salmi).
\beginsong{Dio mio, Dio mio}[by={P. Spoladore}]
\beginchorus
Dio mio, Dio mio, perché mi hai abbandonato? \rep{2}
\endchorus
\beginverse
Si fanno beffe di me quelli che mi vedono,
storcono le labbra, scuotono il capo:
«Si rivolga al Signore, lui lo liberi,
lo porti in salvo se davvero lo ama!».
\endverse
\beginverse
Un branco di cani mi circonda,
mi accerchia una banda di malfattori;
hanno scavato le mie mani e i miei piedi.
Posso contare tutte le mie ossa.
\endverse
\beginverse
Si dividono le mie vesti,
sulla mia tunica gettano la sorte.
Ma tu, Signore, non stare lontano,
mia forza, vieni presto in mio aiuto.
\endverse
\beginverse
Annuncerò il tuo nome ai miei fratelli,
ti loderò in mezzo all'assemblea.
Lodate il Signore, voi suoi fedeli,
gli dia gloria tutta la discendenza di Giacobbe,
lo tema tutta la discendenza di Israele.
\endverse
\endsong

% Fonte: LibrettoAccordi.pdf (Canti Coro), D12.
% Nell'elenco: Salmo, Anno A.
\beginsong{Donaci, o Signore}[by={P. Spoladore}]
\beginchorus
Donaci, o Signore, il tuo amore: in te speriamo. \rep{2}
\endchorus
\beginverse
Retta è la parola del Signore e fedele ogni sua opera.
Egli ama la giustizia e il diritto; dell'amore del Signore è piena la terra.
\endverse
\beginverse
Ecco, l'occhio del Signore è su chi lo teme, su chi spera nel suo amore,
per liberarlo dalla morte e nutrirlo in tempo di fame.
\endverse
\beginverse
L'anima nostra attende il Signore: egli è nostro aiuto e nostro scudo.
Su di noi sia il tuo amore, Signore, come da te noi speriamo.
\endverse
\endsong

% Fonte: LibrettoAccordi.pdf (Canti Coro), D26.
% Nell'elenco: Salmo, Anno C.
\beginsong{Grandi cose}[by={P. Spoladore}]
\beginchorus
Grandi cose ha fatto il Signore per noi. \rep{2}
\endchorus
\beginverse
Quando il Signore ristabilì la sorte di Sion,
ci sembrava di sognare.
Allora la nostra bocca si riempì di sorriso,
la nostra lingua di gioia.
\endverse
\beginverse
Allora si diceva tra le genti:
«Il Signore ha fatto grandi cose per loro».
Grandi cose ha fatto il Signore per noi:
eravamo pieni di gioia.
\endverse
\beginverse
Ristabilisci, Signore, la nostra sorte,
come i torrenti del Negheb.
Chi semina nelle lacrime
mieterà nella gioia.
\endverse
\beginverse
Nell'andare, se ne va piangendo,
portando la semente da gettare,
ma nel tornare, viene con gioia,
portando i suoi covoni.
\endverse
\endsong

% Fonte: LibrettoAccordi.pdf (Canti Coro), D2. Ordine delle righe ricostruito e confermato dall'utente.
% Nell'elenco: Salmo 29, Pasqua. Messo tra i Salmi (regola dell'utente: i salmi stanno nei Salmi).
\beginsong{Grida di vittoria}[by={FdV}]
\beginchorus
Grida di vittoria, grida di esultanza! \rep{2}
La destra di Dio si esaltò! \rep{4}
\endchorus
\beginverse
Ti esalto, Signore, mi hai liberato,
non hai permesso che sul mio capo
potessero esultare i miei nemici,
ti rendo grazie, mio Dio.
Signore Dio, a te ho gridato,
Signore Dio, mi hai guarito.
Dagli inferi mi hai fatto risalire,
mi hai dato vita, Signore.
\endverse
\beginverse
Cantate a Dio, o suoi fedeli,
rendete grazie al suo santo nome.
Dura un istante la sua ira,
la sua bontà in eterno.
Alla sera sopraggiunge il pianto
ed al mattino ecco la gioia.
Signore Dio, mi hai salvato,
mi hai liberato dalla morte.
\endverse
\beginverse
Ascolta, Signore, misericordia!
Vieni, Signore, in mio aiuto!
A te io grido, Signore Dio!
Tu dalla morte mi salverai!
Il mio lamento hai mutato in danza,
mi hai messo l'abito della gioia.
Senza fine io canterò,
per sempre, mio Dio, ti loderò.
\endverse
\endsong

% Fonte: LibrettoAccordi.pdf (Canti Coro), D25.
% Nell'elenco: Salmo, Anno C.
\beginsong{Gustate e vedete}[by={P. Spoladore}]
\beginchorus
Gustate e vedete com'è buono il Signore. \rep{2}
\endchorus
\beginverse
Benedirò il Signore in ogni tempo,
sulla mia bocca sempre la sua lode.
Io mi glorio nel Signore:
i poveri ascoltino e si rallegrino.
\endverse
\beginverse
Magnificate con me il Signore,
esaltiamo insieme il suo nome.
Ho cercato il Signore: mi ha risposto
e da ogni mia paura mi ha liberato.
\endverse
\beginverse
Guardate a lui e sarete raggianti,
i vostri volti non dovranno arrossire.
Questo povero grida e il Signore lo ascolta,
lo salva da tutte le sue angosce.
\endverse
\endsong

% Fonte: "4 - Salmi/Il cantico dei redenti.pdf" (melodia con accordi). Isaia 12.
% Nell'elenco: Salmo. Strofe 3 e 4 (tagliate nel foglio) dal testo fornito dall'utente.
\beginsong{Il cantico dei redenti}[by={A. Marani}]
\beginchorus
Il Signore è la mia salvezza
e con lui non temo più,
perché ho nel cuore la certezza:
la salvezza è qui con me.
\endchorus
\beginverse
Ti lodo, Signore, perché
un giorno eri lontano da me,
ora invece sei tornato
e mi hai preso con te.
\endverse
\beginverse
Berrete con gioia alle fonti,
alle fonti della salvezza,
e quel giorno voi direte:
lodate il Signore, invocate il suo nome.
\endverse
\beginverse
Fate conoscere ai popoli
tutto quello che lui ha compiuto,
e ricordino per sempre,
ricordino sempre che il suo nome è grande.
\endverse
\beginverse
Cantate a chi ha fatto grandezze
e sia fatto sapere nel mondo;
grida forte la tua gioia, abitante di Sion,
perché grande con te è il Signore.
\endverse
\endsong

% Fonte: LibrettoAccordi.pdf (Canti Coro), D3.
% Nell'elenco: Salmo, Pasqua. Messo tra i Salmi (regola dell'utente: i salmi stanno nei Salmi).
\beginsong{Il canto del mare}[by={M. Frisina}]
\beginchorus
Cantiamo al Signore, stupenda è la sua vittoria.
Signore è il suo nome. Alleluja!
\endchorus
\beginverse
Voglio cantare in onore del Signore,
perché ha trionfato, alleluja.
Ha gettato in mare cavallo e cavaliere.
Mia forza e mio canto è il Signore.
Il mio Salvatore, il Dio di mio padre,
ed io lo voglio esaltare.
\endverse
\beginverse
Dio è prode in guerra, si chiama Signore.
Travolse nel mare gli eserciti,
i carri d'Egitto sommerse nel Mar Rosso,
abissi profondi li coprono.
La tua destra, Signore, si è innalzata,
la tua potenza è terribile.
\endverse
\beginverse
Si accumularon le acque al tuo soffio,
s'alzaron le onde come un argine.
Si raggelaron gli abissi in fondo al mare.
Chi è come te, o Signore?
Guidasti con forza il popolo redento,
e lo conducesti verso Sion.
\endverse
\endsong

% Fonte: LibrettoAccordi.pdf (Canti Coro), D28.
% Nell'elenco: Salmo.
\beginsong{Il dramma messianico}[by={M. Maneri}]
\beginverse
Perché le genti congiurano,
perché invano cospirano i popoli?
Insorgono i re della terra
congiurano insieme contro il Signore e contro il suo Messia:
«Spezziamo le loro catene,
gettiamo via i loro legami!».
Se ne ride chi abita i cieli,
li schernisce dall'alto il Signore.
Egli parla con ira
e li spaventa con il suo sdegno:
«Io l'ho costituito sovrano sul mio santo monte».
Annunzierò il decreto di Dio. Egli mi ha detto:
«Tu sei mio figlio, oggi io ti ho generato.
Io ti darò in possesso le genti di tutta la terra.
Le spezzerai con scettro di ferro,
come vasi di argilla e le frantumerai».
E ora sovrani siate saggi, istruitevi, giudici della terra;
servite Dio con timore e con tremore esultate.
che non si sdegni e voi perdiate la via.
Improvvisa divampa la sua ira.
Beato chi in lui si rifugia.
\endverse
\endsong

% Fonte: LibrettoAccordi.pdf (Canti Coro), D20.
% Nell'elenco: Salmo, Anno B.
\beginsong{Il ricordo di te, Signore}[by={P. Spoladore}]
\beginchorus
Il ricordo di te, Signore,
è la nostra gioia. \rep{2}
\endchorus
\beginverse
Lungo i fiumi di Babilonia,
là sedevamo e piangevamo
ricordandoci di Sion.
Ai salici di quella terra appendemmo le nostre cetre.
\endverse
\beginverse
Perché là ci chiedevano parole di canto coloro che ci avevano deportato,
allegre canzoni, i nostri oppressori: «Cantateci canti di Sion!».
\endverse
\beginverse
Come cantare i canti del Signore
in terra straniera?
Se mi dimentico di te, Gerusalemme, si dimentichi di me la mia destra.
\endverse
\beginverse
Mi si attacchi la lingua al palato
se lascio cadere il tuo ricordo,
se non innalzo Gerusalemme al di sopra di ogni mia gioia.
\endverse
\endsong

% Fonte: LibrettoAccordi.pdf (Canti Coro), D15.
% Nell'elenco: Salmo, Anno A.
\beginsong{Il Signore è bontà e misericordia}[by={P. Spoladore}]
\beginchorus
Il Signore è bontà e misericordia. Il Signore è bontà e misericordia. \rep{2}
\endchorus
\beginverse
Dal profondo a te grido, o Signore;
Signore, ascolta la mia voce.
Siano i tuoi orecchi attenti
alla voce della mia supplica.
\endverse
\beginverse
Se consideri le colpe, o Signore,
Signore, chi ti può resistere?
Ma con te è il perdono:
così avremo il tuo timore.
\endverse
\beginverse
Io spero, Signore.
Spera l'anima mia, attende la sua parola.
L'anima mia è rivolta al Signore
più che le sentinelle all'aurora.
\endverse
\beginverse
Più che le sentinelle l'aurora.
Israele attenda il Signore,
perché con il Signore è la misericordia e grande è con lui la redenzione.
Egli redimerà Israele da tutte le sue colpe.
\endverse
\endsong

% Fonte: LibrettoAccordi.pdf (Canti Coro), D14.
% Nell'elenco: Salmo, Anno A.
\beginsong{Il Signore è il mio pastore}[by={P. Spoladore}]
\beginchorus
Il Signore è il mio pastore, non manco di nulla. \rep{2}
\endchorus
\beginverse
Il Signore è il mio pastore, non manco di nulla.
Su pascoli erbosi mi fa riposare,
ad acque tranquille mi conduce.
Rinfranca l'anima mia.
\endverse
\beginverse
Mi guida per il giusto cammino a motivo del suo nome.
Anche se vado per una valle oscura
non temo alcun male, perché tu sei con me.
Il tuo bastone e il tuo vincastro mi danno sicurezza.
\endverse
\beginverse
Davanti a me tu prepari una mensa
sotto gli occhi dei miei nemici.
Ungi di olio il tuo capo;
il mio calice trabocca.
\endverse
\beginverse
Sì, bontà e fedeltà mi saranno compagne
tutti i giorni della mia vita,
abiterò ancora nella casa del Signore
per lunghi giorni.
\endverse
\endsong

% Fonte: LibrettoAccordi.pdf (Canti Coro), D23.
% Nell'elenco: Salmo, Anno C.
\beginsong{Il Signore è mia luce}[by={P. Spoladore}]
\beginchorus
Il Signore è mia luce e mia salvezza.
Il Signore è mia luce e mia salvezza. \rep{2}
\endchorus
\beginverse
Il Signore è mia luce e mia salvezza:
di chi avrò timore?
Il Signore è difesa della mia vita:
di chi avrò paura?
\endverse
\beginverse
Ascolta, Signore, la mia voce.
Io grido: abbi pietà di me, rispondimi!
Il mio cuore ripete il tuo invito:
«Cercate il mio volto!» Il tuo volto, Signore, io cerco.
\endverse
\beginverse
Non nascondermi il tuo volto,
non respingere con ira il tuo servo.
Sei tu il mio aiuto, non lasciarmi,
non abbandonarmi, Dio della mia salvezza.
\endverse
\beginverse
Sono certo di contemplare la bontà del Signore
nella terra dei viventi.
Spera nel Signore, sii forte,
si rinsaldi il tuo cuore e spera nel Signore.
\endverse
\endsong

% Fonte: LibrettoAccordi.pdf (Canti Coro), D24.
% Nell'elenco: Salmo, Anno C.
\beginsong{Il Signore ha pietà del suo popolo}[by={P. Spoladore}]
\beginchorus
Il Signore ha pietà del suo popolo.
Il Signore ha pietà del suo popolo. \rep{2}
\endchorus
\beginverse
Benedici il Signore, anima mia,
quanto è in me benedica il suo santo nome.
Benedici il Signore, anima mia,
non dimenticare tutti i suoi benefici.
\endverse
\beginverse
Egli perdona tutte le tue colpe,
guarisce tutte le tue infermità,
salva dalla fossa la tua vita,
ti circonda di bontà e misericordia.
\endverse
\beginverse
Il Signore compie cose giuste,
difende i diritti di tutti gli oppressi.
Ha fatto conoscere a Mosè le sue vie,
le sue opere ai figli d'Israele.
\endverse
\beginverse
Misericordioso e pietoso è il Signore,
lento all'ira e grande nell'amore.
Perché quanto il cielo è alto sulla terra,
così la sua misericordia è potente su quelli che lo temono.
\endverse
\endsong

% Fonte: "18 - Frisian/Lacqua_viva_accordi.pdf" (foglio accordi).
% Nell'elenco: Salmo.
\beginsong{L'acqua viva}[by={M. Frisina}]
\beginchorus
Chi berrà la mia acqua
non avrà più sete in eterno
e quest'acqua sarà per lui
fonte di vita per l'eternità.
\endchorus
\beginverse
Affannati e stanchi,
voi oppressi e poveri venite,
attingete con gioia a Lui
alla sorgente di felicità.
\endverse
\beginverse
Fiumi di acqua viva
sgorgheranno in colui che crederà
nel Signore che dona a noi
l'acqua di vita e di verità.
\endverse
\beginverse
Percuotendo una roccia
dissetasti il popolo in cammino,
fa' che sempre noi camminiam
nel tuo timore e nella fedeltà.
\endverse
\beginverse
Fonte inesauribile,
pace eterna, carità perfetta,
noi a mensa con te sediam,
dolce, immensa, santa Trinità.
\endverse
\beginverse*
Amen.
\endverse
\endsong

% Fonte: LibrettoAccordi.pdf (Canti Coro), D10.
% Nell'elenco: Salmo, Ordinario.
\beginsong{Loda il Signore, anima mia}[by={Gen Rosso}]
\beginchorus
Loda il Signore, anima mia,
loda il Signore, anima mia! \rep{2}
\endchorus
\beginverse
Il Signore rimane fedele per sempre,
Lui fa giustizia agli umiliati.
Agli affamati dà ogni giorno il pane,
e libertà a chi è chiuso in prigione.
\endverse
\beginverse
Il Signore dà luce a chi vive nel buio,
Lui risolleva chi è caduto.
E tutti i giusti colma col suo amore,
ai forestieri dà un tetto sicuro.
\endverse
\beginverse
Il Signore è vicino a chi non ha nessuno,
Lui non sorride a chi è violento.
Dovunque arriva il regno del Signore,
vive in principio, ora e per sempre.
\endverse
\endsong

% Fonte: "99 - Raccolte e libretti accordi/Il mistero pasquale_ gv.pdf" (pag. 61). Testo dal Salmo 30.
% Nell'elenco: Salmo, Settimana Santa (Venerdì Santo). CD: Il mistero pasquale (2011). Messo tra i Salmi (regola dell'utente: i salmi stanno nei Salmi).
\beginsong{Nelle tue mani\soloindice{ (Gen Verde)}}[by={Gen Verde (CD Il mistero pasquale)}]
\beginchorus
Padre, nelle tue mani consegno il mio spirito.
\endchorus
\beginverse
In te, Signore, mi sono rifugiato,
mai sarò deluso;
difendimi per la tua giustizia.
Affido alle tue mani il mio spirito;
tu mi hai riscattato, Signore, Dio fedele.
\endverse
\beginverse
Per i nemici e per i miei vicini
io sono un rifiuto,
l'obbrobrio ed il terrore di tutta la gente,
un corpo rinsecchito, un cadavere;
sono un vaso rotto, un coccio da gettare.
\endverse
\beginverse
Ma io confido in te, o mio Signore:
tu sei il mio Dio
e nelle mani tue sono tutti i miei giorni,
dai miei persecutori mi liberi;
dalle loro mani mi liberi, Signore.
\endverse
\beginverse
Sul tuo servo fa' splendere il tuo volto,
salvami, Signore,
per la tua fedeltà e la misericordia.
Coraggio, il vostro cuore si rianimi;
tutti voi sperate, sperate nel Signore.
\endverse
\endsong

% Fonte: LibrettoAccordi.pdf (Canti Coro), A15.
% Nell'elenco: Salmo, Natale. Messo tra i Salmi (regola dell'utente: i salmi stanno nei Salmi).
\beginsong{Oggi è nato per noi}[by={F. Baggio}]
\beginchorus
Oggi è nato per noi il Salvatore del mondo. \rep{2}
\endchorus
\beginverse
Cantate al Signore un cantico nuovo,
cantate al Signore da tutta la terra.
Cantate al Signore e benedite il suo nome.
\endverse
\beginverse
Cantate ogni giorno la sua salvezza,
in mezzo alle genti narrate la gloria,
a tutti i popoli dite le sue meraviglie.
\endverse
\beginverse
Gioiscano i cieli, esulti la terra,
risuoni il mare e quanto racchiude;
sia festa in campagna come nella foresta.
\endverse
\beginverse
Davanti al Signore che viene dal cielo
per giudicare tutta la terra;
giudicherà con giustizia e nella sua fedeltà.
\endverse
\endsong

% Fonte: LibrettoAccordi.pdf (Canti Coro), D11. Salmo 50 (51).
% Nell'elenco: Salmo, Anno A (I domenica di Quaresima).
\beginsong{Perdonaci, Signore}[by={P. Spoladore}]
\beginchorus
Perdonaci, Signore: abbiamo peccato. \rep{2}
\endchorus
\beginverse
Pietà di me, o Dio, nel tuo amore; nella tua grande misericordia
cancella la mia iniquità.
Lavami tutto dalla mia colpa,
dal mio peccato rendimi puro.
\endverse
\beginverse
Sì, le mie iniquità io le riconosco,
il mio peccato mi sta sempre dinanzi.
Contro di te, contro te solo ho peccato,
quello che è male ai tuoi occhi io l'ho fatto.
\endverse
\beginverse
Crea in me, o Dio, un cuore puro,
rinnova in me uno spirito saldo.
Non scacciarmi dalla tua presenza
e non privarmi del tuo santo spirito.
\endverse
\beginverse
Rendimi la gioia della tua salvezza,
sostienimi con uno spirito generoso.
Signore, apri le mie labbra
e la mia bocca proclami la tua lode.
\endverse
\endsong

% Fonte: testo con accordi fornito dall'utente (accordi tolti). Salmo 15.
% Nell'elenco: "Proteggimi o Dio", Salmo (autore in elenco: Gallotta; l'utente indica Ritinha, Segatori).
\beginsong{Proteggimi, o Dio}[by={Ritinha, Segatori}]
\beginverse
Proteggimi, o Dio: in te io mi rifugio.
Ho detto a lui: «Sei tu il mio Signore,
senza di te non ho bene alcuno».
Nelle tue mani, Signore, è la mia vita!
\endverse
\beginchorus
Tu mi indicherai il sentiero della vita,
gioia piena nella tua presenza,
dolcezza senza fine. \rep{2}
\endchorus
\beginverse
Benedico Dio che m'ha dato consiglio;
anche di notte il cuore m'istruisce.
Innanzi a me sempre il Signore,
sta alla mia destra, non posso vacillare.
\endverse
\beginchorus
Tu mi indicherai il sentiero della vita,
gioia piena nella tua presenza,
dolcezza senza fine. \rep{2}
\endchorus
\beginverse
Mia eredità, mio calice è il Signore,
per me la sorte è su luoghi deliziosi.
Lieto e sereno è il cuore mio,
luce e speranza ai miei passi tu darai.
\endverse
\beginchorus
Tu mi indicherai il sentiero della vita,
gioia piena nella tua presenza,
dolcezza senza fine. \rep{2}
\endchorus
\endsong

% Fonte: LibrettoAccordi.pdf (Canti Coro), D22. Salmo 90 (91).
% Nell'elenco: Salmo, Anno C (I domenica di Quaresima).
\beginsong{Resta con noi}[by={P. Spoladore}]
\beginchorus
Resta con noi, Signore, nell'ora della prova, Signore.
\endchorus
\beginverse
Chi abita al riparo dell'Altissimo
passerà la notte all'ombra dell'Onnipotente.
Io dico al Signore: «Mio rifugio e mia fortezza,
mio Dio in cui confido».
\endverse
\beginverse
Non ti potrà colpire la sventura,
nessun colpo cadrà sulla tua tenda.
Egli per te darà ordine ai suoi angeli
di custodirti in tutte le tue vie.
\endverse
\beginverse
Sulle mani essi ti porteranno,
perché il tuo piede non inciampi nella pietra.
Calpesterai leoni e vipere,
schiaccerai leoncelli e draghi.
\endverse
\beginverse
«Lo libererò, perché a me si è legato,
lo porrò al sicuro, perché ha conosciuto il mio nome.
Mi invocherà e io gli darò risposta; nell'angoscia io sarò con lui,
lo libererò e lo renderò glorioso».
\endverse
\endsong

% Fonte: LibrettoAccordi.pdf (Canti Coro), D9.
% Nell'elenco: Salmo.
\beginsong{Salmo 50}[by={RnS}]
\beginverse
Pietà di me, o Dio, secondo la tua misericordia;
nel tuo grande amore cancella il mio peccato.
Lavami dalle mie colpe, mondami dal mio peccato.
Riconosco la mia colpa, il mio peccato mi sta davanti.
\endverse
\beginchorus
Nel tuo amore fa' grazia a Sion, rialza le mura di Gerusalemme. \rep{2}
\endchorus
\beginverse
Crea in me, o Dio, un cuore puro;
rinnova in me uno spirito saldo.
Non allontanarmi dalla tua presenza,
e non privarmi del tuo Santo Spirito.
\endverse
\beginverse
Purificami con issopo ed io sarò mondato;
lavami e sarò più bianco della neve.
Fammi sentire gioia e letizia,
esulteranno le ossa che hai spezzato.
\endverse
\endsong

% Fonte: LibrettoAccordi.pdf (Canti Coro), D19. Salmo 18 (19).
% Nell'elenco: Salmo, Anno B (III domenica di Quaresima).
\beginsong{Signore, tu hai parole di vita eterna}[by={P. Spoladore}]
\beginchorus
Signore, tu hai parole di vita eterna.
Signore, tu hai parole di vita eterna.
\endchorus
\beginverse
La legge del Signore è perfetta,
rinfranca l'anima;
la testimonianza del Signore è stabile,
rende saggio il semplice.
\endverse
\beginverse
I precetti del Signore sono retti,
fanno gioire il cuore;
il comando del Signore è limpido,
illumina gli occhi.
\endverse
\beginverse
Il timore del Signore è puro,
rimane per sempre;
i giudizi del Signore sono fedeli,
sono tutti giusti.
\endverse
\beginverse
Più preziosi dell'oro,
di molto oro fino,
più dolci del miele
e di un favo stillante.
\endverse
\endsong

% Fonte: LibrettoAccordi.pdf (Canti Coro), D27. Salmo 29 (30).
% Nell'elenco: Salmo.
\beginsong{Ti esalterò}[by={P. Spoladore}]
\beginchorus
Ti esalterò, Signore! Ti esalterò, Signore! \rep{2}
Perché mi hai risollevato. \rep{4}
\endchorus
\beginverse
Ti esalterò, Signore, perché mi hai risollevato,
non hai permesso ai miei nemici di gioire su di me.
Signore mio Dio, a te ho gridato e mi hai guarito.
Signore, hai fatto risalire la mia vita dagli inferi.
Mi hai fatto rivivere perché non scendessi nella fossa.
\endverse
\beginverse
Cantate inni al Signore, o suoi fedeli,
della sua santità celebrate il ricordo.
Perché la sua collera dura un istante,
la sua bontà per tutta la vita.
Alla sera è ospite il pianto, al mattino la gioia, la gioia.
\endverse
\beginverse
Ascolta, Signore, abbi pietà di me,
Signore, vieni in mio aiuto.
Hai mutato il mio lamento in danza,
Signore mio Dio, ti renderò grazie.
Signore mio Dio, ti renderò grazie per sempre.
\endverse
\endsong

% Fonte: LibrettoAccordi.pdf (Canti Coro), D17. Salmo 24 (25).
% Nell'elenco: Salmo, Anno B (I domenica di Quaresima).
\beginsong{Tutti i sentieri}[by={P. Spoladore}]
\beginchorus
Tutti i sentieri del Signore
sono amore e fedeltà, sono amore e fedeltà.
\endchorus
\beginverse
Fammi conoscere, Signore, le tue vie,
insegnami i tuoi sentieri.
Guidami nella tua fedeltà e istruiscimi,
perché sei tu il Dio della mia salvezza.
\endverse
\beginverse
Ricordati, Signore, della tua misericordia
e del tuo amore, che è da sempre.
Ricordati di me nella tua misericordia,
per la tua bontà, Signore.
\endverse
\beginverse
Buono e retto è il Signore,
indica ai peccatori la via giusta;
guida i poveri secondo giustizia,
insegna ai poveri la sua via.
\endverse
\endsong

% Fonte: LibrettoAccordi.pdf (Canti Coro), A20 (li' tra i canti d'ingresso).
% Nell'elenco: Salmo, Ordinario. CD: Vita nuova con te (1995).
\beginsong{Un cuore nuovo}[by={F. Buttazzo (CD Vita nuova con te)}]
\beginchorus
Ti darò un cuore nuovo, popolo mio,
il mio Spirito effonderò in te.
Toglierò da te il cuore di pietra,
un cuore di carne ti darò, popolo mio.
\endchorus
\beginverse
Da tutte le nazioni vi radunerò,
vi mostrerò la strada della vita.
E vivrà chi la seguirà.
\endverse
\beginverse
Vi aspergerò con acqua, e puri vi farò,
dagli idoli sarete liberati.
Questa è la mia libertà.
\endverse
\beginverse
Mio popolo sarete, le genti lo vedranno,
abiterete dentro la mia casa,
e vedrete il mio volto.
\endverse
\endsong

\end{songs}

% !TEX root = ../libretto.tex
% Sezione: Offertorio
\momento{Offertorio}
\begin{songs}{titoli}

% Fonte: LibrettoAccordi.pdf (Canti Coro), G6.
% Nell'elenco: Offertorio, Ordinario.
\beginsong{Accogli i nostri doni}
\beginverse
Accogli Signore i nostri doni
in questo misterioso incontro,
tra la nostra povertà e la tua grandezza.
\endverse
\beginverse
Noi ti offriamo le cose che tu stesso ci hai dato
e tu in cambio donaci, donaci te stesso.
\endverse
\endsong

% Fonte: LibrettoAccordi.pdf (Canti Coro), G11.
% Nell'elenco: Offertorio, Ordinario.
\beginsong{Antica eterna danza}[by={Gen Verde}]
\beginverse
Spighe d'oro al vento, antica eterna danza
per fare un solo pane spezzato sulla mensa.
Grappoli dei colli, profumo di letizia
per fare un solo vino, bevanda della grazia.
\endverse
\beginverse
Con il pane e il vino, Signore, ti doniamo
le nostre gioie pure, le attese e le paure,
frutti del lavoro e fede nel futuro,
la voglia di cambiare e di ricominciare.
\endverse
\beginverse
Dio della speranza, sorgente d'ogni dono,
accogli questa offerta che insieme ti portiamo.
Dio dell'universo, raccogli chi è disperso
e facci tutti Chiesa, una cosa in te.
\endverse
\endsong

% Fonte: LibrettoAccordi.pdf (Canti Coro), G15.
% Nell'elenco: Offertorio, Ordinario.
\beginsong{Benedici, o Signore}[by={Gen Rosso}]
\beginverse
Nebbia e freddo, giorni lunghi e amari mentre il seme muore.
Poi il prodigio antico e sempre nuovo del primo filo d'erba.
E nel vento dell'estate ondeggiano le spighe:
avremo ancora pane.
\endverse
\beginchorus
Benedici, o Signore,
questa offerta che portiamo a Te,
facci «uno» come il pane
che anche oggi hai dato a noi.
\endchorus
\beginverse
Nei filari dopo il lungo inverno fremono le viti.
La rugiada avvolge nel silenzio i primi tralci verdi.
Poi i colori dell'autunno coi grappoli maturi
avremo ancora vino.
\endverse
\endsong

% Fonte: LibrettoAccordi.pdf (Canti Coro), G10.
% Nell'elenco: Offertorio, Ordinario.
\beginsong{Come tanti fiumi}[by={Grasso, Selleri}]
\beginverse
Come tanti fiumi che nel mare si rispecchiano,
come tanti raggi che dal sole si diramano,
così noi, Signore Dio, dalle nostre mani
un solo pane ti offriamo, perché in te, in te
segno d'unità diventerà.
\endverse
\beginchorus
Fa' di noi, Signore, fiumi di speranza,
fa' che i nostri occhi vedano l'invisibile;
e nel volto dei fratelli e nelle nostre città
siamo semi di pace e d'unità.
\endchorus
\beginverse
Come dai sentieri una grande strada partirà,
come tante voci una parola nascerà,
così noi, Signore Dio, dalle nostre mani
un po' di vino ti offriamo, perché in te, in te
segno d'unità diventerà.
\endverse
\beginverse
Fa' di noi, Signore, offerta gradita a te,
facci vivere e sognare un mondo che
ama ogni uomo, e nella diversità
scopre te, che sei amore ed unità.
Scopre te, che sei amore ed unità.
\endverse
\endsong

% Fonte: LibrettoAccordi.pdf (Canti Coro), G9.
% Nell'elenco: Offertorio, Ordinario.
\beginsong{Cosa offrirti}[by={D. Branca, L. Ciancio}]
\beginverse
Cosa offrirti, o Dio, cosa posso darti,
eccomi son qui davanti a te.
Le gioie ed i dolori, gli affanni di ogni giorno,
tutto voglio vivere in te.
\endverse
\beginchorus
Accetta, mio re, questo poco che ho,
offro a Te la mia vita.
Gioia è per me far la tua volontà,
il mio unico bene sei solo Tu, solo Tu.
\endchorus
\beginverse
Vengo a te mio Dio, apro le mie braccia,
che la tua letizia riempirà.
Rinnova questo cuore perché ti sappia amare,
e nella tua pace io vivrò.
\endverse
\endsong

% Fonte: LibrettoAccordi.pdf (Canti Coro), G7.
% Nell'elenco: Offertorio, Ordinario.
\beginsong{Ecco quel che abbiamo}[by={Gen Verde}]
\beginverse
Ecco quel che abbiamo, nulla ci appartiene ormai.
Ecco i frutti della terra, che Tu moltiplicherai.
Ecco queste mani, puoi usarle se lo vuoi,
per dividere nel mondo il pane che Tu hai dato a noi.
\endverse
\beginverse
Solo una goccia hai messo fra le mani mie,
solo una goccia che Tu ora chiedi a me.
Una goccia che, in mano a Te, una pioggia diventerà
e la terra feconderà.
\endverse
\beginverse
Sulle strade il vento da lontano porterà
il profumo del frumento, che tutti avvolgerà.
E sarà l'amore che il raccolto spartirà,
e il miracolo del pane in terra si ripeterà.
\endverse
\beginverse
Le nostre gocce, pioggia tra le mani Tue,
saranno linfa di una nuova civiltà.
E la terra preparerà la festa del pane che
ogni uomo condividerà.
\endverse
\endsong

% Fonte: LibrettoAccordi.pdf (Canti Coro), G12.
% Nell'elenco: Offertorio, Ordinario.
\beginsong{Frutto della nostra terra}[by={F. Buttazzo}]
\beginverse
Frutto della nostra terra, del lavoro di ogni uomo:
pane della nostra vita, cibo della quotidianità.
Tu che lo prendevi un giorno, lo spezzavi per i tuoi,
oggi vieni in questo pane, cibo vero dell'umanità.
\endverse
\beginchorus
E sarò pane, e sarò vino, nella mia vita, nelle tue mani.
Ti accoglierò dentro di me farò di me un'offerta viva,
un sacrificio gradito a te.
\endchorus
\beginverse
Frutto della nostra terra, del lavoro di ogni uomo:
vino delle nostre vigne sulla mensa dei fratelli tuoi.
Tu che lo prendevi un giorno, lo bevevi con i tuoi,
oggi vieni in questo vino e ti doni per la vita mia.
\endverse
\endsong

% Fonte: testo fornito dall'utente (tolti accordi e stacchetti musicali).
% Nell'elenco: "I frutti che ti offriamo". Finale: ultimi due versi ripetuti.
\beginsong{I frutti che ti offriamo}[by={F. Volpintesta}]
\beginverse
Pane bianco della terra qui per te, Signore,
ed il frutto della vite ti doniamo lieti:
la tua forza benedica questi frutti che,
con l'impegno di ogni giorno,
noi offriamo a te.
\endverse
\beginverse
Ogni ora di lavoro, ogni affanno in più,
il cercare ogni giorno la tua volontà
sono istanti della vita, sono frutti che,
con l'impegno di ogni giorno,
noi offriamo a te.
\endverse
\beginverse
Tanti sbagli da bambino e la nostalgia,
una gioia, un dolore, quella povertà:
tu li accetti, o Signore, sono frutti che,
con l'impegno di ogni giorno,
noi offriamo a te.
\endverse
\beginverse*
Con l'impegno di ogni giorno
noi offriamo a te.
\endverse
\endsong

% Fonte: LibrettoAccordi.pdf (Canti Coro), G1.
% Nell'elenco: Offertorio, Ordinario (autore in elenco: Gen Verde).
\beginsong{Il seme del tuo campo}[by={Gen Verde; testo P. A. Sequeri}]
\beginverse
Per ogni volta che ci doni la parola di Luce, noi offriremo la pace.
Per ogni volta che ci nutre il tuo Pane di vita, noi sazieremo la fame.
Per ogni volta che ci allieta il tuo vino di gioia, noi guariremo ferite.
\endverse
\beginchorus
Offriamo a te, sinceramente, la vita, benediciamo la tua pace fra noi.
Saremo l'eco del tuo canto, il seme del tuo campo,
il lievito del tuo perdono, il lievito del tuo perdono.
\endchorus
\beginverse
Non ci separa dalla fede l'incertezza del cuore, quando ci parli, Signore.
Non ci separa dall'amore la potenza del male, quando rimani con noi.
Non ci separa dall'attesa del tuo giorno la morte, quando ci tieni per mano.
\endverse
\endsong

% Fonte: LibrettoAccordi.pdf (Canti Coro), G16; c'e' anche "6 - Offertorio/06 Quelli che amano Te.pdf".
% Nell'elenco: Offertorio, Ordinario. CD: Voce del mio canto.
\beginsong{Quelli che amano te}[by={Gen Rosso (CD Voce del mio canto)}]
\beginverse
Ecco il frumento che noi ti portiamo,
pane della madre terra
e del lavoro e sudore di tanti,
quelli che amano te.
\endverse
\beginverse
Portiamo il vino, frutto della vite,
sangue della madre terra
e del lavoro e la gioia di tanti,
quelli che amano te.
\endverse
\beginverse
Insieme ai doni veniamo noi stessi,
figli di tutta la terra,
per riscaldarci attorno alla mensa
che hai preparato per noi.
\endverse
\endsong

% Fonte: LibrettoAccordi.pdf (Canti Coro), G5.
% Nell'elenco: Offertorio, Ordinario. CD: Come fuoco vivo.
\beginsong{Segni del tuo amore}[by={Gen Rosso, Gen Verde (CD Come fuoco vivo)}]
\beginverse
Mille e mille grani nelle spighe d'oro
mandano fragranza e danno gioia al cuore,
quando, macinati, fanno un pane solo:
pane quotidiano, dono tuo, Signore.
\endverse
\beginchorus
Ecco il pane e il vino, segni del tuo amore.
Ecco questa offerta, accoglila, Signore:
tu di mille e mille cuori fai un cuore solo, un corpo solo in te
e il Figlio tuo verrà, vivrà ancora in mezzo a noi.
\endchorus
\beginverse
Mille grappoli maturi sotto il sole,
festa della terra, donano vigore,
quando da ogni perla stilla il vino nuovo:
vino della gioia, dono tuo, Signore.
\endverse
\endsong

% Fonte: LibrettoAccordi.pdf (Canti Coro), G14.
% Nell'elenco: Offertorio, Ordinario. CD: Giorno senza tramonto (2013).
\beginsong{Seme della terra}[by={M. Balduzzi (CD Giorno senza tramonto)}]
\beginverse
Seme chiuso nella terra, sotto il manto della neve,
dorme, e morendo forse non lo sa, spiga alta verso il sole diverrà.
\endverse
\beginverse
Uomo curvo sulla terra sotto il peso della vita,
grida, e morendo forse non lo sa, cielo alto come il sole diverrà.
\endverse
\beginverse
Uomo, seme della terra, sotto il manto delle stelle,
ama, e morendo forse non lo sa, vive un istante di immortalità.
\endverse
\endsong

% Fonte: LibrettoAccordi.pdf (Canti Coro), G3.
% Nell'elenco: Offertorio, Ordinario.
\beginsong{Servo per amore}[by={Gen Rosso}]
\beginverse
Una notte di sudore sulla barca in mezzo al mare
e mentre il cielo si imbianca già, tu guardi le tue reti vuote.
Ma la voce che ti chiama un altro mare ti mostrerà
e sulle rive di ogni cuore le tue reti getterai.
\endverse
\beginchorus
Offri la vita tua
come Maria ai piedi della croce
e sarai servo di ogni uomo, servo per amore, sacerdote dell'umanità.
\endchorus
\beginverse
Avanzavi nel silenzio fra le lacrime e speravi
che il seme sparso davanti a te cadesse sulla buona terra.
Ora il cuore tuo è in festa perché il grano biondeggia ormai,
è maturato sotto il sole, puoi riporlo nei granai.
\endverse
\endsong

% Fonte: LibrettoAccordi.pdf (Canti Coro), G13.
% Nell'elenco: Offertorio, Ordinario.
\beginsong{Signore della vita}[by={D. Ricci}]
\beginverse
Signore della vita, Padre. Benedetto sei tu, Dio dell'universo,
dalla tua bontà abbiamo ricevuto il pane,
dalle spighe che ci doni, dalle mani che ci hai dato.
E questo pane noi lo portiamo a te, Signore, accoglilo.
Diventerà il cibo della vita eterna, il corpo tuo, Signore.
\endverse
\beginverse
Signore della vita, Padre. Benedetto sei tu, Dio dell'universo,
dalla tua bontà abbiamo ricevuto il vino,
dalle viti che ci doni, dalle mani che ci hai dato.
E questo vino noi lo portiamo a te, Signore, accoglilo.
Diventerà la bevanda di salvezza, il sangue tuo, Signore.
\endverse
\beginverse*
Signore della vita, Padre. Benedetto sei tu, Dio dell'universo.
\endverse
\endsong

% Fonte: LibrettoAccordi.pdf (Canti Coro), G2. Omesso l'intro "Na, na, na...".
% Nell'elenco: Offertorio, Matrimonio.
\beginsong{Vita nuova nell'amore}[by={F. Buttazzo}]
\beginverse
Vita nuova germogliata in noi, come grano che ha piantato Dio,
come il pane per l'umanità: è l'amore suo, è l'amore in noi.
Vita nuova di felicità, vino buono per l'umanità,
con la forza dello Spirito sarà via di santità. \rep{2}
\endverse
\beginverse
È l'amore eterno di Gesù, per la Chiesa, amata sposa sua,
il mistero che si celebra, che risplende qui, che risplende in noi.
Il Signore ha chiamato noi all'amore nella fedeltà,
benedice i nostri passi qui nella sua comunità. \rep{2}
\endverse
\beginverse
È una casa che si costruirà sulla roccia della verità,
fondamento che non cede mai: il Signore in noi, il Signore in noi.
È famiglia che feconderà questo mondo di vitalità;
è promessa di felicità, vita piena data a noi. \rep{2}
\endverse
\endsong

% Fonte: LibrettoAccordi.pdf (Canti Coro), G4. Tolte le indicazioni Uomini/Donne.
% Nell'elenco: Offertorio, Ordinario.
\beginsong{Voi siete di Dio}[by={M. Balduzzi (CD Verbum Panis)}]
\beginverse
Tutte le stelle della notte,
le nebulose e le comete,
il sole su una ragnatela, è tutto vostro e voi siete di Dio.
\endverse
\beginverse
Tutte le rose della vita, il grano, i prati, i fili d'erba,
il mare, i fiumi, le montagne, è tutto vostro e voi siete di Dio.
\endverse
\beginverse
Tutte le musiche e le danze, i grattacieli, le astronavi,
i quadri, i libri, le culture, è tutto vostro e voi siete di Dio.
\endverse
\beginverse
Tutte le volte che perdono, quando sorrido, quando piango,
quando mi accorgo di chi sono, è tutto vostro e voi siete di Dio.
È tutto nostro e noi siamo di Dio.
\endverse
\endsong

\end{songs}

% !TEX root = ../libretto.tex
% Sezione: Pace e Agnello di Dio
\momento{Pace e Agnello di Dio}
\begin{songs}{titoli}

% Fonte: "8 - Agnello di Dio/Agnello di Dio - RnS (Mi).pdf" (spartito; testo liturgico).
% Nell'elenco: Agnus Dei, Ordinario.
\beginsong{Agnello di Dio}[by={D. De Dominicis (RnS)}]
\beginverse*
Agnello di Dio, che togli i peccati del mondo,
abbi pietà di noi. \rep{2}
Agnello di Dio, che togli i peccati del mondo,
dona a noi la pace.
\endverse
\endsong

% Fonte: LibrettoAccordi.pdf (Canti Coro), I1.
% Nell'elenco: Agnus Dei, Ordinario. Due voci sovrapposte: ostinato "Dona nobis pacem" e invocazione.
\beginsong{Agnus Dei}[by={M. Balduzzi (CD Verbum Panis)}]
\beginchorus
Dona nobis, dona nobis pacem.
\endchorus
\beginverse*
Agnus Dei, qui tollis peccata mundi,
miserere, miserere nobis. \rep{2}
\endverse
\endsong

\end{songs}

% !TEX root = ../libretto.tex
% Sezione: Comunione
\momento{Comunione}
\begin{songs}{titoli}

% Fonte: "9 - Comunione/ALTISSIMO - Spoladore.pdf" (foglio accordi).
% Nell'elenco: senza tipologia.
\beginsong{Altissimo}[by={P. Spoladore}]
\beginverse
Altissimo, onnipotente, buon Signore,
tue son le lodi, la gloria, l'onore ed ogni benedizione
che a te solo e al tuo nome, Altissimo, possiamo elevare
e nessun uomo può credersi degno di poterti nominare.
\endverse
\beginverse
Laudato sii, mi Signore, con tutte le tue creature,
specialmente per frate sole così bello e radioso,
con la sua luce illumini il giorno ed illumini noi
e con grande splendore ci parla di te, Signore.
\endverse
\beginchorus
Lodate e benedite il Signore,
ringraziate e servite con grande umiltà,
lodate, lodate, benedite il Signore con grande umiltà,
ringraziate e servite con grande umiltà.
\endchorus
\beginverse
Laudato sii, mi Signore, per sora luna e le stelle,
le hai formate nel cielo così chiare, preziose e belle,
per frate vento, per l'aria, il sereno ed ogni tempo,
così la vita tu cresci e sostieni in ogni tua creatura.
\endverse
\beginverse
Laudato sii, mi Signore, per sora acqua così preziosa,
per frate fuoco giocoso e potente che c'illumina la notte,
laudato sii, mi Signore, per sora nostra madre la terra,
ci sostiene, governa e ci dona fiori, frutti ed erba.
\endverse
\beginverse
Laudato sii, mi Signore, per quelli che per il tuo amore
perdonano e sopportano in pace ogni persecuzione,
che sora morte ha trovato viventi nella tua volontà,
da te, Altissimo, un giorno saranno da te incoronati.
\endverse
\endsong

% Fonte: LibrettoAccordi.pdf (Canti Coro), J15.
% Nell'elenco: Comunione, Pentecoste; l'utente lo vuole tra i canti di Comunione.
\beginsong{Come fuoco vivo}[by={Gen Rosso, Gen Verde (CD Come fuoco vivo)}]
\beginchorus
Come fuoco vivo si accende in noi un'immensa felicità,
che mai più nessuno ci toglierà, perché tu sei ritornato.
Chi potrà tacere, da ora in poi, che sei tu in cammino con noi,
che la morte è vinta per sempre, che ci hai ridonato la vita.
\endchorus
\beginverse
Spezzi il pane davanti a noi, mentre il sole è al tramonto:
ora gli occhi ti vedono, sei tu! Resta con noi.
\endverse
\beginverse
E per sempre ti mostrerai in quel gesto d'amore:
mani che ancora spezzano pane d'eternità.
\endverse
\endsong

% Fonte: LibrettoAccordi.pdf (Canti Coro), A13.
% Nell'elenco: Comunione, Ordinario.
\beginsong{Come l'aurora verrai}[by={Gen Verde}]
\beginverse
Come l'aurora verrai, le tenebre in luce cambierai
tu per noi, Signore.
Come la pioggia cadrai, sui nostri deserti scenderai,
scorrerà l'amore.
\endverse
\beginchorus
Tutti i nostri sentieri percorrerai,
tutti i figli dispersi raccoglierai,
chiamerai da ogni terra il tuo popolo,
in eterno ti avremo con noi.
\endchorus
\beginverse
Re di giustizia sarai, le spade in aratri forgerai:
ci darai la pace.
Lupo ed agnello vedrai, insieme sui prati dove mai
tornerà la notte.
\endverse
\beginverse
Dio di salvezza tu sei e come una stella sorgerai
su di noi per sempre.
E chi non vede, vedrà, chi ha chiusi gli orecchi sentirà,
canterà di gioia.
\endverse
\endsong

% Fonte: LibrettoAccordi.pdf (Canti Coro), J36.
% Nell'elenco: Comunione, Ordinario.
\beginsong{Dietro un semplice volto}[by={M. Balduzzi (CD Giorno senza tramonto)}]
\beginchorus
Ogni volta voi lo farete a me, dietro il volto di un uomo.
Ogni volta voi scoprirete me, dietro un semplice volto.
\endchorus
\beginverse
Ogni volta che darai una goccia d'acqua,
ogni volta che sarai seme di speranza.
\endverse
\beginverse
Ma quando mai, Signore, ti abbiam visto in chi ha sete e fame?
Quando ti abbiamo visto sfigurato in chi non ha nessuno?
\endverse
\beginverse
Ogni volta che sarai una mano tesa,
ogni volta che darai forza di lottare.
\endverse
\beginverse
Ma quando mai, Signore, ti abbiam visto in chi non ha più casa?
Quando ti abbiamo visto disperato in chi chiedeva aiuto?
\endverse
\endsong

% Fonte: testo fornito dall'utente (tolti gli accordi). Non presente nell'elenco.
% Brano che dà il titolo al CD "Dove tu sei" (1982). La strofa si canta due volte, poi il finale.
\beginsong{Dove tu sei}[by={Gen Rosso (CD Dove tu sei)}]
\beginverse*
Dove tu sei torna la vita,
dove tu passi fiorisce il deserto,
dove tu guardi si rischiara il cielo
e in fondo al cuore torna il sereno.
Dove tu sei,
dove tu sei. \rep{2}
\endverse
\beginverse*
Dove tu sei
torna la vita.
\endverse
\endsong

% Fonte: LibrettoAccordi.pdf (Canti Coro), J23.
% Nell'elenco: Comunione, Ordinario.
\beginsong{E sono solo un uomo}[by={P. A. Sequeri}]
\beginverse
Io lo so, Signore, che vengo da lontano,
prima nel pensiero e poi nella tua mano,
io mi rendo conto che Tu sei la mia vita,
e non mi sembra vero di pregarti così.
\endverse
\beginverse
Padre di ogni uomo, e non ti ho visto mai,
spirito di vita, e nacqui da una donna,
figlio, mio fratello, e sono solo un uomo,
eppure io capisco che Tu sei verità.
\endverse
\beginchorus
E imparerò a guardare tutto il mondo
con gli occhi trasparenti di un bambino,
e insegnerò a chiamarti Padre nostro
ad ogni figlio che diventa uomo. \rep{2}
\endchorus
\beginverse
Io lo so Signore che Tu mi sei vicino,
Luce alla mia mente guida al mio cammino,
mano che sorregge, sguardo che perdona
e non mi sembra vero che Tu esista così.
\endverse
\beginverse
Dove nasce amore, Tu sei la Sorgente;
dove c'è una croce, Tu sei la Speranza;
dove il tempo ha fine, Tu sei Vita eterna,
e so che posso sempre contare su di Te.
\endverse
\beginchorus
E accoglierò la vita come un dono,
e avrò il coraggio di morire anch'io,
e incontro a Te verrò col mio fratello
che non si sente amato da nessuno. \rep{2}
\endchorus
\endsong

% Fonte: LibrettoAccordi.pdf (Canti Coro), J10.
% Nell'elenco: Comunione, Ordinario.
\beginsong{Fratello sole, sorella luna}[by={R. Ortolani}]
\beginverse
Dolce sentire come nel mio cuore,
ora umilmente sta nascendo amore.
Dolce capire che non son più solo
ma che son parte di una immensa vita,
che generosa risplende intorno a me;
dono di Lui del suo immenso amor.
\endverse
\beginverse
Ci ha dato il cielo e le chiare stelle,
fratello sole e sorella luna;
la madre terra, con frutti, prati e fiori,
il fuoco, il vento, l'aria e l'acqua pura,
fonte di vita per le sue creature,
dono di Lui del suo immenso amor. \rep{2}
\endverse
\endsong

% Fonte: LibrettoAccordi.pdf (Canti Coro), J31.
% Nell'elenco: Comunione, Ordinario.
\beginsong{Insieme a te}[by={F. Buttazzo}]
\beginchorus
Insieme a te, uniti a te,
la nostra vita si trasformerà.
Insieme a noi, accanto a noi,
ti sentiremo ogni giorno Gesù.
\endchorus
\beginverse
Spezzando il pane hai detto ai tuoi:
«Questo è il mio corpo donato a voi»,
prendendo il calice hai detto Gesù:
«Ecco il mio sangue versato per voi».
\endverse
\beginverse
Se celebriamo la Pasqua con te,
diventeremo discepoli tuoi,
tu hai donato la vita per noi,
perché viviamo in eterno con te.
\endverse
\beginverse
Tu hai mandato i discepoli tuoi,
in tutto il mondo a parlare di te,
a rinnovare il tuo gesto d'amore:
«Fate questo in memoria di me».
\endverse
\beginverse
E con la forza che viene da te,
cammineremo nel mondo Signor,
con questo pane che hai dato a noi,
riceveremo la vita di Dio.
\endverse
\endsong

% Fonte: LibrettoAccordi.pdf (Canti Coro), J20.
% Nell'elenco: Comunione, Ordinario.
\beginsong{L'unico maestro}[by={M. Civico}]
\beginverse
Le mie mani, con le tue, possono fare meraviglie,
possono stringere e perdonare e costruire cattedrali.
Possono dare da mangiare e far fiorire una preghiera.
\endverse
\beginchorus
Perché tu, solo tu, solo tu sei il mio maestro e insegnami
ad amare come hai fatto tu con me,
se lo vuoi, io lo grido a tutto il mondo che tu sei,
l'unico maestro sei per me.
\endchorus
\beginverse
Questi occhi, con i tuoi, potranno vedere meraviglie,
potranno piangere e luccicare, guardare oltre ogni frontiera.
Potranno amare più di ieri se sanno insieme a te sognare.
\endverse
\beginverse
Questi piedi, con i tuoi possono fare strade nuove,
possono correre e riposare, sentirsi a casa in questo mondo.
Possono mettere radici e passo a passo camminare.
\endverse
\beginverse
Tu sei il capo, noi le membra, diciamo un'unica preghiera,
Tu sei il maestro noi, testimoni della parola del vangelo.
Possiamo vivere felici in questa chiesa che rinasce.
\endverse
\endsong

% Fonte: testo fornito dall'utente (in Canti Coro c'e' solo l'immagine "1 Canti di Ingresso/Le legge della vita.jpg").
% Nell'elenco: Comunione.
\beginsong{La legge della vita}[by={Gen Rosso}]
\beginverse*
C'è una legge vera nella vita, impressa in ogni cosa:
legge che muove gli astri del cielo in un concerto d'armonia;
e in terra canta nei colori della natura, canta nella natura.
\endverse
\beginverse*
Il giorno cede il passo alla notte per amore,
la notte saluta il giorno per amore.
Dal mare sale l'acqua al cielo per amore,
e l'acqua discende dal cielo al mare per amore.
La pianta dà le foglie alla terra per amore,
la terra ridona le foglie per amore,
un seme cade in terra e muore per amore,
la vita germoglia dal solco della morte per amore!
\endverse
\beginchorus
Questa è la legge eterna, legge di Dio,
un Dio che per amore ha creato ogni cosa
ed ha nascosto amore, ed ha nascosto amore,
dietro apparenze di morte e di dolore.
\endchorus
\beginverse*
È la legge vera della vita impressa in ogni cosa,
legge che muove gli astri del cielo in un concerto d'armonia;
e in terra canta nei colori della natura, canta nella natura. \rep{2}
\endverse
\endsong

% Fonte: LibrettoAccordi.pdf (Canti Coro), A3.
% Nell'elenco: Ingresso, Ordinario (RnS); confermato dall'utente: Comunione, M. Balduzzi.
\beginsong{Nel tuo amore siamo un popolo}[by={M. Balduzzi (CD Xaipe)}]
\beginverse
Come note di una musica, pagine di un solo libro,
come il suono della pioggia, alberi di una foresta,
senza più barriere, senza più stranieri,
nel tuo amore siamo un popolo.
\endverse
\beginverse
Dove anche solo due o tre sono uniti nel tuo amore,
dove il pane della vita hai diviso insieme a noi,
senza più barriere, senza più stranieri,
nel tuo amore siamo un popolo.
\endverse
\beginverse
Quando il seme della pace spunta lieve sulla terra,
quando il soffio del tuo vento muove i giorni della storia,
senza più barriere, senza più stranieri,
nel tuo amore siamo un popolo. \rep{2}
\endverse
\endsong

% Fonte: LibrettoAccordi.pdf (Canti Coro), J8.
% Nell'elenco: Comunione, Ordinario.
\beginsong{Nel tuo silenzio}[by={Gen Verde}]
\beginverse
Nel tuo silenzio accolgo il mistero venuto a vivere dentro di me.
Sei tu che vieni, o forse è più vero che tu mi accogli in te, Gesù.
\endverse
\beginverse
Sorgente viva che nasce dal cuore, è questo dono che abita in me.
La tua presenza è un Fuoco d'amore che avvolge l'anima mia, Gesù.
\endverse
\beginverse
Ora il tuo Spirito in me dice: «Padre», non sono io a parlare, sei tu.
Nell'infinito oceano di pace Tu vivi in me, io in te, Gesù.
\endverse
\endsong

% Fonte: LibrettoAccordi.pdf (Canti Coro), J34.
% Nell'elenco: Comunione, Ordinario.
\beginsong{Nelle tue mani\soloindice{ (RnS)}}[by={RnS}]
\beginverse
Nelle tue mani affido la vita, Dio, mia salvezza sei tu.
Volgi il tuo sguardo al mio cuore, con te al sicuro sarò.
\endverse
\beginchorus
Padre del cielo, per il tuo nome vivrò,
un sacrificio con la mia lode io ti offrirò,
per tutto ciò che sempre mi dai.
Padre del cielo, per il tuo nome vivrò,
scruta il mio cuore e la mia lode io ti offrirò,
per tutto ciò che sempre mi donerai,
o Dio di ogni bontà.
\endchorus
\beginverse
Nelle tue mani è la mia vita, Dio, mia speranza sei tu.
Donami pace, o Signore, con te al sicuro vivrò.
\endverse
\endsong

% Fonte: LibrettoAccordi.pdf (Canti Coro), J5.
% Nell'elenco: Comunione (anche Ingresso).
\beginsong{Ogni mia parola}[by={Gen Verde}]
\beginverse
Come la pioggia e la neve
scendono giù dal cielo
e non vi ritornano senza irrigare
e far germogliare la terra,
così ogni mia parola non ritornerà a me
senza operare quanto desidero,
senza aver compiuto ciò per cui l'avevo mandata.
Ogni mia parola, ogni mia parola.
\endverse
\endsong

% Fonte: LibrettoAccordi.pdf (Canti Coro), J7.
% Nell'elenco: Comunione, Ordinario. CD: Dove tu sei (1982).
\beginsong{Padre mio}[by={Gen Rosso (CD Dove tu sei)}]
\beginverse
Padre mio, m'abbandono a te, di me fa' quello che ti piace;
grazie di ciò che fai per me,
spero solamente in te.
Purché si compia il tuo volere in me e in tutti i miei fratelli
niente desidero di più:
fare quello che vuoi tu.
\endverse
\beginchorus
Dammi che ti riconosca, dammi che ti possa amare sempre più.
Dammi che ti resti accanto, dammi d'essere l'amor.
\endchorus
\beginverse
Fra le tue mani depongo la mia anima, con tutto l'amore del mio cuore:
mio Dio, la dono a te,
perché ti amo immensamente.
Sì, ho bisogno di donarmi a te senza misura, affidarmi alle tue mani,
perché sei il Padre mio,
perché sei il Padre mio.
\endverse
\endsong

% Fonte: testo fornito dall'utente.
% Nell'elenco: Comunione, Ordinario. CD: Dove tu sei (1982).
\beginsong{Pane del cielo}[by={Gen Rosso (CD Dove tu sei)}]
\beginchorus
Pane del cielo sei tu, Gesù,
via d'amore: tu ci fai come te. \rep{2}
\endchorus
\beginverse
No, non è rimasta fredda la terra:
tu sei rimasto con noi
per nutrirci di te,
Pane di vita,
ed infiammare col tuo amore
tutta l'umanità.
\endverse
\beginverse
Sì, il cielo è qui su questa terra:
tu sei rimasto con noi
ma ci porti con te
nella tua casa,
dove vivremo insieme a te
tutta l'eternità.
\endverse
\beginverse
No, la morte non può farci paura:
tu sei rimasto con noi
e chi vive in te
vive per sempre.
Sei Dio per noi, sei Dio con noi,
Dio in mezzo a noi.
\endverse
\endsong

% Fonte: immagine del testo fornita dall'utente.
% Nell'elenco: Comunione, Ordinario.
\beginsong{Pane di vita}[by={Parisi}]
\beginverse
Pane di vita tu sei, Signore,
pane del cielo che nutre il cammino;
noi che veniamo alla mensa del regno
oggi con fede annunciamo il mistero.
\endverse
\beginverse
Questo è il mio corpo: prendete e mangiate;
questo è il mio sangue: prendete e bevete;
voi che accogliete la mia parola
della mia vita in eterno vivrete.
\endverse
\beginverse
Dio d'amore tu sei, o Signore,
nuova alleanza e fonte di vita;
noi ti accogliamo, o Dono del Padre,
noi ti portiamo una sete infinita.
\endverse
\beginverse
Siamo il tuo corpo che vive nel tempo,
ma non ci basta l'intero universo;
Cristo Signore, tu sei nutrimento,
placa la fame di gioia e d'eterno.
\endverse
\beginverse
Vieni tra noi e rinnova l'incontro,
pace rinnova tra il cielo e la terra;
ad ogni cosa ridona bellezza,
ad ogni uomo ridona amicizia.
\endverse
\beginverse
Figli del giorno chiamati alla luce,
nel tuo splendore viviamo, Signore;
tu sei il sole che vince la notte,
tu sei la vita che vince la morte.
\endverse
\beginverse
Popolo santo in cammino nel mondo,
tempio che accoglie la tua presenza
noi ti cantiamo la lode del cuore:
tu sei, Signore, la nostra speranza.
\endverse
\endsong

% Fonte: testo fornito dall'utente.
% Nell'elenco: Comunione, Ordinario.
\beginsong{Pane di vita nuova}[by={M. Frisina}]
\beginverse
Pane di vita nuova,
vero cibo dato agli uomini,
nutrimento che sostiene il mondo,
dono splendido di grazia.
\endverse
\beginverse
Tu sei sublime frutto
di quell'albero di vita
che Adamo non poté toccare:
ora è in Cristo a noi donato.
\endverse
\beginchorus
Pane della vita,
sangue di salvezza,
vero corpo, vera bevanda,
cibo di grazia per il mondo.
\endchorus
\beginverse
Sei l'agnello immolato
nel cui sangue è la salvezza,
memoriale della vera Pasqua
della nuova alleanza.
\endverse
\beginverse
Manna che nel deserto
nutri il popolo in cammino,
sei sostegno e forza nella prova
per la Chiesa in mezzo al mondo.
\endverse
\beginverse
Vino che ci dà gioia,
che riscalda il nostro cuore,
sei per noi il prezioso frutto
della vigna del Signore.
\endverse
\beginverse
Dalla vite ai tralci
scorre la vitale linfa
che ci dona la vita divina,
scorre il sangue dell'amore.
\endverse
\beginverse*
Amen.
\endverse
\endsong

% Testo fornito dall'utente. Parafrasi del Salmo 23 (22).
% Album: "Cerco il tuo volto" (2000).
% Struttura cantata: Rit. - Strofa - Rit. - ... (ritornello dopo ogni strofa).
% Nell'elenco: Comunione, Ordinario.
\beginsong{Perché tu sei con me}[by={Gen Verde}]
\beginchorus
Solo tu sei il mio pastore,
niente mai mi mancherà.
Solo tu sei il mio pastore, o Signore (o Signore).
\endchorus
\beginverse
Mi conduci dietro te sulle verdi alture,
ai ruscelli tranquilli lassù,
dov'è più limpida l'acqua per me,
dove mi fai riposare.
\endverse
\beginverse
Anche fra le tenebre d'un abisso oscuro
io non temo alcun male perché
tu mi sostieni, sei sempre con me,
rendi il sentiero sicuro.
\endverse
\beginverse
Siedo alla tua tavola che mi hai preparato
ed il calice è colmo per me
di quella linfa di felicità
che per amore hai versato.
\endverse
\beginverse
Sempre mi accompagnano lungo estati e inverni
la tua grazia, la tua fedeltà.
Nella tua casa io abiterò
fino alla fine dei giorni.
\endverse
\endsong

% Fonte: LibrettoAccordi.pdf (Canti Coro), J13; c'e' anche "9 - Comunione/PREGHIERA - mgc D.pdf".
% Omessa la parte finale a due voci. Nell'elenco: Comunione, Ordinario.
\beginsong{Preghiera}[by={Movimento Giovanile Costruire}]
\beginverse
Voglio ringraziarti, mio Signore, perché sei oggi ancora con me.
C'è tanta sete nel mio cuore, dammi l'acqua del tuo amore.
Sei Tu la mia certezza, sei Tu che fai sperare che
ancora il seme germoglierà, una nuova terra rivivrà,
un lavoro nuovo nascerà, tanta gente finalmente sorriderà,
un bambino forte ti pregherà, ti ringrazierà, o mio Signore.
\endverse
\beginverse
Anche quando il buio è dentro me, ed io non vedo la tua luce,
so che Tu aspetti là, e mi dai ancora la tua pace.
Fa', Signore, che non manchi mai, fa' che resti qui con noi,
facci uniti nel tuo nome, dacci ancora del tuo pane,
effondi amore su ogni cuore, dona gioia a chi è nel dolore.
Ancora il seme germoglierà, una nuova terra rivivrà,
un lavoro nuovo nascerà, tanta gente finalmente sorriderà,
ed ognuno di noi risorgerà, ed ognuno di noi per te sarà.
\endverse
\beginchorus
Preghiera, per Te sarà, preghiera. \rep{2}
\endchorus
\endsong

% Fonte: LibrettoAccordi.pdf (Canti Coro), J18.
% Nell'elenco: Comunione (autore RnS). CD: Cantiamo con gioia (1998).
\beginsong{Re di gloria}[by={F. Marranzino, A. De Luca (RnS) (CD Cantiamo con gioia)}]
\beginverse
Ho incontrato te, Gesù, e ogni cosa in me è cambiata,
tutta la mia vita ora ti appartiene.
Tutto il mio passato io lo affido a te,
Gesù, Re di gloria, mio Signor.
\endverse
\beginchorus
Tutto in te riposa, la mia mente, il mio cuore,
trovo pace in te, Signor, tu mi dai la gioia.
Voglio stare insieme a te, non lasciarti mai,
Gesù, Re di gloria, mio Signor.
\endchorus
\beginverse
Dal tuo amore chi mi separerà? Sulla croce hai dato la vita per me.
Una corona di gloria mi darai quando un giorno ti vedrò.
\endverse
\beginverse*
Io ti aspetto, mio Signor!
Io ti aspetto, mio Signor!
Io ti aspetto, mio Re!
\endverse
\endsong

% Fonte: LibrettoAccordi.pdf (Canti Coro), O22.
% Nell'elenco: Comunione, Pasqua; l'utente lo vuole tra i canti di Comunione. CD: Il mistero pasquale (2011).
\beginsong{Resto con te}[by={Gen Verde (CD Il mistero pasquale)}]
\beginverse
Seme gettato nel mondo,
Figlio donato alla terra,
il tuo silenzio custodirò.
In ciò che vive e che muore
vedo il tuo volto d'amore:
sei il mio Signore e sei il mio Dio.
\endverse
\beginchorus
Io lo so che tu sfidi la mia morte.
Io lo so che tu abiti il mio buio.
Nell'attesa del giorno che verrà
resto con te.
\endchorus
\beginverse
Nube di mandorlo in fiore
dentro gli inverni del cuore
è questo pane che tu ci dai.
Vena di cielo profondo
dentro le notti del mondo
è questo vino che tu ci dai.
\endverse
\beginverse*
Tu sei Re di stellate immensità
e sei tu il futuro che verrà,
sei l'amore che muove ogni realtà
e tu sei qui.
Resto con te.
\endverse
\endsong

% Fonte: LibrettoAccordi.pdf (Canti Coro), O20 (li' in Settimana Santa).
% Nell'elenco: Comunione, Ordinario (due righe). CD: Dove tu sei (1982).
\beginsong{Sei Dio}[by={Gen Rosso (CD Dove tu sei)}]
\beginverse
Perché avessimo la luce, ti facesti buio.
Perché avessimo la vita tu provasti la morte.
Ci basta, Signore, vederci simili a te
e offrire col tuo il nostro dolore.
\endverse
\beginchorus
Sei Dio, sei il mio Dio,
il nostro Dio
d'amore infinito. \rep{2}
\endchorus
\beginverse
Perché fossimo innocenti ti sentisti peccato.
Perché fosse nostro il cielo fosti abbandonato.
Per darci, Signore, la vera vita qui in terra,
sembrò che il Padre fosse lontano da te.
\endverse
\endsong

% Fonte: testo fornito dall'utente (tolti accordi e intermezzo). Non presente nell'elenco.
% Finale: ritornello ripetuto un tono sopra.
\beginsong{Sei il Pane della vita}[by={F. Volpintesta}]
\beginverse
A te veniamo per dirti grazie,
solo chiediamo il tuo amor,
e nasce un canto che ci innamora:
la nostra vita doniamo a te.
Sei tu soltanto il nostro bene,
vogliam restare accanto a te.
La vita è un volo che giunge a sera:
nessuno è solo, se spera in te.
\endverse
\beginchorus
Sei il Pane della vita, sei l'Amore, tu.
Sei il Giorno che non muore, sei la Vita, tu.
\endchorus
\beginverse
Ogni dolore nel tuo dolor,
ogni fatica l'hai presa tu,
nel tuo silenzio un solo canto:
ogni respiro sarà per te.
\endverse
\beginchorus
Sei il Pane della vita, sei l'Amore, tu.
Sei il Giorno che non muore, sei la Vita, tu. \rep{2}
\endchorus
\endsong

% Fonte: LibrettoAccordi.pdf (Canti Coro), J16.
% Nell'elenco: Comunione, Ordinario. CD: In cerca d'autore (1977).
\beginsong{Symbolum '77}[by={P. A. Sequeri (CD In cerca d'autore)}]
\beginverse
Tu sei la mia vita, altro io non ho; Tu sei la mia strada, la mia verità.
Nella tua parola io camminerò
finché avrò respiro, fino a quando Tu vorrai.
Non avrò paura, sai, se Tu sei con me; io Ti prego resta con me.
\endverse
\beginverse
Credo in Te, Signore, nato da Maria, Figlio eterno e santo, uomo come noi,
morto per amore, vivo in mezzo a noi,
una cosa sola con il Padre e con i tuoi;
fino a quando, io lo so, Tu ritornerai per aprirci il Regno di Dio.
\endverse
\beginverse
Tu sei la mia forza, altro io non ho; Tu sei la mia pace, la mia libertà.
Niente nella vita ci separerà.
So che la tua mano forte non mi lascerà,
so che da ogni male Tu mi libererai; e nel tuo perdono vivrò.
\endverse
\beginverse
Padre della vita, noi crediamo in Te; Figlio Salvatore, noi speriamo in Te;
Spirito d'amore, vieni in mezzo a noi.
Tu da mille strade ci raduni in unità
e per mille strade, poi, dove Tu vorrai, noi saremo il seme di Dio.
\endverse
\endsong

% Fonte: LibrettoAccordi.pdf (Canti Coro), J3; c'e' anche
% "Canti Coro/9 - Comunione/208TE, AL CENTRO DEL MIO CUORE.pdf".
% Nell'elenco: Comunione, Ordinario.
\beginsong{Te al centro del mio cuore}[by={Gen Verde}]
\beginverse
Ho bisogno di incontrarti nel mio cuore,
di trovare te, di stare insieme a te.
Unico riferimento del mio andare,
unica ragione tu, unico sostegno, tu;
al centro del mio cuore ci sei solo tu.
\endverse
\beginverse
Anche il cielo gira intorno e non ha pace,
ma c'è un punto fermo, è quella stella là:
la stella polare è fissa, ed è la sola:
la stella polare, tu, la stella sicura, tu;
al centro del mio cuore ci sei solo tu.
Tutto ruota intorno a te e in funzione di te,
e poi non importa il come, il dove, il se.
\endverse
\beginchorus
Che tu splenda sempre al centro del mio cuore,
il significato allora sarai tu.
Quello che farò sarà soltanto amore,
unico sostegno tu, la stella polare tu,
al centro del mio cuore ci sei solo tu.
\endchorus
\endsong

% Fonte: LibrettoAccordi.pdf (Canti Coro), J22.
% Nell'elenco: Comunione, Ordinario.
\beginsong{Tra voi la carità}[by={S. Zampa}]
\beginverse
Cos'è quest'aria che ci avvolge, questo mistero che ci supera,
e fa di noi un solo cuore, un'unica famiglia:
è l'amore di chi sa dare la vita, di chi non misura mai.
\endverse
\beginchorus
Ci sia tra voi la carità,
l'amore vero che mai finirà.
E tra la gente un fuoco che
il mondo intero trasformerà.
\endchorus
\beginverse
Sentirti ancora vivo tra di noi, al centro della nostra vita,
e ritrovare come nuova quell'unità più vera,
vivendo così, già sulla terra, un angolo di cielo.
\endverse
\endsong

% Fonte: LibrettoAccordi.pdf (Canti Coro), J25.
% Nell'elenco: Comunione, Ordinario.
\beginsong{Tu sei il pane di vita}[by={F. Baggio}]
\beginverse
Tu sei il pane di vita, sei la Parola di verità.
Tu sei luce del mondo, sei la certezza di eternità.
Tu, sollievo dell'anima, stella del nostro cammino.
Acqua dalla sorgente che sgorga dal cuore di Dio.
\endverse
\beginverse
Tu sei il nostro pastore che ci conduce lungo la via;
nulla c'è da temere se la tua guida mai mancherà.
Tu ci sveli nell'intimo il mistero di Dio.
Ci fai figli nel Figlio: noi, sacerdoti, profeti e re.
\endverse
\beginverse
Tu sei il nostro Maestro, sei la sapienza, unico Dio.
Sei la strada e la porta che apre il tempo all'eternità.
Tu, sostegno del debole, vincitore del male,
tu, ricchezza del povero ed eterna felicità.
\endverse
\beginverse
Tu sei pietra angolare, in te tutti siamo Chiesa di Dio.
Tu sei vite feconda e noi siamo i tralci uniti a te.
Sei l'amore più candido, sacrificio perfetto.
La bellezza che salva il mondo sei tu, Signore Gesù.
\endverse
\endsong

% Fonte: LibrettoAccordi.pdf (Canti Coro), J44.
% Nell'elenco: Comunione, Ordinario. CD: Alla tua mensa (2006).
\beginsong{Tu sei per noi}[by={P. Auricchio (CD Alla tua mensa)}]
\beginchorus
Tu sei per noi sorgente d'acqua viva,
Tu sei per noi la vita senza fine,
Tu sei per noi il pane del cielo,
Cristo Gesù Signore.
\endchorus
\beginverse
Pane del cielo, cibo di vita,
solo tu doni la gioia del cuore,
doni vigore, slancio nel dare,
spingi la Chiesa a gesti d'amore.
\endverse
\beginverse
Unico pane che nutre il mondo,
unico vino che dà speranza;
sei comunione che unisce i cuori,
sei comunione che ci fa uno.
\endverse
\beginverse
Verbo di Dio, Maestro e Signore,
guarda il tuo popolo sulle tue vie;
dona coraggio di scelte feconde,
mostra il tuo volto ad ogni uomo.
\endverse
\beginverse
Nella tua Chiesa risuona il Vangelo,
buona notizia che salva il mondo;
unico annuncio che infrange ogni buio,
vera speranza che illumina il cuore.
\endverse
\endsong

% Fonte: LibrettoAccordi.pdf (Canti Coro), J12. Omesse le indicazioni per il coro.
% Nell'elenco: Comunione, Ordinario.
\beginsong{Una città che ti assomiglia}[by={Gen Rosso}]
\beginverse
E mi fermai dove solo tu sai, dove al buio succede l'aurora, e mi calmai.
Non dubitai, niente al mondo vorrei che salire la cima del monte dove sei.
\endverse
\beginverse
Anche i passeri nel vento che volano hanno un nido ed al sicuro riposano.
Anche i piccoli del cervo hanno tane dove tornano.
\endverse
\beginchorus
E mi fermerò all'ombra dei tuoi portici;
abiterò in una grande città, che ti assomiglia e mi scalderà. \rep{2}
\endchorus
\beginverse*
E mi fermai dove solo tu sai, dove al buio succede l'aurora e mi calmai.
\endverse
\endsong

% Fonte: LibrettoAccordi.pdf (Canti Coro), J14.
% Nell'elenco: Comunione, Ordinario. CD: Verbum Panis (2000).
\beginsong{Verbum panis}[by={C. Casucci, M. Balduzzi (CD Verbum Panis)}]
\beginverse
Prima del tempo, prima ancora che la terra cominciasse a vivere
il Verbo era presso Dio.
Venne nel mondo, e per non abbandonarci in questo viaggio ci lasciò
tutto sé stesso come pane.
\endverse
\beginchorus
Verbum caro factum est, Verbum panis factum est. \rep{2}
\endchorus
\beginverse
Qui spezzi ancora il pane in mezzo a noi,
e chiunque mangerà non avrà più fame.
Qui vive la tua Chiesa intorno a Te,
dove ognuno troverà la sua vera casa.
\endverse
\beginverse
Prima del tempo, quando l'universo fu creato dall'oscurità
il Verbo era presso Dio.
Venne nel mondo, nella sua misericordia Dio ha mandato il Figlio suo,
tutto sé stesso come pane.
\endverse
\endsong

\end{songs}

% !TEX root = ../libretto.tex
% Sezione: Canti finali
\momento{Canti finali}
\begin{songs}{titoli}

% Fonte: LibrettoAccordi.pdf (Canti Coro), J26.
% Nell'elenco: Finale, Ordinario.
\beginsong{Al centro del presente}[by={Gen Rosso}]
\beginchorus
Oh, oh, oh, vivi il momento che va
cogli l'attimo adesso, eh, eh eh
Oh, oh, oh, vivi il momento che va
cogli l'attimo adesso.
\endchorus
\beginverse
Lui va a scuola, ma non sa perché, butta i giorni nella noia;
vive dentro una fotografia inseguendo un sogno fragile.
Lui non ha tempo per nessuno, sta viaggiando nel futuro.
\endverse
\beginverse
Lei ripensa ai passi che non ha sta vagando nel passato;
ha negli occhi un che di arido, se potesse cambierebbe tutto di sé.
Ma ormai l'acqua che è passata non fa più girar la ruota.
\endverse
\beginverse
Al centro di questo presente scorre l'eterno nelle cose.
E se in quest'attimo tu ami incontri Dio nella tua vita.
\endverse
\endsong

% Fonte: LibrettoAccordi.pdf (Canti Coro), Q1.
% Nell'elenco: Finale, Ordinario.
\beginsong{Alma misionera}[by={P. E. García Vélez}]
\beginverse
Señor, toma mi vida nueva,
antes de que la espera desgaste años en mí.
Estoy dispuesto a lo que quieras,
no importa lo que sea, Tú llámame a servir.
\endverse
\beginchorus
Llévame donde los hombres necesiten Tus palabras,
necesiten mis ganas de vivir.
Donde falte la esperanza, donde falte la alegría
simplemente por no saber de ti.
\endchorus
\beginverse
Te doy mi corazón sincero,
para gritar sin miedo lo hermoso que es tu amor.
Señor, tengo alma misionera,
condúceme a la tierra que tenga sed de ti.
\endverse
\beginverse
Así, en marcha iré cantando,
por calles predicando tu grandeza, Señor.
Tendré mis manos sin cansancio,
tu historia entre mis labios y fuerza en la oración.
\endverse
\endsong

% Fonte: LibrettoAccordi.pdf (Canti Coro), K5.
% Nell'elenco: senza tipologia (nel libretto accordi tra i canti finali).
\beginsong{E la strada si apre}[by={Gen Arcobaleno}]
\beginverse
Raggio che buca le nubi ed è già cielo aperto,
acqua che scende decisa scavando da sé l'argine per la vita,
la traiettoria di un volo che sull'orizzonte di sera,
tutto di questa natura è una strada per sé.
\endverse
\beginverse
Attimo che segue attimo, un salto nel tempo,
passo di un mondo che tende oramai all'unità, che non è più domani,
usiamo allora queste mani, scaviamo a fondo nel cuore,
solo scegliendo l'amore il mondo vedrà.
\endverse
\beginchorus
Che la strada si apre, passo dopo passo, ora su questa strada noi.
E si spalanca un cielo, un mondo che rinasce, si può vivere nell'unità.
\endchorus
\beginverse
Nave che segue una rotta in mezzo alle onde,
uomo che s'apre la strada in una giungla di idee, seguendo sempre il sole,
quando si sente assetato deve raggiungere l'acqua,
sabbia che nella risacca ritorna al mare.
usiamo allora queste mani, scaviamo a fondo nel cuore,
solo scegliendo l'amore il mondo vedrà.
\endverse
\endsong

% Fonte: LibrettoAccordi.pdf (Canti Coro), K10.
% Nell'elenco: Inno GMG 2000.
\beginsong{Emmanuel}[by={M. Mammoli}]
\beginverse
Dall'orizzonte una grande luce viaggia nella storia
e lungo gli anni ha vinto il buio facendosi memoria,
e illuminando la nostra vita chiaro ci rivela
che non si vive se non si cerca la Verità.
\endverse
\beginverse
Da mille strade arriviamo a Roma sui passi della fede,
sentiamo l'eco della Parola che risuona ancora
da queste mura, da questo cielo per il mondo intero:
è vivo oggi, è l'Uomo Vero, Cristo tra noi.
\endverse
\beginchorus
Siamo qui, sotto la stessa luce, sotto la sua croce, cantando ad una voce:
è l'Emmanuel, Emmanuel, Emmanuel.
È l'Emmanuel, Emmanuel.
\endchorus
\beginverse
Dalla città di chi ha versato il sangue per amore,
ed ha cambiato il vecchio mondo vogliamo ripartire.
Seguendo Cristo, insieme a Pietro, rinasce in noi la fede,
Parola viva che ci rinnova e cresce in noi.
\endverse
\beginverse
Un grande dono che Dio ci ha fatto è Cristo il suo Figlio,
e l'umanità è rinnovata, è in Lui salvata.
È vero uomo, è vero Dio, è il Pane della Vita,
che ad ogni uomo, ai suoi fratelli ridonerà.
\endverse
\beginverse
La morte è uccisa, la vita ha vinto, è Pasqua in tutto il mondo,
un vento soffia in ogni uomo lo Spirito fecondo,
Che porta avanti nella storia la Chiesa sua sposa,
sotto lo sguardo di Maria Comunità.
\endverse
\endsong

% Fonte: LibrettoAccordi.pdf (Canti Coro), J32.
% Nell'elenco: Finale, Ordinario.
\beginsong{Il mio spirito in festa}[by={M. Balduzzi (CD Xaipe)}]
\beginverse
Nessuno mai ha guardato a me,
con gli occhi nuovi del mattino,
nessuno mai ha pensato a me,
come ha pensato il mio Signore,
ha pensato il mio Signore.
\endverse
\beginchorus
L'anima mia esulta,
e il mio spirito in festa canta per te.
\endchorus
\beginverse
Nessuno mai solo resterà,
in ogni angolo del mondo,
nessuno mai ci separerà,
dalla presenza del tuo amore,
la presenza del tuo amore.
\endverse
\endsong

% Fonte: LibrettoAccordi.pdf (Canti Coro), J11.
% Nell'elenco: Finale, Ordinario.
\beginsong{Inno del missionario}
\beginverse
Signor, prendi la mia vita,
prima che l'attesa consumi anni in me.
Io sono aperto a ciò che mi chiedi,
qualunque cosa sia, Tu chiamami a servir.
\endverse
\beginchorus
Portami là dove all'uomo manchino le Tue parole
e la voglia di vivere con Te.
Dove manchi la speranza, dove manchi l'allegria
solamente perché non san di Te.
\endchorus
\beginverse
Ti do il mio cuor sincero,
per dire con coraggio quanto è bello il tuo amor.
Signor, ho un cuore missionario,
conducimi alla terra che si disseti in Te.
\endverse
\beginverse
Così io partirò cantando,
a tutti predicando la Tua immensità.
Sarò Tuo servo senza sosta,
il Tuo predicatore, la forza del Tuo amor.
\endverse
\endsong

% Fonte: LibrettoAccordi.pdf (Canti Coro), K7.
% Nell'elenco: Finale, Ordinario.
\beginsong{Le tue meraviglie}[by={M. Balduzzi, C. Casucci (CD Verbum Panis)}]
\beginchorus
Ora lascia o Signore che io vada in pace perché ho visto le tue meraviglie.
Il tuo popolo in festa per le strade correrà a portare le tue meraviglie.
\endchorus
\beginverse
La tua presenza ha riempito d'amore le nostre vite e le nostre giornate.
In te una sola anima, un solo cuore siamo noi
con te la luce risplende, splende più chiara che mai.
\endverse
\beginverse
La tua presenza ha inondato d'amore le nostre vite e le nostre giornate.
Fra la tua gente resterai, per sempre vivo in mezzo a noi
fino ai confini del tempo, così ci accompagnerai.
\endverse
\endsong

% Fonte: LibrettoAccordi.pdf (Canti Coro), A12. Le strofe "Ti ringraziamo..." sono parlate:
% tenute nel libretto su indicazione dell'utente.
% Nell'elenco: Finale, Ordinario.
\beginsong{Lodate Dio}[by={Gen Rosso}]
\beginchorus
Lodate Dio, cieli immensi ed infiniti.
Lodate Dio, cori eterni d'angeli.
Lodate Dio, santi del suo Regno.
Lodatelo uomini, Dio vi ama.
Lodatelo uomini, Dio è con voi.
\endchorus
\beginverse
Ti ringraziamo, Dio nostro Padre
perché sei Amore. Tu ci hai fatto dono della vita
e ci hai creati per essere figli tuoi.
Ti ringraziamo perché ci fai partecipi
della tua opera creatrice
dandoci un mondo da plasmare
con le nostre mani.
\endverse
\beginverse
Ti ringraziamo, Signore Gesù,
perché per amore nostro sei venuto in questo mondo
per riscattarci da ogni male e riportarci al Padre.
Ti ringraziamo perché sei rimasto
in mezzo a noi per sempre
per far di tutti gli uomini una sola famiglia,
un corpo solo in te.
\endverse
\beginverse
Ti ringraziamo, Spirito d'amore,
perché rinnovi la faccia della terra.
Tu dai luce e conforto ad ogni cuore.
Ti ringraziamo perché con i tuoi santi doni
ci dai la forza di avanzare nel nostro cammino
per giungere uniti alla gioia della tua casa.
\endverse
\endsong

% Fonte: LibrettoAccordi.pdf (Canti Coro), A14.
% Nell'elenco: Finale, Ordinario.
\beginsong{Musica di festa}[by={E. Sands, P. Jotti}]
\beginverse
Cantate al Signore un cantico nuovo, grande è la sua gloria!
Grande la sua forza, grande la sua pace, grande la sua santità!
\endverse
\beginchorus
In tutta la terra, popoli del mondo, gridate la sua fedeltà!
Musica di festa, musica di lode, musica di libertà!
\endchorus
\beginverse
Agli occhi del mondo ha manifestato la sua salvezza!
Per questo si canti, per questo si danzi, per questo si celebri!
\endverse
\beginverse
Con l'arpa ed il corno, con timpani e flauti, con tutta la voce,
canti di dolcezza, canti di salvezza, canti d'immortalità.
\endverse
\beginverse
I fiumi ed i monti battono le mani davanti al Signore.
La sua giustizia giudica la terra, giudica le genti!
\endverse
\beginverse
Al Dio che ci salva, gloria in eterno. Amen! Alleluja!
Gloria a Dio Padre, gloria a Dio Figlio, gloria a Dio Spirito!
\endverse
\endsong

% Fonte: testo fornito dall'utente.
% Nell'elenco: Comunione; l'utente lo indica come Finale.
\beginsong{Nelle mani tue}[by={Gen Verde}]
\beginverse
Sopra un mare di luce,
il vento nelle nostre vele,
sulla rotta di un sogno
nell'immenso blu.
Navighiamo sicuri,
siamo uniti, siamo insieme
e il sole splende sempre più.
\endverse
\beginchorus
E tutto parla di te,
del tuo amore per noi,
e noi crediamo in te, Signore,
sei tu il nostro Dio:
noi siamo nelle mani tue.
\endchorus
\beginverse*
We are, we are, we are,
we are in your hands. \rep{2}
\endverse
\beginverse
Ora il vento è più forte,
si sollevano le onde,
è calato il sole,
c'è solo freddo e buio.
Dove sei, o Signore?
Siamo persi nella notte,
ci puoi salvare solo tu.
\endverse
\beginchorus
E noi guardiamo a te,
al tuo amore per noi,
e noi crediamo in te, Signore,
sei tu il nostro Dio:
noi siamo nelle mani tue.
\endchorus
\beginverse
E quando il mare infuria
e ci prende la paura
e intorno a noi c'è solo oscurità,
tu cammini sulle acque
e ci chiami verso te.
Signore, sei davvero tu? \rep{2}
\endverse
\beginchorus
E tutto parla di te,
del tuo amore per noi,
cammineremo sulle acque
perché tu sei con noi:
noi siamo nelle mani tue.
E noi guardiamo a te,
al tuo amore per noi,
cammineremo sulle acque
perché tu sei con noi:
noi siamo nelle mani tue.
\endchorus
\beginverse*
We are, we are, we are,
we are in your hands. \rep{2}
\endverse
\beginverse*
Cammineremo sulle acque
perché tu sei con noi:
noi siamo nelle mani tue. \rep{2}
\endverse
\endsong

% Fonte: LibrettoAccordi.pdf (Canti Coro), J1.
% Nell'elenco: Finale, Ordinario.
\beginsong{Nozze a Cana}[by={RnS}]
\beginchorus
Danzano con gioia le figlie di Israele,
le nozze sono pronte per il figlio del re.
Alleluja, alleluja per il figlio del re. \rep{2}
\endchorus
\beginverse
A Cana in Galilea nel nome del Signore
gli sposi han giurato amore e fedeltà.
Coi cembali e coi flauti le danze fan corona,
il vino sulla mensa il cuor rallegrerà.
\endverse
\beginverse
Nel mezzo della festa il vino viene meno,
non hanno più la gioia, la danza finirà.
La Madre dice ai servi: «Udite la Parola
che il Figlio mio vi dona, lui vi disseterà».
\endverse
\beginverse
Si arrossano le coppe di vino nuovo colme,
è il dono dello Sposo per nuova fedeltà.
Danziamo allor fratelli: del Re noi siamo figli,
a lui cantiamo lodi per sempre, alleluja!
\endverse
\beginverse*
Danzano con gioia per il figlio del re. \rep{3}
\endverse
\endsong

% Fonte: LibrettoAccordi.pdf (Canti Coro), K8.
% Nell'elenco: Finale, Ordinario (due righe). CD: Cerco il tuo volto (2000).
\beginsong{Resta accanto a me}[by={Gen Verde (CD Cerco il tuo volto)}]
\beginchorus
Ora vado sulla mia strada, con l'amore tuo che mi guida.
O Signore, ovunque io vada resta accanto a me!
Io ti prego, stammi vicino ogni passo del mio cammino.
Ogni notte, ogni mattino, resta accanto a me.
\endchorus
\beginverse
Il tuo sguardo puro sia luce per me
e la tua parola sia voce per me.
Che io trovi il senso del mio andare solo in te,
nel tuo fedele amare il mio perché.
\endverse
\beginverse
Fa' che chi mi guarda non veda che te,
fa' che chi mi ascolta non senta che te.
E chi pensa a me, fa' che nel cuore pensi a te
e trovi quell'amore che hai dato a me.
\endverse
\endsong

% Fonte: LibrettoAccordi.pdf (Canti Coro), K2.
% Nell'elenco: Finale, Ordinario. CD: Una storia che cambia (1984).
\beginsong{Resta qui con noi}[by={Gen Rosso (CD Una storia che cambia)}]
\beginverse
Le ombre si distendono, scende ormai la sera
e s'allontanano dietro i monti
i riflessi di un giorno che non finirà,
di un giorno che ora correrà sempre,
perché sappiamo che una nuova vita
da qui è partita e mai più si fermerà.
\endverse
\beginchorus
Resta qui con noi, il sole scende già,
resta qui con noi, Signore, è sera ormai.
Resta qui con noi, il sole scende già,
se tu sei fra noi la notte non verrà.
\endchorus
\beginverse
S'allarga verso il mare il tuo cerchio d'onda,
che il vento spingerà fino a quando
giungerà ai confini di ogni cuore,
alle porte dell'amore vero.
Come una fiamma che dove passa brucia,
così il tuo amore tutto il mondo invaderà.
\endverse
\beginverse
Davanti a noi l'umanità lotta, soffre e spera
come una terra che nell'arsura
chiede l'acqua da un cielo senza nuvole,
ma che sempre le può dare vita.
Con te saremo sorgente di acqua pura,
con te fra noi il deserto fiorirà.
\endverse
\endsong

% Fonte: LibrettoAccordi.pdf (Canti Coro), K4 (li' senza autore).
% Nell'elenco: Finale (autore Baggio).
\beginsong{Siamo Chiesa del Signore}[by={F. Baggio}]
\beginchorus
Siamo Chiesa del Signore, siamo l'unica famiglia di Dio!
Uno spirito, un solo corpo, siamo tutti comunità.
\endchorus
\beginverse
Nella casa del Padre soli non siamo mai: cittadini del Regno, familiari di Dio,
con gli apostoli edificati sullo stesso Gesù.
\endverse
\beginverse
E la pietra angolare è il Signore Gesù: costruzione ordinata, tempio santo di Dio,
dolce dimora dello Spirito sino all'eternità.
\endverse
\beginverse
Una sola speranza nel Signore Gesù: divenire conformi all'immagine sua,
è chiamata alla santità tutta l'umanità.
\endverse
\endsong

% Fonte: "10 - Finali/vieni e seguimi Sol.pdf" e "Vieni e Seguimi (La).pdf".
% Nell'elenco: senza tipologia/periodo (il file e' tra i Finali).
\beginsong{Vieni e seguimi}[by={Gen Rosso}]
\beginverse
Lascia che il mondo vada per la sua strada.
Lascia che l'uomo ritorni alla sua casa.
Lascia che la gente accumuli la sua fortuna.
Ma tu, tu vieni e seguimi, tu, vieni e seguimi.
\endverse
\beginverse
Lascia che la barca in mare spieghi la vela.
Lascia che trovi affetto chi segue il cuore.
Lascia che dall'albero cadano i frutti maturi.
Ma tu, tu vieni e seguimi, tu, vieni e seguimi.
\endverse
\beginchorus
E sarai luce per gli uomini
e sarai sale della terra
e nel mondo deserto aprirai
una strada nuova. \rep{2}
\endchorus
\beginverse*
E per questa strada va', va',
e non voltarti indietro, va',
e non voltarti indietro.
\endverse
\endsong

% Fonte: LibrettoAccordi.pdf (Canti Coro), A7 (li' tra i canti d'ingresso).
% Nell'elenco: Finale.
\beginsong{Vivere la vita}[by={Gen Verde}]
\beginverse
Vivere la vita con le gioie e i dolori di ogni giorno,
è quello che Dio vuole da te.
Vivere la vita e inabissarti nell'amore è il tuo destino,
è quello che Dio vuole da te.
Fare insieme agli altri la tua strada verso lui,
correre con i fratelli tuoi.
Scoprirai allora il cielo dentro di te,
una scia di luce lascerai.
\endverse
\beginverse
Vivere la vita è l'avventura più stupenda dell'amore,
è quello che Dio vuole da te.
Vivere la vita è generare ogni momento il paradiso:
è quello che Dio vuole da te.
Vivere perché ritorni al mondo l'unità,
perché Dio sta nei fratelli tuoi.
Scoprirai allora il cielo dentro di te,
una scia di luce lascerai. \rep{2}
\endverse
\endsong

% Fonte: LibrettoAccordi.pdf (Canti Coro), J24 (li' tra i canti di comunione).
% Nell'elenco: Finale, Ordinario. CD: In cerca d'autore (1977).
\beginsong{Vocazione}[by={P. A. Sequeri (CD In cerca d'autore)}]
\beginverse
Era un giorno come tanti altri, e quel giorno lui passò.
Era un giorno come tutti gli altri e, passando, mi chiamò.
Come lo sapesse che il mio nome era proprio quello,
come mai vedesse proprio me nella sua vita non lo so.
Era un giorno come tanti altri, e quel giorno mi chiamò.
\endverse
\beginchorus
Tu, Dio, che conosci il nome mio, fa' che ascoltando la tua voce
io ricordi dove porta la mia strada, nella vita all'incontro con te. \rep{2}
\endchorus
\beginverse
Era l'alba triste e senza vita, e qualcuno mi chiamò:
era un uomo come tanti altri, ma la voce, quella no.
Quante volte un uomo con il nome giusto mi ha chiamato,
una volta sola l'ho sentito pronunciare con amore.
Era un uomo come nessun altro, e quel giorno mi chiamò.
\endverse
\endsong

\end{songs}

% !TEX root = ../libretto.tex
% Sezione: Canti mariani
\momento{Canti mariani}
\begin{songs}{titoli}

% Fonte: LibrettoAccordi.pdf (Canti Coro), M2.
% Nell'elenco: Canto Mariano, Ordinario.
\beginsong{Ave, Maria}[by={M. Balduzzi, C. Casucci (CD Verbum Panis)}]
\beginchorus
Ave Maria! Ave!
Ave Maria! Ave!
\endchorus
\beginverse
Donna dell'attesa e madre di speranza, ora pro nobis.
Donna del sorriso e madre del silenzio, ora pro nobis.
Donna di frontiera e madre dell'ardore ora pro nobis.
Donna del riposo e madre del sentiero, ora pro nobis.
\endverse
\beginverse
Donna del deserto e madre del respiro, ora pro nobis.
Donna della sera e madre del ricordo, ora pro nobis.
Donna del presente e madre del ritorno, ora pro nobis.
Donna della terra e madre dell'amore, ora pro nobis.
\endverse
\endsong

% Fonte: LibrettoAccordi.pdf (Canti Coro), M12.
% Nell'elenco: Canto Mariano, Ordinario.
\beginsong{Ave Maria}[by={Gen Rosso}]
\beginverse
Ave Maria, piena di grazia,
Il Signore è con te, benedetta fra le donne.
Tu sei la madre di ogni uomo,
tu sei la luce che ci guida nella nostra notte.
Regina del cielo, Madre di Dio,
dolce respiro dell'anima mia.
\endverse
\beginchorus
Ave Maria!
Ave Maria! \rep{2}
\endchorus
\endsong

% Fonte: "11 - Canti Mariani/dell_aurora tu sorgi piu bella ORIGINALE.pdf" (foglio accordi, n. 420).
% Nell'elenco: senza tipologia (Compagnia Aquero).
\beginsong{Dell'aurora tu sorgi più bella}[by={Compagnia Aquero}]
\beginverse
Dell'aurora tu sorgi più bella, coi tuoi raggi a far lieta la terra
e fra gli astri che il cielo rinserra non v'è stella come te.
Gli occhi tuoi son più fondi del mare,
la tua fronte ha il profumo del giglio,
il tuo viso ricorda tuo Figlio, sui tuoi passi nascon fiori.
\endverse
\beginchorus
Bella tu sei qual sole, bianca più della luna
e le stelle più belle non son belle come te. \rep{2}
\endchorus
\beginverse
Ti coronano tutte le stelle, al tuo canto risponderà il vento,
della luna si curva l'argento, si rivolge verso te;
quando tutto dintorno è rovina
e la voce del pianto non tace,
il tuo sguardo riporta la pace, la concordia in fondo ai cuori.
\endverse
\endsong

% Fonte: LibrettoAccordi.pdf (Canti Coro), M6.
% Nell'elenco: Canto Mariano, Ordinario.
\beginsong{Giovane donna}[by={L. Scaglianti, L. Bancolini}]
\beginverse
Giovane donna, attesa dell'umanità,
un desiderio d'amore e pura libertà.
Il Dio lontano è qui vicino a te,
voce e silenzio, annuncio di novità.
\endverse
\beginchorus
Ave Maria! \rep{2}
\endchorus
\beginverse
Dio t'ha prescelta qual Madre piena di bellezza,
ed il suo amore ti avvolgerà con la sua ombra.
Grembo di Dio, venuto sulla terra,
tu sarai madre di un uomo nuovo.
\endverse
\beginverse
Ecco l'ancella che vive della sua parola,
libero il cuore perché l'amore trovi casa.
Ora l'attesa è densa di preghiera
e l'uomo nuovo è qui in mezzo a noi.
\endverse
\endsong

% Fonte: LibrettoAccordi.pdf (Canti Coro), M11.
% Nell'elenco: Canto Mariano, Avvento; l'utente lo vuole tra i Mariani.
\beginsong{L'anima mia magnifica}[by={P. Dargenio, M. Zambuto}]
\beginchorus
L'anima mia magnifica, l'anima mia magnifica.
L'anima mia magnifica, magnifica il Signore. \rep{2}
\endchorus
\beginverse
Perché si è chinato quaggiù e ha guardato alla mia umiltà,
trasformando il mio lamento in una danza.
E in eterno, da ora in poi, canteranno le generazioni
chiameranno il mio nome e per sempre beata sarò.
\endverse
\beginverse
Grandi cose ha compiuto per me il Signore, l'Onnipotente,
che dispiega la Sua pietà su chi lo teme.
La potenza della Sua Maestà ha disperso i superbi di cuore,
rovesciando i potenti e innalzando gli umili.
\endverse
\beginverse
Perché con la Sua carità ha sfamato chi non ha più pane,
rimandando indietro i ricchi a mani vuote.
Ha soccorso il popolo suo, come aveva promesso ad Abramo,
ricordando il Suo Amore fedele per l'umanità.
\endverse
\endsong

% Fonte: LibrettoAccordi.pdf (Canti Coro), M4.
% Nell'elenco: Canto Mariano, Ordinario.
\beginsong{La mia anima canta}[by={Gen Verde}]
\beginchorus
La mia anima canta la grandezza del Signore,
il mio spirito esulta nel mio salvatore.
Nella mia povertà l'Infinito mi ha guardata,
in eterno ogni creatura mi chiamerà beata.
\endchorus
\beginverse
La mia gioia è nel Signore che ha compiuto grandi cose in me,
la mia lode al Dio fedele che ha soccorso il suo popolo
e non ha dimenticato le sue promesse d'amore.
\endverse
\beginverse
Ha disperso i superbi nei pensieri inconfessabili,
ha deposto i potenti, ha risollevato gli umili,
ha saziato gli affamati e aperto ai ricchi le mani.
\endverse
\endsong

% Fonte: LibrettoAccordi.pdf (Canti Coro), M8.
% Nell'elenco: Canto Mariano, Avvento; l'utente lo vuole tra i Mariani.
\beginsong{La mia anima esulta}[by={D. Ricci}]
\beginverse
L'anima mia magnifica il Signore
e il mio spirito esulta in Dio, mio salvatore,
perché ha guardato l'umile sua ancella,
d'ora in poi l'umanità mi chiamerà beata.
Grandi cose, grandi cose ha fatto in me
L'onnipotente, e santo è il nome suo.
La misericordia sua si estenderà
su tutti quelli che lo temono.
\endverse
\beginchorus
L'anima esulta in Dio. La mia anima esulta in Dio.
L'anima esulta in Dio. La mia anima esulta in Dio.
\endchorus
\beginverse
L'anima mia magnifica il Signore
e il mio spirito esulta in Dio, mio salvatore,
perché ha guardato l'umile sua ancella,
d'ora in poi l'umanità mi chiamerà beata.
Ha disteso onnipotente il braccio suo
disperdendo le ambizioni dei superbi, lui.
Giù dai troni ha rovesciato i grandi e
ha posto in alto invece gli umili.
\endverse
\beginverse
L'anima mia magnifica il Signore
e il mio spirito esulta in Dio, mio salvatore,
perché ha guardato l'umile sua ancella,
d'ora in poi l'umanità mi chiamerà beata.
Ha versato beni sopra chi non ha
e i ricchi li ha mandati a mani vuote via.
Ha soccorso noi, suo popolo, perché
lui è fedele nell'amore suo.
\endverse
\endsong

% Fonte: LibrettoAccordi.pdf (Canti Coro), M3.
% Nell'elenco: Canto Mariano, Ordinario.
\beginsong{Madre dolcissima}[by={Gen Rosso}]
\beginverse
Se io fossi poeta, scriverei di te le cose che nessuno ha detto mai.
Dell'arte avessi il dono, io rapirei la luce per poter disegnare gli occhi tuoi.
Se grandi avessi le ali, le scioglierei nel volo per arrivare in alto fino a te.
\endverse
\beginchorus
Madre Dolcissima,
tesoro dei più poveri, coraggio di chi ancora spera. Tu,
Madre Dolcissima, consola queste lacrime, ascolta chi parole più non ha.
\endchorus
\beginverse
Riposo dei più deboli, silenzio di chi dà, fontana per la nostra siccità,
splendore della sera, stella del mattino, neve bianca sopra le città,
regina della pace, vetta dell'amore, radice della nostra vita, tu
\endverse
\beginverse
Dei naufraghi tu l'ancora, del canto poesia, dei nostri sogni unica realtà,
riflesso d'infinito, chiarore delle stelle, immagine del cosmo che sarà,
custode della terra, sorgente fra le rocce, segreto schiuso da ogni fiore, tu,
\endverse
\endsong

% Fonte: LibrettoAccordi.pdf (Canti Coro), M5.
% Nell'elenco: Canto Mariano, Ordinario.
\beginsong{Maria}[by={RnS}]
\beginverse
Ora che sei qui, la tua dolcezza sento dentro me.
Come l'aurora sorgi, risplendi fra noi.
Ora che sei qui, tu figlia e madre, fra le braccia tue
diffondi la tua luce celeste di pace e di amor.
\endverse
\beginchorus
Tu Maria, tu sei dimora per Gesù figlio tuo.
Regina che dal cielo speranza e amore ci dai.
\endchorus
\beginverse
Ora che sei qui, tu dolce madre sono figlio tuo
trasforma il mio pianto nel canto d'amore per te.
\endverse
\beginverse
Ora che sei qui, insegnami ad amare come sai,
radiosa tra le stelle, preghiera ora sei.
\endverse
\endsong

% Fonte: LibrettoAccordi.pdf (Canti Coro), M10.
% Nell'elenco: Canto Mariano, Ordinario.
\beginsong{Mi alma canta}[by={Gen Verde}]
\beginchorus
Mi alma canta, canta la grandeza del Señor,
mi espíritu se estremece de gozo en Dios,
mi Salvador.
\endchorus
\beginverse
Porque miró con bondad
la pequeñez de su servidora.
Porque miró con bondad
la pequeñez de su servidora.
En adelante todas las gentes
me llamarán feliz!
Me llamarán feliz!
Me llamarán feliz!
\endverse
\beginverse
Derribó del trono a los poderosos,
y elevó a los humildes,
colmó de bienes a los hambrientos
y despidió a los ricos con las manos vacías.
Mi alma canta la grandeza del Señor
y mi espíritu se estremece
de gozo en Dios, mi Salvador.
\endverse
\endsong

% Fonte: LibrettoAccordi.pdf (Canti Coro), M9.
% Nell'elenco: Canto Mariano, Ordinario. CD: Giorno senza tramonto.
\beginsong{Regina della benevolenza}[by={M. Balduzzi (CD Giorno senza tramonto)}]
\beginverse
Stella del mare, signora del cielo,
porto sicuro dell'anima nelle tempeste.
Casa dei popoli, donna d'amore,
dolce regina che dona il re dell'universo.
\endverse
\beginchorus
Madre di Dio, prega per noi,
prega, prega, prega per noi.
\endchorus
\beginverse
Mano che veglia sui cuori nella notte buia,
sguardo di benevolenza sopra la città,
aurora chiara nel giorno nuovo che ci aspetta,
mistica rosa fiorita sulle ingiurie della terra.
\endverse
\endsong

% Fonte: LibrettoAccordi.pdf (Canti Coro), M1. C'e' anche "11 - Canti Mariani/RIMANI CON ME.pdf".
% Nell'elenco: Canto Mariano, Ordinario.
\beginsong{Rimani con me}[by={Gen Rosso}]
\beginverse
Maria, tu sei la vita per me,
sei la speranza, la gioia, l'amore: tutto sei.
Maria tu sai quello che vuoi,
sai con che forza d'amore in cielo mi porterai.
\endverse
\beginchorus
Maria, ti do il mio cuore per sempre se vuoi,
tu dammi l'amore che non passa mai.
Rimani con me e andiamo nel mondo insieme.
La tua presenza sarà goccia di paradiso per l'umanità.
\endchorus
\beginverse
Maria, con te sempre vivrò,
in ogni momento giocando, cantando, ti amerò.
Seguendo i tuoi passi in te io avrò
la luce che illumina i giorni e le notti dell'anima.
\endverse
\endsong

% Fonte: LibrettoAccordi.pdf (Canti Coro), M7.
% Nell'elenco: Canto Mariano, Ordinario.
\beginsong{Santa Maria del cammino}[by={R. D'Andrea, J. A. Espinosa}]
\beginverse
Mentre trascorre la vita
solo tu non sei mai;
Santa Maria del cammino
sempre sarà con te.
\endverse
\beginchorus
Vieni, o Madre, in mezzo a noi, vieni Maria quaggiù.
Cammineremo insieme a te
verso la libertà.
\endchorus
\beginverse
Quando qualcuno ti dice:
«Nulla mai cambierà»,
lotta per un mondo nuovo,
lotta per la verità!
\endverse
\beginverse
Lungo la strada la gente
chiusa in se stessa va;
offri per primo la mano
a chi è vicino a te.
\endverse
\beginverse
Quando ti senti ormai stanco
e sembra inutile andar,
tu vai tracciando un cammino:
un altro ti seguirà.
\endverse
\endsong

\end{songs}

% !TEX root = ../libretto.tex
% Sezione: Adorazione eucaristica
\momento{Adorazione eucaristica}
\begin{songs}{titoli}

% Fonte: testo fornito dall'utente.
% Nell'elenco: "10000 Reasons", Redman, Adorazione, Ordinario.
\beginsong{10000 Reasons (Bless the Lord)}[by={M. Redman}]
\beginchorus
Bless the Lord, O my soul, O my soul,
worship his holy name.
Sing like never before, O my soul,
I'll worship your holy name.
\endchorus
\beginverse
The sun comes up, it's a new day dawning,
it's time to sing your song again.
Whatever may pass and whatever lies before me,
let me be singing when the evening comes.
\endverse
\beginverse
You're rich in love and you're slow to anger,
your name is great and your heart is kind.
For all your goodness I will keep on singing,
ten thousand reasons for my heart to find.
\endverse
\beginverse
And on that day when my strength is failing,
the end draws near and my time has come,
still my soul will sing your praise unending,
ten thousand years and then forevermore.
\endverse
\beginverse*
Sing like never before, O my soul,
I'll worship your holy name.
O I'll worship your holy name,
I will worship your holy name,
holy name.
\endverse
\endsong

% Fonte: testo fornito dall'utente. Non presente nell'elenco. Autore non indicato.
% Titolo dal primo verso. Ritornello finale ripetuto.
\beginsong{All'aurora}
\beginverse
All'aurora il mio spirito ti cerca
come terra arida senz'acqua.
O Signore, vieni in mio aiuto,
la tua mano poni su di me.
\endverse
\beginverse*
Vengo a te, Signore della vita,
la mia anima anela solo a te,
i miei occhi cercano il tuo volto,
mio Dio, mio Re.
\endverse
\beginchorus
Donami il tuo amore, o Signor,
con il tuo potere toccami,
manda la tua grazia come un fiume su di me
e liberami, consolami, guariscimi.
\endchorus
\beginverse
Ho bisogno di te, o mio Signore,
la tua pace scenda nel mio cuore.
Come un cieco cerco la tua luce,
la tua mano guidi i passi miei.
\endverse
\beginverse*
Vengo a te, Signore della vita,
la mia anima anela solo a te,
i miei occhi cercano il tuo volto,
mio Dio, mio Re.
\endverse
\beginchorus
Donami il tuo amore, o Signor,
con il tuo potere toccami,
manda la tua grazia come un fiume su di me
e liberami, consolami, guariscimi. \rep{2}
\endchorus
\endsong

% Fonte: LibrettoAccordi.pdf (Canti Coro), J30.
% Nell'elenco: Adorazione, Ordinario.
\beginsong{Ascolta il mio cuore}[by={RnS}]
\beginverse
Ascolta il mio cuore, rispondi al mio grido,
t'invoco nella mia povertà.
Custodisci i miei pensieri, con il tuo immenso amore,
rinnova in me la tua fedeltà.
\endverse
\beginchorus
Per te canterò, per te io vivrò,
perché sei, oltre i confini,
il Signore che vita mi dà.
Per te danzerò, con te cambierò
il mio pianto nel canto nuovo
perché vivi in me, sei risorto e con me resterai.
\endchorus
\beginverse
Ritorno a te mio Dio, accolgo il tuo perdono,
la tua misericordia è per me.
Dal profondo della terra mi hai chiamato a nuova vita,
per sempre grazie ti renderò.
\endverse
\endsong

% Fonte: testo fornito dall'utente. Non presente nell'elenco.
\beginsong{Come in cielo così in terra}[by={D. Vezzani}]
\beginverse
Io ti chiamo piangendo,
soffocando d'amore per te.
Tu che sei mio figlio, dimmi
dove sei: ti voglio con me.
Ho fatto tutto per te,
voglio solo la tua volontà per me.
Le tue miserie, le tue debolezze,
i mali, le colpe, la vita passata:
tu donami tutto,
io brucerò tutto
col mio immenso amore per te.
\endverse
\beginchorus
Vengo io a cercare te,
a portarti il grande dono della mia volontà.
Se mi accetterai, sarò con te
con la pace, con la gioia e con la piena felicità.
\endchorus
\beginverse
Io ti sto preparando,
ho donato mia madre per te:
lei ti ama, tu sei suo figlio,
ti affido a lei.
Lezioni di cielo
per darti amore e luce
e formare la mia santità in te.
Le mie ricchezze,
la gioia più pura,
la pace, il coraggio e forza divina,
la felicità e la mia stessa vita:
ti dono la mia volontà.
\endverse
\beginchorus
Vengo io a cercare te,
a portarti il grande dono della mia volontà.
Se mi accetterai, sarò con te
con la pace, con la gioia e con la piena felicità.
\endchorus
\beginverse*
Come in cielo così in terra
io voglio farti vivere:
è ora di vivere. \rep{2}
\endverse
\beginchorus
Vengo io a cercare te,
a portarti il grande dono della mia volontà.
Se mi accetterai, sarò con te
con la pace, con la gioia e con la piena felicità.
\endchorus
\beginverse*
Se mi accetterai, sarò con te
con la pace, con la gioia e con la piena felicità.
Con la pace, con la gioia e con la piena felicità.
\endverse
\endsong

% Fonte: LibrettoAccordi.pdf (Canti Coro), J43.
% Nell'elenco: Adorazione.
\beginsong{Come tu mi vuoi}[by={RnS}]
\beginverse
Eccomi Signor, vengo a te mio Re
che si compia in me la tua volontà
Eccomi Signor, vengo a te mio Dio
plasma il cuore mio e di Te vivrò
Se tu lo vuoi Signore manda me
e il tuo nome annuncerò.
\endverse
\beginchorus
Come tu mi vuoi io sarò,
dove tu mi vuoi io andrò.
Questa vita io voglio donarla a te
per dar gloria al tuo nome mio Re.
Come tu mi vuoi io sarò,
dove tu mi vuoi io andrò.
Se mi guida il tuo amore paura non ho,
per sempre io sarò
come tu mi vuoi.
come tu mi vuoi, come tu mi vuoi,
come tu mi vuoi, come tu mi vuoi,
come tu mi vuoi, come tu mi vuoi,
come tu mi vuoi.
\endchorus
\beginverse
Eccomi Signor, vengo a te mio Re
che si compia in me la tua volontà
Eccomi Signor, vengo a te mio Dio
plasma il cuore mio e di Te vivrò
Tra le tue mani mai più vacillerò
e strumento tuo sarò.
\endverse
\endsong

% Fonte: LibrettoAccordi.pdf (Canti Coro), P6.
% Nell'elenco: Adorazione.
\beginsong{Come un prodigio}[by={D. Vezzani}]
\beginverse
Signore, Tu mi scruti e conosci,
sai quando seggo e quando mi alzo.
Riesci a vedere i miei pensieri,
sai quando io cammino e quando riposo.
Ti sono note tutte le mie vie:
la mia parola non è ancora sulla lingua
e Tu Signore già la conosci tutta
\endverse
\beginchorus
Sei Tu che mi hai creato
e mi hai tessuto nel seno di mia madre.
Tu mi hai fatto come un prodigio:
le Tue opere sono stupende e per questo Ti lodo.
\endchorus
\beginverse
Di fronte e alle spalle Tu mi circondi,
poni su me la Tua mano.
La Tua saggezza, stupenda per me,
è troppo alta e io non la comprendo.
Che sia in cielo o agli inferi, ci sei:
non si può mai fuggire dalla Tua presenza.
Ovunque la Tua mano guiderà la mia
\endverse
\beginverse
E nel segreto Tu mi hai formato,
mi hai intessuto dalla terra.
Neanche le ossa ti eran nascoste,
ancora informe mi hanno visto i tuoi occhi
I miei giorni erano fissati
quando ancora non ne esisteva uno,
e tutto quanto era scritto nel tuo libro.
\endverse
\endsong

% Fonte: testo (senza accordi) del foglio accordi fornito dall'utente e da lui indicato come testo giusto.
% Non presente nell'elenco.
\beginsong{En mi Getsemaní}[by={E. Meana}]
\beginverse
Para que mi amor no sea un sentimiento
tan sólo de deslumbramiento pasajero;
para no gastar mis palabras más mías
ni vaciar de contenido mi «te quiero».
Quiero hundir más hondo mi raíz en ti
y cimentar en solidez éste, mi afecto;
pues mi corazón, que es inquieto y es frágil,
sólo acierta, si se abraza a tu proyecto.
\endverse
\beginchorus
Más allá de mis miedos,
más allá de mi inseguridad,
quiero darte mi respuesta:
aquí estoy, para hacer tu voluntad,
para que mi amor sea decirte sí,
hasta el final.
\endchorus
\beginverse
Duermen su sopor y temen en el huerto:
ni sus amigos acompañan al Maestro.
Si es hora de cruz, es de fidelidades;
pero el mundo nunca quiere aceptar esto.
Dame comprender, Señor, tu amor tan puro:
amor que persevera en cruz, amor perfecto.
Dame serte fiel cuando todo es oscuro,
para que mi amor sea más que un sentimiento.
\endverse
\beginchorus
Más allá de mis miedos,
más allá de mi inseguridad,
quiero darte mi respuesta:
aquí estoy, para hacer tu voluntad,
para que mi amor sea decirte sí,
hasta el final.
\endchorus
\beginverse
No es en las palabras ni es en las promesas
donde la historia tiene su motor secreto:
sólo es el amor en la cruz madurado,
el amor que mueve todo el universo.
\endverse
\beginchorus
Más allá de mis miedos,
más allá de mi inseguridad,
quiero darte mi respuesta:
aquí estoy, para hacer tu voluntad,
para que mi amor sea decirte sí,
hasta el final.
\endchorus
\beginverse
Pongo mi pequeña vida hoy en tus manos,
por sobre mis seguridades y mis miedos;
y, para elegir tu querer y no el mío,
hazme, en mi Getsemaní, fiel y despierto.
\endverse
\beginchorus
Más allá de mis miedos,
más allá de mi inseguridad,
quiero darte mi respuesta:
aquí estoy, para hacer tu voluntad,
para que mi amor sea decirte sí,
hasta el final.
\endchorus
\endsong

% Fonte: LibrettoAccordi.pdf (Canti Coro), O19.
% Nell'elenco: Adorazione.
\beginsong{Fiumi di misericordia}[by={RnS}]
\beginverse
Fiumi di misericordia dalla croce lavano la colpa mia.
La tua salvezza, mio Re, luce che rinnova,
e dona alla mia vita la tua eternità.
La tua misericordia, Signor, è il sole che mi illumina, Gesù.
\endverse
\beginchorus
Sei amore che perdona, Parola che dà vita,
ritrovo la mia gioia in te, Gesù.
La voce tua mi attira, la grazia tua mi libera,
in te ritorno a vivere.
\endchorus
\beginverse
Sei tu che visiti il mio cuore, metti in me la pace del tuo Spirito.
È la tua grazia, mio Re, balsamo d'amore
sopra le ferite della vita mia.
La tua misericordia, Signor, è il sole che mi illumina, Gesù.
\endverse
\beginverse
Perdoni il mio peccato, grande è la tua salvezza,
Dio di misericordia e d'infinita tenerezza. \rep{2}
\endverse
\endsong

% Fonte: LibrettoAccordi.pdf (Canti Coro), O16.
% Nell'elenco: Adorazione.
\beginsong{Golgota}[by={Kantiere Kairòs}]
\beginverse
Misteriosamente tutto accade qui:
si pieghino le ginocchia, in croce è messo il Re. In croce è messo il Re.
\endverse
\beginverse
L'altare diventa il Golgota, il pane è la tua carne viva,
l'altare si apre all'eternità,
il vino unito all'acqua è il sangue tuo, Gesù, che ti lega a noi.
\endverse
\beginverse
A Te, Dio Padre, si offre il Figlio tuo,
che muore in croce in nome dell'amore e della carità.
A Te, Dio Padre, si offre il Figlio tuo,
che muore in croce al posto di ogni uomo,
al posto che toccava a me.
\endverse
\endsong

% Fonte: LibrettoAccordi.pdf (Canti Coro), P7.
% Nell'elenco: Adorazione.
\beginsong{I.N.R.I. Io non ritorno indietro}[by={D. Vezzani}]
\beginverse
Ha senso solo così, perdonando proprio tutto.
Ha senso solo così, amando fino in fondo.
Ha senso solo così, anche quando hai tutti contro.
Ha senso solo così, e ora io lo difendo.
\endverse
\beginchorus
Io non ritorno indietro. Io non ritorno indietro.
Da un Amore gigante così io non posso tornare indietro.
Io non ritorno indietro. Io no che non ritorno indietro.
Da un Amore gigante così io resto appesa e non scendo.
Io non scendo.
\endchorus
\beginverse
Ha senso solo così, rischiando proprio tutto.
Ha senso solo così, continuando fino in fondo.
Ha senso solo così, soprattutto quando hai il buio addosso.
Ha senso solo così, e ora io lo difendo.
\endverse
\beginverse
È questa la natura del vero Amore
l'operare da Dio, la fermezza e il non ritirarsi mai.
È questa la natura del vero Amore
l'operare da Dio, la fermezza e il non ritirarsi mai,
a costo di qualunque sacrificio io non ritorno indietro.
\endverse
\endsong

% Fonte: testo fornito dall'utente; autore RnS indicato dall'utente. Non presente nell'elenco.
% Il ritornello si canta due volte: la prima senza "perché tu sei il mio re".
\beginsong{Io scelgo te}[by={RnS}]
\beginverse
Io scelgo te
come Signore della mia vita, scelgo te,
perché il mondo amore non mi dà.
Io voglio amare te,
guardare il tuo volto e vivere per te.
\endverse
\beginchorus
La mia voce sale a te, ascoltami, Signor,
per sempre in te vivrò, io ti benedirò,
accogli la mia lode.
La mia voce sale a te, ascoltami, Signor,
per sempre in te vivrò, io ti benedirò,
perché tu sei il mio re, accogli la mia lode.
\endchorus
\beginverse
Io cerco te,
tu sei il Signore della mia vita, cerco te,
perché il mondo pace non mi dà.
Io voglio amare te,
guardare il tuo volto e vivere per te.
\endverse
\beginchorus
La mia voce sale a te, ascoltami, Signor,
per sempre in te vivrò, io ti benedirò,
accogli la mia lode.
La mia voce sale a te, ascoltami, Signor,
per sempre in te vivrò, io ti benedirò,
perché tu sei il mio re, accogli la mia lode, mio Signor.
La mia voce sale a te, ascoltami, Signor,
per sempre in te vivrò, io ti benedirò,
perché tu sei il mio re, accogli la mia lode.
\endchorus
\endsong

% Fonte: LibrettoAccordi.pdf (Canti Coro), A6.
% Nell'elenco: Adorazione.
\beginsong{Io ti seguirò}[by={RnS}]
\beginverse
Mostrami la via per seguire Te, apri i miei occhi, Gesù;
donami la forza per camminare sulla via che hai tracciato per me.
\endverse
\beginchorus
La Tua croce, o Dio, amerò, e con Te nel mondo la porterò.
O Signore, mia vera libertà, se con me sarai io ti seguirò.
\endchorus
\beginverse
Mostrami la via per raggiungere Te, venga il Tuo Spirito in me;
donami la grazia per rimanere sulla via che mi porta a Te.
\endverse
\endsong

% Fonte: "18 - Frisian/lodi all altissimo.pdf" (partitura per coro). Testo dalla voce principale;
% divisione in 5 strofe confermata dall'utente.
% Nell'elenco: Ingresso, Ordinario (riga 136; riga 131 senza tipologia). CD: Tu sei bellezza (1996).
% Spostato in Adorazione su indicazione dell'utente (CD di preghiera, non per la Messa).
\beginsong{Lodi all'Altissimo}[by={M. Frisina (CD Tu sei bellezza)}]
\beginverse*
Tu sei Santo, Signore Dio,
tu sei forte, tu sei grande,
tu sei l'Altissimo, l'Onnipotente,
tu Padre Santo, Re del cielo.
\endverse
\beginverse*
Tu sei trino, uno Signore,
tu sei il bene, tutto il bene,
tu sei l'Amore, tu sei il vero,
tu sei umiltà, tu sei sapienza.
\endverse
\beginverse*
Tu sei bellezza, tu sei la pace,
la sicurezza, il gaudio, la letizia,
tu sei speranza, tu sei giustizia,
tu temperanza e ogni ricchezza.
\endverse
\beginverse*
Tu sei il Custode, tu sei mitezza,
tu sei rifugio, tu sei fortezza,
tu carità, fede, tu speranza,
tu sei tutta la nostra dolcezza.
\endverse
\beginverse*
Tu sei la vita, eterno gaudio,
Signore grande, Dio ammirabile,
Onnipotente, o Creatore,
o Salvatore di misericordia.
\endverse
\endsong

% Fonte: testo fornito dall'utente; autori dallo spartito "10 - Finali/luce-gentile-spartito-tronconi-fiaschi-newman.pdf"
% (musica M. Fiaschi, testo J. H. Newman; il nome del file cita anche Tronconi). Non presente nell'elenco.
\beginsong{Luce gentile}[by={M. Fiaschi; testo da J. H. Newman}]
\beginverse
Sostieni il mio cammino nella notte,
non chiedo di vedere già la meta.
Tu mostrami soltanto il primo passo
e oggi questo basterà.
Da solo decidevo la mia strada,
l'orgoglio mi guidava nelle scelte.
Dimentica quegli anni in cui non ti pregai:
tu ora sei la guida mia.
\endverse
\beginchorus
Guidami tu, luce gentile,
guidami tu in questo buio che mi stringe.
Scura è la notte e la mia casa è lontana,
ma so che con te io la troverò.
Guidami tu, luce gentile,
guidami tu fino alla luce del mattino.
Scura è la notte e la mia casa è lontana,
ma so che con te io la troverò.
\endchorus
\beginverse
Mi ha sempre sostenuto la tua forza,
per monti e fiumi ancora lei mi porterà.
Sarà con me finché la notte passerà,
mi guiderà sicuro a te.
\endverse
\beginverse*
Guidami tu.
\endverse
\endsong

% Fonte: testo fornito dall'utente (tolti i cori tra parentesi). Non presente nell'elenco.
% Il ritornello si canta più volte; qui scritto due volte, con l'ultimo verso ripetuto.
\beginsong{Mi rialzerai}[by={G. Ammirabile}]
\beginverse*
Se dentro me ho perso la speranza
e sento che certezze più non ho,
non temerò, ma aspetterò in silenzio
perché io so che sei vicino a me.
\endverse
\beginchorus
Mi rialzerai se non avrò più forze,
mi rialzerai, con te ce la farò.
Sarai con me nel buio della notte,
mi rialzerai e in alto volerò.
\endchorus
\beginchorus
Mi rialzerai se non avrò più forze,
mi rialzerai, con te ce la farò.
Sarai con me nel buio della notte,
mi rialzerai e in alto volerò. \rep{2}
\endchorus
\endsong

% Fonte: testo fornito dall'utente ("persona" corretto in "perdona", "cogli" in "accogli").
% Nell'elenco: "Potente sei mio Signor", RnS, Adorazione. CD: Misericordia sei.
\beginsong{Potente sei, mio Signor}[by={RnS (CD Misericordia sei)}]
\beginverse
Dio di misericordia,
amore che non muore,
riempi il mio cuor.
Dio che perdona sempre,
Padre di tenerezza,
speranza eterna.
\endverse
\beginchorus
Mio Dio, muovi le montagne,
potente sei, mio Signor.
So che mi salverai per sempre,
Dio della salvezza:
la morte hai vinto per noi,
tu hai vinto per noi.
\endchorus
\beginverse
Mi accogli come sono,
non guardi ai miei errori,
riempi il mio cuor.
Ti dono la mia vita,
mi arrendo a te, Signore,
speranza eterna.
\endverse
\beginchorus
Mio Dio, muovi le montagne,
potente sei, mio Signor.
So che mi salverai per sempre,
Dio della salvezza:
la morte hai vinto per noi,
tu hai vinto per noi.
\endchorus
\beginverse*
Gloria, la tua luce sempre splenderà,
gloria, noi cantiamo a Gesù il Signor.
Gloria, la tua luce sempre splenderà,
gloria, noi cantiamo a Gesù il Signor.
\endverse
\beginchorus
Mio Dio, muovi le montagne,
potente sei, mio Signor.
So che mi salverai per sempre,
Dio della salvezza:
la morte hai vinto per noi,
tu hai vinto per noi.
\endchorus
\beginverse*
Gloria, la tua luce sempre splenderà,
gloria, noi cantiamo a Gesù il Signor.
Gloria, la tua luce sempre splenderà,
gloria, noi cantiamo a Gesù il Signor.
\endverse
\endsong

% Fonte: LibrettoAccordi.pdf (Canti Coro), O15 (anche Q4). Messo in Adorazione su indicazione dell'utente.
\beginsong{Potrò mai dimenticare}[by={S. Zampa}]
\beginverse
L'ho cercata la felicità
troppo tempo nelle vanità, lontano.
Tante volte nel passato
il mio cuore si è slanciato verso te.
Potrò mai dimenticare quelle lacrime amare
che solcavano il mio viso, così all'improvviso,
e una danza di gioia e dolore dentro me,
ma ora tu sei l'unico mio bene.
\endverse
\beginchorus
Tu, sulla croce tu, hai colmato ogni vuoto.
Tu, sulla croce tu, una gioia che non può finire.
Lì, sulla croce tu, in un vortice d'amore,
doni al mondo senza vita il segreto di ogni cosa.
\endchorus
\beginverse
Il tuo sguardo su di me,
e l'Amore si riaccende nel cuore.
Io non mi appartengo più,
voglio darti ciò che sono, ora.
Per te, che mi hai pensato, voglio vivere per te,
voglio amare solo te in chi mi passa accanto.
Nei miei occhi brilla la tua luce e scopro che
l'umanità cerca ma non sa.
\endverse
\beginchorus
Tu, sulla croce tu, hai colmato ogni vuoto.
Tu, sulla croce tu, una gioia che non può finire.
Lì, sulla croce tu, in un vortice d'amore,
doni al mondo senza vita il segreto di ogni cosa.
\endchorus
\endsong

% Fonte: testo fornito dall'utente. Non presente nell'elenco.
\beginsong{Preghiera (Progetto 1)}[by={Progetto 1}]
\beginverse
Disegni il tempo con le tue grandi mani
e lo racchiudi dentro te.
I tuoi respiri sembrano lontani
eppure avvolgono anche me.
E cercherò di indovinare
le cose che mi chiederai,
quando ti avrò davanti al cuore, tu lo sai:
non ti potrei mentire.
\endverse
\beginverse
Sorvoli il mondo con le tue grandi ali
e suoni eterne musiche
e poi meravigliosamente ami
eppure hai scelto ancora me.
Ed anche se non so capire,
se non riesco a vincere,
verrà il momento di partire, tu lo sai:
ti prego non lasciarmi.
\endverse
\beginverse*
Ed anche se non so capire,
se non riesco a vincere,
verrà il momento di partire, tu lo sai:
ti prego non lasciarmi, ti prego,
non lasciarmi, ti prego non lasciarmi… mai.
\endverse
\endsong

% Fonte: "13 - RNS e adorazione/QUESTO IO CREDO.pdf" (foglio con accordi).
% Nell'elenco: Adorazione.
\beginsong{Questo io credo}[by={RnS}]
\beginverse
Tu, Padre immortale,
o Dio creatore,
Dio potente,
lo Spirito hai mandato
per generare il Figlio, Cristo Signore.
\endverse
\beginchorus
Credo in Te che sei Dio Padre,
credo in Te, Figlio Gesù,
credo in Te, Spirito Santo,
è trino il nostro Dio.
Credo che noi risorgeremo,
credo nell'eternità,
io credo in Te, nel tuo nome, o Cristo.
\endchorus
\beginverse
Tu, il nostro difensore,
su quella croce tu sei misericordia,
disceso negli abissi, risorto nella gloria,
tu regni in eterno.
\endverse
\beginverse*
E io credo in te,
credo che ritornerai,
credo che Gesù è il mio Signor. \rep{2}
\endverse
\endsong

% Fonte: testo fornito dall'utente ("il nostri" corretto in "i nostri"). Non presente nell'elenco. Autore non indicato.
\beginsong{Re dei re}
\beginverse
Hai sollevato i nostri volti dalla polvere,
le nostre colpe hai portato su di te.
Signore, ti sei fatto uomo
in tutto come noi per amore.
\endverse
\beginchorus
Figlio dell'Altissimo, povero tra i poveri,
vieni a dimorare tra noi.
Dio dell'impossibile, re di tutti i secoli,
vieni nella tua maestà.
Re dei re, i popoli ti acclamano,
i cieli ti proclamano re dei re.
Luce degli uomini,
regna col tuo amore tra noi, noi, noi, noi.
\endchorus
\beginverse
Ci hai riscattati dalla stretta delle tenebre
perché potessimo glorificare te.
Hai riversato in noi la vita
del tuo Spirito per amore.
\endverse
\beginchorus
Figlio dell'Altissimo, povero tra i poveri,
vieni a dimorare tra noi.
Dio dell'impossibile, re di tutti i secoli,
vieni nella tua maestà.
Re dei re, i popoli ti acclamano,
i cieli ti proclamano re dei re.
Luce degli uomini,
regna col tuo amore tra noi, noi, noi, noi.
\endchorus
\beginverse*
Tua è la gloria per sempre,
tua è la gloria per sempre.
Gloria, gloria,
gloria, gloria.
\endverse
\beginchorus
Figlio dell'Altissimo, povero tra i poveri,
vieni a dimorare tra noi.
Dio dell'impossibile, re di tutti i secoli,
vieni nella tua maestà.
Re dei re, i popoli ti acclamano,
i cieli ti proclamano re dei re.
Luce degli uomini,
regna col tuo amore tra noi, noi, noi, noi.
\endchorus
\endsong

% Fonte: LibrettoAccordi.pdf (Canti Coro), J39 (li' tra i canti di comunione).
% Nell'elenco: Adorazione, Ordinario. CD: Io credo.
% Omesse le indicazioni Uomini/Donne del finale.
\beginsong{Sono qui a lodarti}[by={RnS (CD Io credo)}]
\beginverse
Luce del mondo, nel buio del cuore vieni ed illuminami.
Tu mia sola speranza di vita resta per sempre con me.
\endverse
\beginchorus
Sono qui a lodarti, qui per adorarti, qui per dirti che Tu sei il mio Dio.
E solo Tu sei santo, sei meraviglioso, degno e glorioso sei per me.
\endchorus
\beginverse
Re della storia e Re nella gloria sei sceso in terra fra noi.
Con umiltà il tuo trono hai lasciato per dimostrarci il tuo amor.
\endverse
\beginverse*
Io mai saprò quanto ti costò lì sulla croce morir per me. \rep{4}
\endverse
\endsong

% Fonte: testo fornito dall'utente (ripetizioni riassunte).
% Nell'elenco: Adorazione.
\beginsong{Stai con me}[by={RnS}]
\beginverse
Stai con me, proteggimi,
coprimi con le tue ali, o Dio.
\endverse
\beginchorus
Quando la tempesta arriverà
volerò più in alto insieme a te.
Nelle avversità sarai con me
ed io saprò che tu sei il mio Re.
\endchorus
\beginverse
Il cuore mio riposa in te,
io vivrò in pace e verità.
\endverse
\beginverse*
Ed io saprò che tu sei il mio Re. \rep{2}
\endverse
\endsong

% Testo latino tradizionale (S. Tommaso d'Aquino, ultime due strofe del Pange lingua):
% non presente nei file di Canti Coro, scritto a memoria. Nell'elenco: senza tipologia.
\beginsong{Tantum ergo}[by={S. Tommaso d'Aquino}]
\beginverse
Tantum ergo Sacramentum
veneremur cernui:
et antiquum documentum
novo cedat ritui:
praestet fides supplementum
sensuum defectui.
\endverse
\beginverse
Genitori, Genitoque
laus et iubilatio,
salus, honor, virtus quoque
sit et benedictio:
procedenti ab utroque
compar sit laudatio. Amen.
\endverse
\endsong

% Fonte: LibrettoAccordi.pdf (Canti Coro), J38 (li' tra i canti di comunione).
% Nell'elenco: Adorazione, Ordinario.
\beginsong{Tu sei}[by={P. Spoladore}]
\beginverse
Tu sei la prima stella del mattino,
tu sei la nostra grande nostalgia,
tu sei il cielo chiaro dopo la paura,
dopo la paura di esserci perduti,
e tornerà la vita in questo mare.
\endverse
\beginchorus
Soffierà, soffierà il vento forte della vita,
soffierà sulle vele e le gonfierà di te! \rep{2}
\endchorus
\beginverse
Tu sei l'unico volto della pace,
tu sei speranza nelle nostre mani,
tu sei il vento nuovo sulle nostre ali,
sulle nostre ali soffierà la vita,
e gonfierà le vele per questo mare.
\endverse
\endsong

% Fonte: testo fornito dall'utente. Non presente nell'elenco; titolo confermato dall'utente.
% Strofa e ritornello si cantano due volte prima del ponte (seconda volta non riscritta: testo identico).
\beginsong{Tu sei la forza}[by={Eman}]
\beginverse*
Proprio quando sono qui con te
tu vinci per me le mie battaglie,
proprio quando sono qui con te
tu vinci per me le mie infermità.
\endverse
\beginverse*
In te, Dio, io trovo la forza
per non gettare la spugna,
perché Cristo ha donato il suo sangue.
In te, Dio, io trovo la forza
per non gettare la spugna,
perché Cristo è in me.
\endverse
\beginchorus
Tu sei la forza nella debolezza,
sei la speranza del cuore mio,
tu sei la certezza in un mondo che è senza:
tu sei il mio Dio, non dubito!
\endchorus
\beginverse*
E se, Gesù, tu sei con me,
chi sarà contro di me?
Se tu, Gesù, sarai con me
io vincerò comunque! \rep{4}
\endverse
\beginchorus
Tu sei la forza nella debolezza,
sei la speranza del cuore mio, del cuore mio!
Tu sei la forza nella debolezza,
sei la speranza del cuore mio,
tu sei la certezza in un mondo che è senza:
tu sei il mio Dio, non dubito! \rep{2}
\endchorus
\endsong

% Fonte: testo fornito dall'utente.
% Nell'elenco: Adorazione.
\beginsong{Tutto è possibile}[by={Nuovi Orizzonti}]
\beginverse
Questo è il luogo che Dio ha scelto per te,
questo è il tempo pensato per te.
Quella che vedi è la strada che Lui traccerà,
quello che senti l'amore che mai finirà.
\endverse
\beginchorus
E andremo e annunceremo che
in Lui tutto è possibile.
E andremo e annunceremo che
nulla ci può vincere,
perché abbiamo udito le sue parole,
perché abbiam veduto vite cambiare,
perché abbiamo visto l'amore vincere.
Sì, abbiamo visto l'amore vincere!
\endchorus
\beginverse
Questo è il momento che Dio ha atteso per te,
questo è il sogno che ha fatto su te.
Quella che vedi è la strada tracciata per te,
quello che senti, l'amore che ti accompagnerà.
\endverse
\beginverse*
Questo è il tempo che Dio ha scelto per te,
questo è il sogno che ha fatto su te.
\endverse
\endsong

% Fonte: testo fornito dall'utente (traduzione italiana ufficiale). Non presente nell'elenco.
% Originale: "Broken Vessels (Amazing Grace)". Pre-coro e coro ripetuti dopo la 2ª strofa.
\beginsong{Vasi rotti (Immensa grazia)}[by={J. Houston, J. Myrin (trad. F. Lorenzi, D. Minutella)}]
\beginverse
Tutti quei pezzi
rotti e dispersi
tu li hai raccolti,
riuniti in te.
A mani vuote
ma mai da solo,
mi hai liberato,
mi hai liberato.
\endverse
\beginverse*
Immensa grazia,
sei la speranza
che ha salvato uno come me.
Io mi ero perso
ma tu mi hai trovato,
ero cieco,
ora vedo te.
\endverse
\beginchorus
Ora posso vederti davvero,
negli occhi tuoi l'amore sincero.
Ti pieghi verso chi era ferito,
rialzi chi si sente spezzato.
\endchorus
\beginverse
I miei insuccessi,
le mie incoerenze:
metti un tesoro
in fragili vasi.
Prendi il mio cuore,
fanne una vela:
che il mondo veda
la tua vita in me.
\endverse
\beginverse*
Immensa grazia,
sei la speranza
che ha salvato uno come me.
Io mi ero perso
ma tu mi hai trovato,
ero cieco,
ora vedo te.
\endverse
\beginchorus
Ora posso vederti davvero,
negli occhi tuoi l'amore sincero.
Ti pieghi verso chi era ferito,
rialzi chi si sente spezzato.
\endchorus
\endsong

% Fonte: testo fornito dall'utente (traduzione italiana ufficiale). Originale: "You'll Come".
% Nell'elenco: "Verrai Signore" (Ingresso, Avvento); messo in Adorazione su indicazione dell'utente.
% Pre-coro scritto in coda alle strofe.
\beginsong{Verrai}[by={B. Ligertwood (trad. F. Lorenzi, N. Uelmen, D. Minutella)}]
\beginverse
Oggi ho deciso, ho scelto ormai,
Signore, aspetto Te.
Mia roccia e salvezza, mia ricompensa,
Signore, aspetto Te.
Il sole certo sorgerà
e verrai a noi,
l'aurora presto nascerà.
\endverse
\beginchorus
Verrai nella Tua gloria,
ci risponderai, oh Dio.
Vieni, Spirito,
disseta i nostri cuori col Tuo amore.
Verrai, verrai.
\endchorus
\beginverse
Non siamo scossi, restiamo saldi,
noi Ti aspettiamo, oh Dio.
Sei il liberatore, il vincitore,
noi Ti aspettiamo, oh Dio.
Il sole certo sorgerà
e verrai a noi,
la Tua Parola resterà.
\endverse
\beginchorus
Verrai nella Tua gloria,
ci risponderai, oh Dio.
Vieni, Spirito,
disseta i nostri cuori col Tuo amore.
Verrai, verrai.
\endchorus
\beginverse*
Mai più catene,
guariti in Te.
Ci fai vedere:
Dio, sei tra noi.
\endverse
\endsong

% Fonte: LibrettoAccordi.pdf (Canti Coro), J35.
% Nell'elenco: Adorazione, Ordinario. CD: Grazie (2012).
\beginsong{Voglio adorare te}[by={RnS (CD Grazie)}]
\beginchorus
Voglio adorare Te, voglio adorare Te,
voglio adorare Te, Signor, solo Te. \rep{2}
\endchorus
\beginverse
Nella gioia e nel dolore, nell'affanno della vita,
quando sono senza forze, adoro Te.
Nella pace e nell'angoscia, nella prova della croce,
quando ho sete del tuo amore, adoro Te, Signore.
\endverse
\beginverse
Nel coraggio e nel timore, nel tormento del peccato
quando il cuore mio vacilla, adoro Te.
Nella fede e nella grazia, nello zelo per il Regno,
quando esulto nel tuo nome, adoro Te, Signore.
\endverse
\endsong

\end{songs}

% !TEX root = ../libretto.tex
% Sezione: Canti allo Spirito
\momento{Canti allo Spirito}
\begin{songs}{titoli}

% Fonte: LibrettoAccordi.pdf (Canti Coro), L4.
% Nell'elenco: Canto allo Spirito / Ingresso, Ordinario.
\beginsong{Invochiamo la tua presenza}[by={RnS}]
\beginverse
Invochiamo la tua presenza vieni Signor.
Invochiamo la tua presenza scendi su di noi.
Vieni Consolatore dona pace e umiltà.
Acqua viva d'amore questo cuore apriamo a Te.
\endverse
\beginchorus
Vieni Spirito, vieni Spirito, scendi su di noi!
Vieni Spirito, vieni Spirito, scendi su di noi!
Vieni su noi Maranatha, vieni su noi Spirito!
Vieni Spirito, vieni Spirito, scendi su di noi!
Vieni Spirito, vieni Spirito, scendi su di noi, scendi su di noi.
\endchorus
\beginverse
Invochiamo la tua presenza, vieni Signor,
invochiamo la tua presenza scendi su di noi.
Vieni luce dei cuori dona forza e fedeltà.
Fuoco eterno d'amore questa vita offriamo a te.
\endverse
\endsong

% Fonte: LibrettoAccordi.pdf (Canti Coro), L2.
% Nell'elenco: Canto allo Spirito, Ordinario.
\beginsong{Lo Spirito di Cristo}
\beginchorus
Lo Spirito di Cristo, fa fiorire il deserto
torna la vita,
noi diventiamo testimoni di luce.
\endchorus
\beginverse
Non abbiamo ricevuto, uno spirito di schiavitù,
ma uno Spirito d'amore, uno Spirito di pace,
nel quale gridiamo: «Abbà Padre! Abbà Padre!»
\endverse
\beginverse
Lo Spirito che Cristo risuscitò,
darà vita ai nostri corpi, corpi mortali
e li renderà strumenti di salvezza, strumenti di salvezza.
\endverse
\beginverse
Non siamo più divisi e chiusi in noi stessi
non siamo più sterili e fuori dalla storia:
lo Spirito invade il cuore, ci rende nuova creatura.
\endverse
\beginverse
Sono venuto a portare il fuoco sulla terra
e come desidero che divampi nel mondo
e porti amore ed entusiasmo in tutti i cuori.
\endverse
\endsong

% Fonte: LibrettoAccordi.pdf (Canti Coro), A18.
% Nell'elenco: Ingresso, Pentecoste.
\beginsong{Vieni qui tra noi}[by={Gen Rosso, Gen Verde (CD Come fuoco vivo)}]
\beginchorus
Vieni qui tra noi, come fiamma che scende dal cielo.
Vieni qui tra noi, rinnova il cuore del mondo.
Vieni qui tra noi, col tuo amore rischiara la terra.
Vieni qui tra noi,
soffio di libertà.
\endchorus
\beginverse
Nel silenzio sei pace, nella notte luce,
Dio nascosto, vita, Dio tu sei, Amore.
Tutto si ricrea in te, tutto vive in te.
Scalda col tuo fuoco terra e cielo.
Tu, che sai raccogliere ogni gemito,
semina nel nostro cuore una speranza d'eternità.
\endverse
\beginverse*
Soffio di libertà,
Amore, Dio in mezzo a noi!
\endverse
\endsong

% Fonte: LibrettoAccordi.pdf (Canti Coro), L3.
% Nell'elenco: Canto allo Spirito, Ordinario.
\beginsong{Vieni, Santo Spirito (Sequenza)}[by={Gen Verde}]
\beginchorus
Vieni, Santo Spirito! Manda a noi dal cielo
un raggio di luce, un raggio di luce.
Vieni, padre dei poveri! Vieni, datore dei doni!
Luce dei cuori! Luce dei cuori.
\endchorus
\beginverse
Consolatore perfetto, ospite dolce dell'anima,
dolcissimo sollievo, dolcissimo sollievo.
Nella fatica, riposo, nella calura, riparo,
nel pianto, conforto, nel pianto, conforto.
\endverse
\beginverse
Luce beatissima, invadi fino all'intimo
il cuore dei tuoi fedeli, il cuore dei tuoi fedeli.
Senza la tua forza, nulla è nell'uomo,
nulla senza colpa, nulla senza colpa.
\endverse
\beginverse
Lava ciò che è sordido, bagna ciò che è arido,
ciò che sanguina risana, ciò che sanguina risana.
Piega ciò che è rigido, scalda ciò che è gelido,
drizza ciò ch'è sviato, drizza ciò ch'è sviato.
\endverse
\beginverse
Dona ai tuoi fedeli che solo in te confidano
i tuoi santi doni, i tuoi santi doni.
Dona virtù e premio, dona morte santa,
dona eterna gioia, dona eterna gioia.
\endverse
\endsong

% Fonte: testo fornito dall'utente (tolti gli accordi). Non presente nell'elenco.
% Messo tra i canti allo Spirito (sezione non indicata dall'utente).
\beginsong{Vieni, Spirito d'amore}[by={F. Volpintesta}]
\beginchorus
Vieni, Spirito d'amore,
manda un raggio di luce.
Dolce ospite dei cuori,
resta per sempre con noi.
\endchorus
\beginverse
Vieni, padre dei poveri,
vieni, datore dei doni,
vieni, luce dei cuori,
vieni, o Spirito Santo!
\endverse
\beginchorus
Vieni, Spirito d'amore,
manda un raggio di luce.
Dolce ospite dei cuori,
resta per sempre con noi.
\endchorus
\beginverse
Consolatore perfetto,
ospite dolce dell'anima,
vero sollievo dell'uomo,
vieni, o Spirito Santo!
\endverse
\beginchorus
Vieni, Spirito d'amore,
manda un raggio di luce.
Dolce ospite dei cuori,
resta per sempre con noi.
\endchorus
\beginverse
Nella fatica riposo,
nella calura riparo,
tu sei del pianto il conforto,
vieni, o Spirito Santo!
\endverse
\beginchorus
Vieni, Spirito d'amore,
manda un raggio di luce.
Dolce ospite dei cuori,
resta per sempre con noi.
\endchorus
\endsong

\end{songs}

% !TEX root = ../libretto.tex
% Sezione: Canoni (Taizé e altri)
\momento{Canoni}
\begin{songs}{titoli}

% Fonte: testo liturgico tradizionale dei canti di Taizé, scritto a memoria; verificato dall'utente.
\beginsong{Adoramus te, Christe}[by={Taizé}]
\beginverse*
Adoramus te, Christe, benedicimus tibi,
quia per crucem tuam redemisti mundum.
\endverse
\endsong

% Fonte: LibrettoAccordi.pdf (Canti Coro), O3.
% Nell'elenco: Canone.
\beginsong{Audite, beati}[by={F. Volpintesta}]
\beginverse*
Audite, beati, verbum Dei!
\endverse
\endsong

% Fonte: testo fornito dall'utente (solo ritornello).
% Non presente nell'elenco. Messo tra i canoni su indicazione dell'utente.
\beginsong{Awesome God}[by={R. Mullins}]
\beginverse*
Our God is an awesome God,
he reigns from heaven above
with wisdom, power and love:
our God is an awesome God.
\endverse
\endsong

% Fonte: testo liturgico tradizionale dei canti di Taizé, scritto a memoria; verificato dall'utente.
\beginsong{Bonum est confidere}[by={Taizé}]
\beginverse*
Bonum est confidere in Domino,
bonum sperare in Domino.
\endverse
\endsong

% Fonte: testo liturgico tradizionale dei canti di Taizé, scritto a memoria; verificato dall'utente.
\beginsong{Christe, lux mundi}[by={Taizé}]
\beginverse*
Christe, lux mundi,
qui sequitur te
habebit lumen vitae, lumen vitae.
\endverse
\endsong

% Fonte: testo liturgico tradizionale dei canti di Taizé, scritto a memoria; verificato dall'utente.
\beginsong{Confitemini Domino}[by={Taizé}]
\beginverse*
Confitemini Domino, quoniam bonus.
Confitemini Domino, alleluia.
\endverse
\endsong

% Fonte: testo liturgico tradizionale dei canti di Taizé, scritto a memoria; verificato dall'utente.
\beginsong{Crucem tuam}[by={Taizé}]
\beginverse*
Crucem tuam adoramus, Domine,
resurrectionem tuam laudamus, Domine.
Laudamus et glorificamus.
Resurrectionem tuam laudamus, Domine.
\endverse
\endsong

% Fonte: testo fornito dall'utente. Non presente nell'elenco.
\beginsong{Evangelizare pauperibus}[by={S. Zampa}]
\beginverse*
Evangelizare pauperibus
misit me, misit me.
Evangelizare pauperibus.
\endverse
\endsong

% Fonte: testo liturgico tradizionale dei canti di Taizé, scritto a memoria; verificato dall'utente.
\beginsong{Gloria, gloria in excelsis Deo}[by={Taizé}]
\beginverse*
Gloria, gloria in excelsis Deo!
Gloria, gloria, alleluia, alleluia!
\endverse
\endsong

% Fonte: testo liturgico tradizionale dei canti di Taizé, scritto a memoria; verificato dall'utente.
\beginsong{Il Signore è la mia forza}[by={Taizé}]
\beginverse*
Il Signore è la mia forza
e io spero in lui.
Il Signor è il Salvator.
In lui confido, non ho timor,
in lui confido, non ho timor.
\endverse
\endsong

% Fonte: testo liturgico tradizionale dei canti di Taizé, scritto a memoria; verificato dall'utente.
\beginsong{In manus tuas, Pater}[by={Taizé}]
\beginverse*
In manus tuas, Pater,
commendo spiritum meum.
\endverse
\endsong

% Fonte: testo liturgico tradizionale dei canti di Taizé, scritto a memoria; verificato dall'utente.
\beginsong{Jubilate Deo}[by={Taizé}]
\beginverse*
Jubilate Deo omnis terra.
Servite Domino in laetitia.
Alleluia, alleluia, in laetitia.
Alleluia, alleluia, in laetitia.
\endverse
\endsong

% Fonte: testo liturgico tradizionale dei canti di Taizé, scritto a memoria; verificato dall'utente.
\beginsong{Laudate Dominum}[by={Taizé}]
\beginverse*
Laudate Dominum, laudate Dominum,
omnes gentes, alleluia.
\endverse
\endsong

% Fonte: testo fornito dall'utente. Non presente nell'elenco. Autore indicato dall'utente.
\beginsong{Let it rain}[by={M. W. Smith}]
\beginverse*
Let it rain, let it rain,
open the floodgates of Heaven.
\endverse
\endsong

% Fonte: testo liturgico tradizionale dei canti di Taizé, scritto a memoria; verificato dall'utente.
\beginsong{Magnificat}[by={Taizé}]
\beginverse*
Magnificat, magnificat,
magnificat anima mea Dominum.
Magnificat, magnificat,
magnificat anima mea.
\endverse
\endsong

% Fonte: testo liturgico tradizionale dei canti di Taizé, scritto a memoria; verificato dall'utente.
\beginsong{Misericordias Domini}[by={Taizé}]
\beginverse*
Misericordias Domini
in aeternum cantabo.
\endverse
\endsong

% Fonte: testo liturgico tradizionale dei canti di Taizé, scritto a memoria; verificato dall'utente.
% Nell'elenco: "Nada te turbe/Nulla ti turbi", Canone. Testo di s. Teresa d'Avila.
\beginsong{Nada te turbe}[by={Taizé}]
\beginverse*
Nada te turbe, nada te espante:
quien a Dios tiene nada le falta.
Nada te turbe, nada te espante:
solo Dios basta.
\endverse
\endsong

% Fonte: testo liturgico tradizionale dei canti di Taizé, scritto a memoria; verificato dall'utente.
% Versione italiana di "Nada te turbe".
\beginsong{Nulla ti turbi}[by={Taizé}]
\beginverse*
Nulla ti turbi, nulla ti spaventi:
chi ha Dio niente gli manca.
Nulla ti turbi, nulla ti spaventi:
solo Dio basta.
\endverse
\endsong

% Fonte: testo liturgico tradizionale dei canti di Taizé, scritto a memoria; verificato dall'utente.
\beginsong{Oculi nostri}[by={Taizé}]
\beginverse*
Oculi nostri ad Dominum Jesum,
oculi nostri ad Dominum nostrum.
\endverse
\endsong

% Fonte: testo fornito dall'utente.
\beginsong{Salvator mundi}[by={Taizé}]
\beginverse*
Salvator mundi, salva nos,
Salvator mundi, salva nos.
Salva nos, salva nos,
Salvator mundi, salva nos.
\endverse
\endsong

% Fonte: testo liturgico tradizionale dei canti di Taizé, scritto a memoria; verificato dall'utente.
\beginsong{Ubi caritas}[by={Taizé}]
\beginverse*
Ubi caritas et amor,
ubi caritas Deus ibi est.
\endverse
\endsong

% Fonte: testo liturgico tradizionale dei canti di Taizé, scritto a memoria; verificato dall'utente.
\beginsong{Veni Sancte Spiritus (Tui amoris ignem)}[by={Taizé}]
\beginverse*
Veni Sancte Spiritus,
tui amoris ignem accende.
Veni Sancte Spiritus,
veni Sancte Spiritus.
\endverse
\endsong

\end{songs}

% !TEX root = ../libretto.tex
% Sezione: Avvento
% Regola: i brani con periodo Avvento vanno qui anche se nell'elenco
% sono di Ingresso/Comunione/Finale.
\momento{Avvento}
\begin{songs}{titoli}

% Fonte: LibrettoAccordi.pdf (Canti Coro), L5.
% Nell'elenco: Canto allo Spirito, Avvento/Natale.
\beginsong{Come brezza}[by={D. Ricci}]
\beginverse
Spirito di luce pura che parli nell'anima,
Spirito di gioia e vita che adombrasti Maria,
Spirito impetuoso e forte, che palpiti ora in noi
come brezza leggera, leggera…
\endverse
\beginverse
Spirito consolatore, carezza dell'anima,
Spirito di fiamma e fuoco che trascini con te,
Spirito che ci travolgi e sussurri impercettibile
come brezza leggera, leggera…
\endverse
\beginchorus
Che non passi mai, un solo giorno senza amarti!
Che non passi senza di te,
senza la bellezza limpida che santifica ogni cosa in noi
e illumina l'oscurità.
\endchorus
\beginverse
Tu che all'alba del creato sull'acque e gli oceani
Aleggiavi come soffio potente di Dio.
Tu parola dei profeti, del Regno, tu l'anelito
come brezza leggera, leggera…
\endverse
\beginverse
Refrigerio di ogni pena, pensiero che illumina,
Tu che hai dato al mondo il Verbo, il cuore di Dio.
Dacci d'ascoltare te, il tuo respiro dentro noi
come brezza leggera, leggera…
\endverse
\endsong

% Fonte: LibrettoAccordi.pdf (Canti Coro), J9.
% Nell'elenco: Comunione, Avvento.
\beginsong{Dal tronco di Jesse}[by={F. Baggio}]
\beginchorus
Dal tronco di Jesse germoglierà un nuovo virgulto domani,
dalle sue radici si eleverà un albero nuovo.
\endchorus
\beginverse
Su di Lui scenderà lo Spirito di Dio, gli regalerà i suoi ricchi doni:
consiglio e sapienza, scienza e fortezza, santo timore di Dio.
\endverse
\beginverse
Non giudicherà per le apparenze, non deciderà per sentito dire.
Ai poveri poi darà con larghezza, farà giustizia agli oppressi.
\endverse
\beginverse
Ed il lupo e l'agnello in pace vivranno, saranno amici la mucca e il leone.
Ed un fanciullo li guiderà, pascoleranno insieme.
\endverse
\endsong

% Fonte: LibrettoAccordi.pdf (Canti Coro), A11.
% Nell'elenco: Ingresso, Avvento.
\beginsong{Gesù verrà}[by={RnS}]
\beginverse
Grande gioia ci sarà,
la speranza rivivrà,
la tua vita cambierà,
devi crederlo perché Gesù verrà.
Il tuo deserto sparirà,
nuove vie lui ti aprirà,
la tristezza passerà
e la notte più non tornerà,
Gesù verrà.
\endverse
\beginchorus
Gesù verrà,
la sua gloria apparirà,
Gesù verrà,
come il sole splenderà,
Gesù verrà,
e con noi Lui resterà,
come una stella Lui sarà ed il tuo cuore orienterà.
Gesù verrà. \rep{2}
\endchorus
\beginverse
La città di Dio verrà,
per gli eletti si aprirà,
la giustizia regnerà
e la pace tornerà, Gesù verrà.
Ogni uomo adorerà
Gesù Cristo che verrà,
e lo Spirito sarà
l'acqua viva che ci riempirà,
Gesù verrà.
\endverse
\endsong

% Fonte: LibrettoAccordi.pdf (Canti Coro), K6.
% Nell'elenco: Finale, Avvento.
\beginsong{Innalzate nei cieli}[by={Albisetti, Martorell}]
\beginverse
Innalzate nei cieli lo sguardo: la salvezza di Dio è vicina.
Risvegliate nei cuori l'attesa per accogliere il Re della Gloria.
\endverse
\beginchorus
Vieni, Gesù! Vieni, Gesù! Discendi dal cielo! Discendi dal cielo!
Vieni, Gesù!
\endchorus
\beginverse
Sorgerà dalla casa di David il Messia da tutti invocato:
prenderà da una Vergine il corpo per potenza di Spirito Santo.
\endverse
\beginverse
Benedetta sei tu, Maria, che rispondi all'attesa del mondo:
come aurora splendente di grazia porti al mondo il sole divino.
\endverse
\beginverse
Vieni, o Re, discendi dal cielo, porta al mondo il sorriso di Dio:
nessun uomo ha visto il Suo volto, solo Tu puoi svelarci il mistero.
\endverse
\endsong

% Fonte: LibrettoAccordi.pdf (Canti Coro), J27.
% Nell'elenco: Comunione, Avvento.
\beginsong{O cieli, piovete dall'alto}[by={K. Arguello}]
\beginchorus
O cieli, piovete dall'alto! O nubi, mandateci il Santo.
O terra, apriti, o terra, e germina il Salvator!
\endchorus
\beginverse
Siamo il deserto, siamo l'arsura: Maranatha! Maranatha!
Siamo il vento, nessuno ci ode: Maranatha! Maranatha!
\endverse
\beginverse
Siamo le tenebre, nessuno ci guida: Maranatha! Maranatha!
Siam le catene, nessuno ci scioglie: Maranatha! Maranatha!
\endverse
\beginverse
Siamo il freddo, nessuno ci copre: Maranatha! Maranatha!
Siamo la fame, nessuno ci nutre: Maranatha! Maranatha!
\endverse
\beginverse
Siamo le lacrime, nessuno ci asciuga: Maranatha! Maranatha!
Siamo il dolore, nessuno ci guarda: Maranatha! Maranatha!
\endverse
\endsong

% Fonte: LibrettoAccordi.pdf (Canti Coro), A22.
% Nell'elenco: Ingresso, Avvento. CD: Rallegratevi (2015).
\beginsong{Prepariamo la via}[by={P. Dargenio, M. Zambuto (CD Rallegratevi)}]
\beginchorus
Prepariamo la via
al Signore che viene a salvarci.
Raddrizziamo i nostri sentieri
e contempleremo il suo volto.
Prepariamo la via
al Signore che viene a salvarci.
Appianiamo le strade di Dio
che viene a donare la vita.
\endchorus
\beginverse
C'è una voce che grida:
«Preparate la strada,
appianate i sentieri
che portano al Regno di Dio.
Il Signore è vicino,
si ode già la sua voce:
sollevate lo sguardo e correte da Lui perché adesso è qui.»
\endverse
\beginverse
Come lampada accesa
la tua voce ci guida:
appianate i sentieri
che portano al Regno di Dio.
Ora è tempo di gioia,
di far sciogliere il canto:
il Signore del cielo cammina con noi, Egli adesso è qui.
\endverse
\endsong

% Fonte: LibrettoAccordi.pdf (Canti Coro), J21.
% Nell'elenco: Comunione, Avvento. CD: Maranatha, Vieni Signore (2002).
\beginsong{Signore del cielo}[by={D. Ricci (CD Maranatha, Vieni Signore)}]
\beginverse
Nell'anima scende il tuo respiro e l'infinito è in mezzo a noi.
Nell'anima scenderà l'immensità, l'amore in mezzo a noi.
E l'anima canta la tua presenza e il paradiso è in mezzo a noi.
E l'anima canterà la verità, la vita in mezzo a noi.
\endverse
\beginchorus
Signore del cielo, degli abissi sereni,
tu, Signore, che vieni per essere con noi che siamo in questa valle!
Signore del cielo, delle altezze più sante,
tu, Signore viandante, cammini a fianco a noi che siamo in questa valle!
\endchorus
\beginverse
E l'anima splende di sole e fiamma e l'universo è in mezzo a noi.
E l'anima splenderà e accenderà la luce in mezzo a noi.
\endverse
\beginverse
Nell'anima un vento che spira lieve: la primavera è in mezzo a noi.
Nell'anima un vento che è la libertà, la gioia in mezzo a noi.
\endverse
\endsong

% Fonte: testo fornito dall'utente.
% Nell'elenco: Ingresso, Avvento (autore in elenco: Ecco perché canto).
% Nell'ultimo ritornello la seconda parte si ripete.
\beginsong{Verrai Signore}[by={A. Benedetti, G. Anselmi}]
\beginverse
In silenzio attendiamo Te,
saremo pronti, saremo pronti.
Sveglieremo la fede in noi,
la Tua promessa resta salda.
\endverse
\beginchorus
E verrai, verrai, Signore,
Tu verrai, verrai per noi!
Ci sarà un sentiero, ci sarà una via.
Resterai per sempre nella nostra vita,
Tu verrai, verrai, verrai!
\endchorus
\beginverse
La Tua voce si leverà
fino ai confini della terra,
la Tua gioia ci invaderà
e non avremo più paura.
\endverse
\beginchorus
E verrai, verrai, Signore,
Tu verrai, verrai per noi!
Ci sarà un sentiero, ci sarà una via.
Resterai per sempre nella nostra vita,
Tu verrai, verrai, verrai!
Ci sarà un sentiero, ci sarà una via.
Resterai per sempre nella nostra vita,
Tu verrai, verrai, verrai!
\endchorus
\endsong

% Fonte: LibrettoAccordi.pdf (Canti Coro), J41.
% Nell'elenco: Comunione, Avvento. CD: Rallegratevi (2015).
% Due ritornelli alternati: dopo ogni strofa e' indicato quale.
\beginsong{Viene in mezzo a noi}[by={P. Dargenio (CD Rallegratevi)}]
\beginchorus
\textit{Rit. 1:} Ecco, viene in mezzo a noi come un'alba nuova,
la speranza dell'umanità.
\textit{Rit. 2:} Ecco, si apre in mezzo a noi una strada nuova:
il Signore è qui. Alleluia!
\endchorus
\beginverse
Vieni in mezzo a noi, Sapienza dell'Altissimo,
la tua dimora si estende ai confini del mondo.
Tu, che disponi con forza e tutto governi con soavità,
vieni a insegnarci la strada che porta da te, alla saggezza. \textit{(Rit. 1)}
\endverse
\beginverse
Vieni in mezzo a noi, Signore d'Israele,
tu che parlasti a Mosè dal roveto ardente.
Tu, che dall'alto del monte ci hai donato la Legge,
vieni a mostrarci la strada per la libertà con la tua potenza. \textit{(Rit. 2)}
\endverse
\beginverse
Vieni in mezzo a noi, Germoglio di Jesse,
tu che ti innalzi a vessillo glorioso nel buio.
A te rivolgono gli occhi tutte le genti del mondo.
Vieni a donare al tuo popolo la libertà, non tardare. \textit{(Rit. 1)}
\endverse
\beginverse
Vieni in mezzo a noi, Chiave di Davide,
tu che spalanchi le porte del regno di Dio.
Tu sei lo scettro potente che scioglie le nostre catene,
vieni a sbarrare le porte dell'oscurità con la tua forza. \textit{(Rit. 2)}
\endverse
\beginverse
Vieni in mezzo a noi, Stella dell'Oriente,
tu che rischiari di luce la notte del mondo.
Illumina il cuore dell'uomo che, nella sua colpa, ricerca te,
vieni a ridare speranza all'umanità, che in te risplende. \textit{(Rit. 1)}
\endverse
\beginverse
Vieni in mezzo a noi, Atteso delle genti,
Pietra angolare che unisci i popoli in uno.
Tu che sei l'Emmanuele nostro fratello, il Dio con noi,
vieni a donare la vita per l'umanità che non ha pace. \textit{(Rit. 1 e 2)}
\endverse
\endsong

% Fonte: LibrettoAccordi.pdf (Canti Coro), A16.
% Nell'elenco: Ingresso, Avvento. CD: Maranatha, Vieni Signore (2011).
\beginsong{Vieni Gesù, Maranatha}[by={F. Baggio (CD Maranatha, Vieni Signore)}]
\beginchorus
Vieni Gesù, Maranatha, speranza di pace per noi.
Vieni Gesù, Maranatha, rivelaci il volto di Dio.
\endchorus
\beginverse
Fra tanto dolore di uomini persi, ti prego, rinasci, Gesù!
Asciuga le lacrime amare di morte e schiudici l'eternità.
\endverse
\beginverse
Fra tanti conflitti e minacce di guerra, riporta la pace, Gesù!
A tutti gli oppressi traditi dall'uomo ridona la libertà.
\endverse
\beginverse
Fra tante ingiustizie e occulti poteri, insegnaci la tua umiltà.
Sei tu il solo giudice dell'universo per tutta l'eternità.
\endverse
\endsong

\end{songs}

% !TEX root = ../libretto.tex
% Sezione: Natale
\momento{Natale}
\begin{songs}{titoli}

% Fonte: LibrettoAccordi.pdf (Canti Coro), N6.
% Nell'elenco: Comunione, Natale.
\beginsong{Adeste, fideles}[by={Canto popolare}]
\beginverse
Adeste, fideles, laeti triumphantes,
venite, venite in Bethlem:
natum videte Regem angelorum.
\endverse
\beginchorus
Venite, adoremus, venite adoremus,
venite adoremus Dominum.
\endchorus
\beginverse
En grege relicto, humiles ad cunas
vocati pastores adproperant:
et nos ovanti gradu festinemus.
\endverse
\beginverse
Aeterni parentis splendorem aeternum
velatum sub carne videbimus:
Deum infantem pannis involutum.
\endverse
\endsong

% Fonte: LibrettoAccordi.pdf (Canti Coro), E1.
% Nell'elenco: Alleluia, Natale.
\beginsong{Alleluia, è nato Gesù}[by={F. Baggio}]
\beginchorus
Alleluja! Alleluja! Alleluja!
Alleluja! Alleluja! Oggi è nato Gesù!
\endchorus
\beginverse
Al mondo annunciamo la gioia più grande: è nato Cristo Signore.
È l'Emanuele, il Dio con noi, è luce per l'umanità.
\endverse
\beginverse
Un giorno santo è spuntato per noi, venite, adoriamo il Signore.
Da tutta la terra un canto di gioia s'innalza al Dio fatto uomo.
\endverse
\endsong

% Fonte: LibrettoAccordi.pdf (Canti Coro), N3.
% Nell'elenco: Comunione, Natale.
\beginsong{Astro del ciel}[by={Canto popolare}]
\beginverse
Astro del Ciel, pargol divin, mite agnello, Redentor,
tu che i Vati da lungi sognar, tu che angeliche voci nunziar,
\endverse
\beginchorus
Luce dona alle menti, pace infondi nei cuor. \rep{2}
\endchorus
\beginverse
Astro del Ciel, pargol divin, mite agnello, Redentor,
tu di stirpe regale decor, tu virgineo, mistico fior,
\endverse
\beginverse
Astro del Ciel, pargol divin, mite agnello, Redentor,
tu disceso a scontare l'error, tu sol nato a parlare d'amor,
\endverse
\endsong

% Fonte: LibrettoAccordi.pdf (Canti Coro), K11.
% Nell'elenco: Finale, Natale.
\beginsong{Cantate al Signore}[by={D. Ricci}]
\beginverse
Cantate al Signore un canto nuovo da tutta la terra!
Voi cantate al Signore, cantate a lui. Benedite il nome suo.
Annunziate ogni giorno di lui la salvezza e la gloria.
A tutti i popoli dite i prodigi di lui. Benedite il nome suo.
\endverse
\beginchorus
È nato per noi, il Signore per noi oggi è nato.
Lui è nato per noi, il salvatore del mondo è nato.
\endchorus
\beginverse
Esultino i campi, esulti la terra, esultino i mari,
esultino i cieli cantando di lui! Benedite il nome suo.
E sia gioia, perché lui stende già sulla terra la mano sua:
a tutti i popoli la sua giustizia darà. Benedite il nome suo.
\endverse
\endsong

% Fonte: LibrettoAccordi.pdf (Canti Coro), J37.
% Nell'elenco: Comunione, Natale.
\beginsong{Emmanuel}[by={G. Cento}]
\beginchorus
Emmanuel, Emmanuel,
Emmanuel è Dio con noi, è Dio con noi. \rep{2}
\endchorus
\beginverse
Il Dio lontano si è fatto vicino, ha preso un volto come noi.
Lui pensa da uomo, lavora da uomo e ama da uomo: è uno di noi.
\endverse
\beginverse
È nato un bambino, è nato il Signore, è nato il nostro Redentore.
La terra è in festa, un uomo crede, la Chiesa canta: Alleluia!
\endverse
\beginverse
Abbiamo toccato il Verbo di vita, abbiamo parlato insieme a lui.
Ha dato la vita per l'uomo che crede, per tutta la nostra umanità.
\endverse
\endsong

% Fonte: LibrettoAccordi.pdf (Canti Coro), N2.
% Nell'elenco: Finale, Natale.
\beginsong{Gli angeli delle campagne}[by={Canto popolare}]
\beginverse
Gli angeli delle campagne, cantano l'inno: Gloria in ciel!
E l'eco delle montagne ripete il canto dei fedel.
\endverse
\beginchorus
Gloria in excelsis Deo! \rep{2}
\endchorus
\beginverse
Oh pastori che cantate, dite il perché di tanto onor.
Qual Signore, qual Profeta merita questo gran splendor?
\endverse
\beginverse
Oggi è nato in una stalla, nella notturna oscurità.
Egli è il Verbo, s'è incarnato, e venne in questa povertà.
\endverse
\endsong

% Fonte: LibrettoAccordi.pdf (Canti Coro), N5.
% Nell'elenco: Comunione, Natale.
\beginsong{Lieto annuncio}[by={Canto popolare}]
\beginverse
Lieto annuncio risuonò, luce nuova sfolgorò: Cristo è nato all'umanità.
\endverse
\beginchorus
Hai vinto ormai le tenebre: sei nato Re del mondo. \rep{2}
\endchorus
\beginverse
Gioia piena brilla in noi; vive qui l'eternità; nasce nuova umanità.
\endverse
\endsong

% Fonte: LibrettoAccordi.pdf (Canti Coro), G8.
% Nell'elenco: Offertorio, Natale.
\beginsong{Nasce l'amor}[by={Gen Rosso}]
\beginverse
Fredda è la notte, buio d'intorno, silente attesa, sacro timor.
Brilla una stella, veglia la terra, nasce la vita, palpita un cuor.
\endverse
\beginchorus
Nasce Gesù, nasce l'Amor. Gli angeli in cielo cantano in cor. \rep{2}
\endchorus
\beginverse
E se l'amore regna fra noi, nei nostri cuori vive Gesù.
Splende la luce, in tutti è pace: il paradiso è in mezzo a noi.
\endverse
\beginverse*
Nato è Gesù, nato è l'Amor. Gli angeli in cielo cantano in cor. \rep{2}
\endverse
\endsong

% Fonte: LibrettoAccordi.pdf (Canti Coro), K9.
% Nell'elenco: Finale, Natale.
\beginsong{Oggi è nato}[by={D. Ricci}]
\beginverse
Oggi per noi è nato il Salvatore. Tutta la terra canta a lui, benedice lui.
E ogni giorno si leva un canto, canto nuovo di gioia e vita.
E ogni giorno si narra la gloria ed osannano i popoli che oggi è nato.
\endverse
\beginchorus
Tutto il mondo canta alleluja. Anche il cielo canta alleluja.
Canta la terra, canta il mare
e nei campi le messi rispondono agli alberi. \rep{2}
\endchorus
\beginverse*
Tutto il mondo canta alleluja. Anche il cielo canta alleluja.
Canta la terra, canta il mare,
alleluja, alleluja, alleluja, alleluja.
\endverse
\endsong

% Fonte: LibrettoAccordi.pdf (Canti Coro), N1. Spazi tra le parole latine ricostruiti.
% Nell'elenco: Comunione, Natale (canto popolare).
\beginsong{Resonet in laudibus}[by={Canto popolare}]
\beginverse
Resonet in laudibus cum iucundis plausibus Sion cum fidelibus.
\endverse
\beginchorus
Apparuit, apparuit quem genuit Maria. \rep{2}
\endchorus
\beginverse
Sion lauda Dominum, Salvatorem omnium: Virgo parit Filium.
\endverse
\beginverse
Pueri concurrite, nato Regi psallite, voce pia dicite.
\endverse
\beginverse
Ergo nostra concio cum sit plena gaudio benedicat Domino.
\endverse
\beginverse
Et nos unanimiter proclamemus dulciter ipse Summus Arbiter.
\endverse
\endsong

% Versione popolare (3 strofe). Fonte: spartito Singing Bell,
% singing-bell.com/wp-content/uploads/2015/11/Tu-scendi-dalle-stelle_Singing-Bell.pdf
% Testo originale completo (7 strofe): it.wikipedia.org/wiki/Tu_scendi_dalle_stelle
% Nell'elenco: Comunione, Natale.
\beginsong{Tu scendi dalle stelle}[by={S. Alfonso Maria de' Liguori}]
\beginverse
Tu scendi dalle stelle, o Re del cielo,
e vieni in una grotta al freddo e al gelo. \rep{2}
O Bambino mio divino,
io ti vedo qui a tremar;
o Dio beato!
Ah, quanto ti costò l'avermi amato! \rep{2}
\endverse
\beginverse
A te, che sei del mondo il Creatore,
mancano panni e fuoco, o mio Signore. \rep{2}
Caro eletto pargoletto,
quanto questa povertà
più m'innamora,
giacché ti fece amor povero ancora. \rep{2}
\endverse
\beginverse
Tu lasci il bel gioir del divin seno,
per giungere a penar su questo fieno. \rep{2}
Dolce amore del mio cuore,
dove amor ti trasportò?
O Gesù mio,
perché tanto patir? Per amor mio! \rep{2}
\endverse
\endsong

\end{songs}

% !TEX root = ../libretto.tex
% Sezione: Quaresima e Pasqua
\momento{Quaresima e Pasqua}
\begin{songs}{titoli}

% Fonte: LibrettoAccordi.pdf (Canti Coro), O14; testo e struttura confermati dall'utente.
% Nell'elenco: Ceneri (l'elenco dice Buttazzo, ma l'utente conferma RnS).
\beginsong{Abbà misericordia}[by={RnS}]
\beginverse
Non sono degno di essere qui, ho abbandonato la tua casa,
ho dissipato i tuoi beni,
Padre, ho peccato contro il cielo e contro di Te.
\endverse
\beginchorus
Abbà, misericordia, Abbà, misericordia, Abbà, Abbà.
Abbà, misericordia, Abbà, misericordia, Abbà, Abbà, Abbà.
\endchorus
\beginverse
Non sono degno di esser tuo figlio, in Gesù Cristo rialzami
e in Lui ridammi la vita,
Padre, ho peccato contro il cielo e contro di Te.
\endverse
\beginverse
Non sono degno del tuo amore, riempi il mio cuore del tuo Spirito,
insieme a Te farò festa per sempre,
Padre, ho peccato contro il cielo e contro di Te.
\endverse
\endsong

% Fonte: LibrettoAccordi.pdf (Canti Coro), E8.
% Nell'elenco: Alleluia, Pasqua.
\beginsong{Alleluia, Cristo è risorto veramente}[by={F. Buttazzo}]
\beginchorus
Alleluia! Alleluja! Alleluja! Alleluja!
Cristo è risorto veramente. Alleluja! Alleluja! \rep{2}
\endchorus
\beginverse
Cristo ha inaugurato il suo regno,
salvezza e vita sono con lui.
Rallegriamoci, esultiamo con lui:
la morte è stata vinta, la vita ora trionfa.
\endverse
\endsong

% Fonte: "99 - Raccolte e libretti accordi/Il mistero pasquale_ gv.pdf" (pag. 77).
% Nell'elenco: Comunione, Settimana Santa. CD: Il mistero pasquale (2011).
\beginsong{Amore abbandonato}[by={Gen Verde; testo P. A. Sequeri (CD Il mistero pasquale)}]
\beginverse
Giorno verrà
del tuo ritorno in festa:
un vino nuovo porterai
ed io berrò con te.
Come il tralcio resto in te
e accolgo in me la vita.
Dal tuo ferito amore io
non mi separerò.
\endverse
\beginverse
Canto per te
che doni la tua vita:
nel sangue mi purifichi,
agnello immacolato.
Del mio cuore arido
tu spezzi le catene.
Più nulla ci separerà:
un cuore nuovo avrò.
\endverse
\beginverse
Resta con me
la sera della vita:
nel segno della croce io
ti riconoscerò.
Nel mio amore fragile
tu accendi la speranza.
Amore abbandonato, tu
non mi abbandonerai.
\endverse
\endsong

% Fonte: LibrettoAccordi.pdf (Canti Coro), O18.
% Nell'elenco: Finale, Settimana Santa.
\beginsong{Ave verum}[by={Gen Verde (CD Il mistero pasquale)}]
\beginchorus
Ave, verum corpus,
natum de Maria virgine,
vere passum immolatum
in cruce pro homine.
Cujus latus perforatum
unda fluxit et sanguine,
esto nobis praegustatum
in mortis examine.
O Iesu dulcis, o Iesu pie,
o Iesu, fili Mariae,
miserere, miserere, miserere mei.
\endchorus
\beginverse*
Amen. Amen. Amen. Amen.
\endverse
\endsong

% Fonte: LibrettoAccordi.pdf (Canti Coro), O21.
% Nell'elenco: Pasqua.
\beginsong{Canto dell'acqua}[by={Gen Verde (CD Il mistero pasquale)}]
\beginverse
L'acqua che sommerge il buio delle colpe scorre qui.
L'acqua che ci rende il cuore trasparente scorre qui.
È zampillo nato da roccia, è sorgente di fecondità,
è fontana del tuo giardino, onda di libertà.
L'acqua che ci immerge dentro la tua morte scorre qui.
L'acqua che ci dà la vita che risorge scorre qui.
È ruscello nato da piaga, fresco fuoco che ci inonderà,
acqua viva e bianca rugiada, goccia d'eternità.
L'acqua viva, grembo della vita nuova, scorre qui.
\endverse
\endsong

% Fonte: "15 - Quaresima, Triduo, Pasqua/come è grande il tuo amore.pdf" (foglio accordi).
% Nell'elenco: Ingresso, Quaresima (autore in elenco: Ecco perché canto).
\beginsong{Com'è grande il tuo amore}[by={A. Benedetti}]
\beginverse
Siamo qui
per cercare il tuo volto, o Dio,
tu, sorgente di vita sei,
acqua viva che sgorga dalla roccia.
Solo in te
ritroviamo il perdono e ci rialziamo di nuovo
se ci tendi la tua mano ancora, Gesù.
\endverse
\beginchorus
Com'è grande il tuo amore
per tutte le creature,
tu ci ami, Signore,
e non ci abbandoni mai. \rep{2}
\endchorus
\beginverse
Su di noi
hai posato lo sguardo, o Dio,
ci hai chiamati per nome, tu
sei la luce che illumina la notte.
Solo in te
ritroviamo il perdono e ci rialziamo di nuovo
se ci tendi la tua mano ancora, Gesù.
\endverse
\beginverse*
E non ci abbandoni mai,
tu non ci abbandoni.
\endverse
\endsong

% Fonte: LibrettoAccordi.pdf (Canti Coro), O2.
% Nell'elenco: Lavanda, Settimana Santa.
\beginsong{Come io ho amato voi}[by={A. Porfiri}]
\beginverse
Amatevi a vicenda come io ho amato voi,
così dice il Signore:
niente divisioni, liti o discordie,
l'amore regni tra voi.
\endverse
\beginverse
Servitevi a vicenda come io ho servito voi,
così dice il Signore:
chi di voi è primo ultimo sarà
e l'amore regni tra voi.
\endverse
\beginverse
Donatevi a vicenda come io ho fatto a voi,
così dice il Signore:
siete un solo corpo, Popolo di Dio,
l'amore regni tra voi.
\endverse
\endsong

% Fonte: LibrettoAccordi.pdf (Canti Coro), G17.
% Nell'elenco: Offertorio, Pasqua.
\beginsong{Come un fiume}[by={Gen Verde (CD Il mistero pasquale)}]
\beginverse
Come un fiume il frutto della tua vittoria
viene a Te.
Tutto quello che ci hai dato nella storia
diamo a Te.
Ecco il mondo, immensa patena
dei fermenti dell'umanità;
ecco il pane ardente dei giorni
che il tuo amore ci dà.
Ecco il mondo, calice immenso
d'ogni pena, d'ogni povertà;
Ecco il vino ardente del tempo,
questo tempo che va.
\endverse
\endsong

% Fonte: LibrettoAccordi.pdf (Canti Coro), O13.
% Nell'elenco: Ingresso, Ceneri.
\beginsong{Dal profondo, o Signore}[by={F. Buttazzo}]
\beginchorus
Dal profondo, o Signore, gridiamo il tuo nome:
dalle strade del mondo vogliamo tornare a te.
Dal profondo, o Signore, gridiamo il tuo nome:
solo in te la vita risplenderà.
\endchorus
\beginverse
E tu ci guiderai sulla via che porta a te;
donerai salvezza all'umanità.
\endverse
\beginverse
E tu accoglierai l'uomo che ritorna a te;
una grande festa tu farai per lui.
\endverse
\beginverse
E tu illuminerai ogni nostra oscurità;
splenderà il tuo volto dentro di noi.
\endverse
\beginverse
E tu ridonerai la tua pace all'umanità;
e verrà il tuo Regno in mezzo a noi.
\endverse
\endsong

% Fonte: LibrettoAccordi.pdf (Canti Coro), O12.
% Nell'elenco: Offertorio, Settimana Santa.
\beginsong{Dayenu Adonai}[by={Gen Verde (CD Il mistero pasquale)}]
\beginverse
Se dopo averci fatti uscire dall'Egitto
Tu non ci avessi sostenuti con la manna,
dayenu, dayenu Adonai,
dayenu, dayenu Adonai.
Se dopo averci sostenuti con la manna
Tu non ci avessi consegnato la tua Legge,
dayenu, dayenu Adonai,
dayenu, dayenu Adonai.
Se dopo averci consegnato la tua Legge
Tu non ci avessi fatti entrare in Israele,
dayenu, dayenu Adonai,
ci sarebbe bastato, Signore.
\endverse
\beginverse
Se dopo averci definito tuoi amici
Tu non ti fossi inginocchiato a lavarci,
dayenu, dayenu Adonai,
dayenu, dayenu Adonai.
Se quando Tu ti sei chinato a lavarci
poi non ti fossi a noi spezzato come pane,
dayenu, dayenu Adonai,
dayenu, dayenu Adonai.
Se quando Tu ti sei spezzato come pane
poi non ti fossi consegnato all'abbandono
dayenu, dayenu Adonai,
ci sarebbe bastato, Signore.
\endverse
\beginverse
Per tutto questo cosa mai potrò donarti?
Se ti rendessi oro e perle senza pari,
non potrà mai bastare, Adonai,
o Signore, non basterà mai.
Se ti rendessi oro e perle senza pari
ed aggiungessi il coro immenso dei respiri,
non potrà mai bastare, Adonai,
o Signore, non basterà mai.
Se ti rendessi il coro immenso dei respiri
ed aggiungessi terre, cieli ed universi,
non potrà mai bastare, Adonai,
o Signore non basterà mai.
\endverse
\beginverse
Noi ti portiamo tralci spogli e solchi vuoti,
Tu ci darai stille d'uva e farina;
solo Tu, solo Tu Adonai,
solo Tu, solo Tu basterai.
Noi ti portiamo stille d'uva e farina,
Tu ne farai il tuo corpo, o Signore;
solo Tu, solo Tu Adonai,
solo Tu, solo Tu basterai.
Noi ti portiamo mille vite e mille cuori,
Tu ci farai un solo corpo, un solo cuore;
solo Tu, solo Tu Adonai,
solo Tu, solo Tu basterai.
Solo Tu, solo Tu basterai.
\endverse
\endsong

% Fonte: LibrettoAccordi.pdf (Canti Coro), J2.
% Nell'elenco: Comunione, Pasqua.
\beginsong{E sei rimasto qui}[by={Gen Rosso}]
\beginverse
Perché la sete d'infinito? Perché la fame d'immortalità?
Sei tu che hai messo dentro il cuore il desiderio dell'eternità.
Ma Tu sapevi che quel vuoto lo colmavi tu,
per questo sei venuto in mezzo a noi.
\endverse
\beginchorus
E sei rimasto qui visibile mistero.
E sei rimasto qui cuore del mondo intero.
E rimarrai con noi finché questo universo girerà, salvezza dell'umanità!
\endchorus
\beginverse
Si apre il cielo del futuro, il muro della morte ormai non c'è
Tu, Pane Vivo, ci fai Uno, richiami tutti i figli attorno a te.
E doni il Tuo Spirito che lascia dentro noi
il germe della sua immortalità.
\endverse
\beginverse
Presenza vera nel mistero, ma più reale di ogni realtà.
Da te ogni cosa prende vita, e tutto un giorno a te ritornerà.
Varcando l'infinito tutti troveremo in Te
un Sole immenso di felicità.
\endverse
\endsong

% Fonte: LibrettoAccordi.pdf (Canti Coro), O17.
% Nell'elenco: Ingresso, Settimana Santa.
\beginsong{È venuta l'ora}[by={Gen Verde (CD Il mistero pasquale)}]
\beginverse
È venuta l'ora; c'è pianto e festa stasera,
ci raccogli insieme con te e ti consegni a noi, Signore.
\endverse
\beginchorus
Kyrie eleison, Christe eleison, Kyrie eleison, eleison.
\endchorus
\beginverse
È venuta l'ora; tu, come servo, Signore
versi un'acqua di libertà su piedi rossi di catene.
\endverse
\beginverse
È venuta l'ora d'amore senza misura;
tutto il tuo mistero ci dai, che sa di cielo e sa di grano.
\endverse
\beginverse
La tua gloria ha il colore delle piaghe,
la tua croce di vittoria ci rivestirà.
\endverse
\endsong

% Fonte: LibrettoAccordi.pdf (Canti Coro), K1.
% Nell'elenco: Finale, Pasqua.
\beginsong{Esulta il cielo}[by={F. Buttazzo}]
\beginchorus
Esulta il cielo e canta a noi la novità:
Gesù è risorto dalla morte!
Esulta l'uomo per la vita che gli donerà
Gesù, il Signore.
Esulta il cielo e canta a noi la novità:
Gesù è risorto dalla morte!
E si spalancano le porte dell'eternità
per tutti noi.
\endchorus
\beginverse
Il vero Agnello che toglie i peccati è lui,
lui che morendo ha distrutto la morte,
e risorgendo ha ridato la vita,
vita per l'eternità.
\endverse
\beginverse
In lui risplende la vita del Padre,
a noi è donata la vita dei figli.
E sono aperte le porte del Regno
per chi crede nel suo nome.
\endverse
\beginverse
In lui, che ha vinto il peccato e la morte,
l'universo si rinnova e risorge;
l'uomo ritorna alle vere sorgenti
della vita insieme a Dio.
\endverse
\beginverse
Ora continua ad offrirsi per tutti noi,
lui ci difende, è nostro «avvocato».
L'uomo che vive amando i fratelli
la salvezza abbraccerà.
\endverse
\endsong

% Fonte: LibrettoAccordi.pdf (Canti Coro), O5.
% Nell'elenco: Settimana Santa.
\beginsong{Il Signore è la luce}[by={M. Giombini}]
\beginverse
Il Signore è la luce che vince la notte.
Gloria, gloria! Cantiamo al Signore!
\endverse
\beginverse
Il Signore è il coraggio che vince il terrore.
Gloria, gloria! Cantiamo al Signore!
\endverse
\beginverse
Il Signore è il sereno che vince la pioggia.
Gloria, gloria! Cantiamo al Signore!
\endverse
\beginverse
Il Signore è la vita che vince la morte.
Gloria, gloria! Cantiamo al Signore!
\endverse
\beginverse
Il Signore è l'amore che vince il peccato.
Gloria, gloria! Cantiamo al Signore!
\endverse
\beginverse
Il Signore è la gioia che vince l'angoscia.
Gloria, gloria! Cantiamo al Signore!
\endverse
\beginverse
Il Signore è la pace che vince la guerra.
Gloria, gloria! Cantiamo al Signore!
\endverse
\endsong

% Fonte: LibrettoAccordi.pdf (Canti Coro), O10.
% Nell'elenco: Lavanda, Settimana Santa.
\beginsong{Io vi do un grande esempio}[by={D. Haas}]
\beginchorus
Io vi do un grande esempio, con amore lavo i vostri piedi,
se ora tutti voi farete come me l'amore crescerà.
\endchorus
\beginverse
Ora, Pietro, il mio gesto non capisci,
un giorno non lontano capirai
cos'è l'amore.
\endverse
\beginverse
Ti stupisce quello che io sto facendo,
io mi chino per lavare i tuoi piedi: questo è l'amore.
\endverse
\beginverse
Guarda e impara da me.
Chi lava i piedi all'altro capirà
cos'è l'amore.
\endverse
\beginverse
Questo è l'amore più grande:
dare con gioia la tua vita
per un amico.
\endverse
\beginverse
Vai e vivi nell'amore;
tu sarai mio discepolo fedele
se amerai.
\endverse
\beginverse
Sono io che ho scelto te.
Io ti mando ad annunciare il mio regno
ai miei fratelli.
\endverse
\endsong

% Fonte: "99 - Raccolte e libretti accordi/Il mistero pasquale_ gv.pdf" (pag. 69).
% Nell'elenco: Adorazione della Croce, Settimana Santa. CD: Il mistero pasquale (2011). In Quaresima e Pasqua su indicazione dell'utente (non tra i Mariani).
\beginsong{Lì con te}[by={Gen Verde; testo P. Stradi (CD Il mistero pasquale)}]
\beginverse
Figlio mio, tu sei qui;
sanguina ancora il tuo perché.
Scende la spada nel cuore,
scende sul mondo la notte;
qui sul mio grembo,
come agnello svenato sei qui.
\endverse
\beginverse
Sono qui, qui con te.
Niente può farti male, ormai;
non urla più l'uragano,
tutto lo strazio è lontano;
qui sul mio grembo,
per un'ultima volta sei qui.
\endverse
\beginchorus
Ed ora ti porteranno via
strappandoti da me,
ma qui sul mio grembo resterà
la tua impronta, Figlio,
ed il sangue tuo su di me.
\endchorus
\beginverse
Ora addio, Figlio mio.
L'ultimo bacio che ti do
ti porta tutto l'amore
di chi non trova parole
per dirti grazie
e per chiedere la tua pietà.
\endverse
\beginchorus
Ed ora ti porteranno via
strappandoti da me,
ma Figlio, non t'abbandonerò
e davanti al sepolcro sbarrato, lì io starò.
\endchorus
\beginchorus
Ed ora ti porteranno via;
in piedi aspetterò
finché da quel grembo nascerai
e per sempre vivo mi porterai
lì con te.
\endchorus
\endsong

% Fonte: LibrettoAccordi.pdf (Canti Coro), O11.
% Nell'elenco: Ingresso, Settimana Santa.
\beginsong{L'ultima cena con i suoi}[by={F. Buttazzo}]
\beginverse
È vicina la Pasqua, festa degli Ebrei,
e Gesù siede a cena coi discepoli suoi.
Raccontava la storia dell'Egitto e di Mosè,
e, spezzando il pane, disse loro così.
\endverse
\beginchorus
Prendete e mangiatene, amici miei: il pane è il mio corpo, donato a voi.
Bevete il mio calice e fate questo in memoria di me.
\endchorus
\beginverse
Quella sera si apriva il cuore di Gesù,
fece un dono speciale ai discepoli suoi:
anche loro saranno sacerdoti come Gesù,
sull'altare diranno ciò che ha detto Gesù.
\endverse
\beginverse
In quell'ultima cena con i Dodici
Lui veniva tradito da un discepolo.
Con immenso dolore anche a Giuda il pane donò
e con vero amore disse loro così.
\endverse
\endsong

% Fonte: "99 - Raccolte e libretti accordi/Il mistero pasquale_ gv.pdf" (pag. 20).
% Nell'elenco: Lavanda, Settimana Santa. CD: Il mistero pasquale (2011).
\beginsong{Non c'è amore più grande}[by={Gen Verde (CD Il mistero pasquale)}]
\beginchorus
Non c'è amore più grande di chi dà la vita per i suoi.
Non c'è amore più grande: io do la mia vita per voi.
\endchorus
\beginverse
Il Signore si alzò da tavola,
versò dell'acqua in un catino,
e cominciò a lavare i piedi ai discepoli:
ad essi volle lasciare questo esempio.
\endverse
\beginverse
«Quello che io faccio, ora non lo comprendi,
ma lo comprenderai un giorno».
«Signore, tu lavi i piedi a me?».
Gesù gli rispose dicendo:
«Se non ti laverò, non avrai parte con me».
\endverse
\beginverse
«Se vi ho lavato i piedi,
io, Signore e Maestro,
quanto più voi avete il dovere
di lavarvi i piedi l'un l'altro».
\endverse
\beginverse
Gesù disse ai suoi discepoli:
«Da questo tutti sapranno
che siete miei discepoli,
se vi amerete gli uni gli altri».
\endverse
\beginverse
«Vi do un comandamento nuovo:
che vi amiate gli uni gli altri
come io ho amato voi»,
dice il Signore.
\endverse
\beginverse
Fede, speranza e carità,
tutte e tre rimangano tra voi:
ma più grande di tutte è la carità.
\endverse
\beginchorus
Non c'è amore più grande di chi dà la vita per i suoi.
Non c'è amore più grande: io do la mia vita per voi, per voi.
\endchorus
\endsong

% Fonte: LibrettoAccordi.pdf (Canti Coro), O1.
% Nell'elenco: Ingresso, Settimana Santa.
\beginsong{Osanna al Figlio di David}[by={Canto popolare}]
\beginchorus
Osanna al figlio di David! Osanna al Redentor. \rep{2}
\endchorus
\beginverse
Apritevi o porte eterne, avanzi il Re della gloria;
adorin cielo e terra l'eterno suo poter.
\endverse
\beginverse
O monti stillate dolcezza, il Re d'amor s'avvicina;
si dona pane vivo ed offre pace al cuor.
\endverse
\beginverse
O Vergine presso l'Altissimo trovasti grazia e onor;
soccorri i tuoi figliuoli donando il Salvator.
\endverse
\beginverse
Verrai un giorno Giudice, o mite e buon Gesù;
rimetti i nostri falli nel tempo del perdon.
\endverse
\endsong

% Fonte: testo fornito dall'utente. Non presente nell'elenco. Autore indicato dall'utente.
% Domenica delle Palme. Il ritornello all'inizio e alla fine si canta due volte (qui scritto una volta).
\beginsong{Rallegrati, Gerusalemme}[by={P. Dargenio}]
\beginchorus
Rallegrati, Gerusalemme,
esulta tra le tue mura,
perché il Signore, adesso, viene da te
tra petali e mantelli, tra canti di gioia:
Osanna, osanna al Re!
\endchorus
\beginverse
Le folle degli ebrei cantavano a gran voce:
Osanna, osanna al Re!
Ti festeggiavano con palme nelle mani:
Osanna, osanna al Re!
\endverse
\beginverse
Le folle degli ebrei non possono tacere:
Osanna, osanna al Re!
Ma se tacessero le pietre griderebbero:
Osanna, osanna al Re!
\endverse
\beginverse
Le folle degli ebrei stendevano i mantelli:
Osanna, osanna al Re!
E ti acclamavano Messia d'Israele:
Osanna, osanna al Re!
\endverse
\beginverse*
Osanna, osanna al Re!
Osanna, osanna al Re!
\endverse
\endsong

% Fonte: LibrettoAccordi.pdf (Canti Coro), K3 (li' tra i canti finali).
% Nell'elenco: Finale, Pasqua. Omesso il finale "a cori scontrati".
\beginsong{Resurrezione}[by={Gen Rosso}]
\beginverse
Che gioia ci hai dato, Signore del cielo, Signore del grande universo!
Che gioia ci hai dato, vestito di luce, vestito di gloria infinita,
vestito di gloria infinita!
\endverse
\beginverse
Vederti risorto, vederti Signore, il cuore sta per impazzire!
Tu sei ritornato, tu sei qui tra noi, e adesso ti avremo per sempre,
e adesso ti avremo per sempre.
\endverse
\beginverse
Chi cercate, donne, quaggiù? Chi cercate, donne, quaggiù?
Quello che era morto non è qui:
è risorto, sì! Come aveva detto anche a voi, voi gridate a tutti che
è risorto Lui, a tutti che è risorto Lui!
\endverse
\beginverse
Tu hai vinto il mondo, Gesù! Tu hai vinto il mondo, Gesù!
Liberiamo la felicità!
E la morte, no, non esiste più, l'hai vinta tu e hai salvato tutti noi,
uomini con te, tutti noi, uomini con te.
\endverse
\endsong

% Fonte: LibrettoAccordi.pdf (Canti Coro), O8. C'e' anche "6 - Offertorio/SERVIRE è REGNARE.pdf".
% Nell'elenco: Lavanda, Settimana Santa. CD: Il mistero pasquale (2011).
\beginsong{Servire è regnare}[by={Gen Verde (CD Il mistero pasquale)}]
\beginverse
Guardiamo a te che sei Maestro e Signore:
chinato a terra stai, ci mostri che l'amore
è cingersi il grembiule, sapersi inginocchiare,
ci insegni che amare
è servire.
\endverse
\beginchorus
Fa' che impariamo, Signore, da te,
che il più grande è chi più sa servire,
chi si abbassa e chi si sa piegare,
perché grande è soltanto l'amore.
\endchorus
\beginverse
E ti vediamo poi, Maestro e Signore,
che lavi i piedi a noi che siamo tue creature;
e, cinto del grembiule, che è il manto tuo regale,
ci insegni che servire
è regnare.
\endverse
\endsong

% Fonte: LibrettoAccordi.pdf (Canti Coro), O6 (sezione Settimana Santa - Ceneri).
% Nell'elenco: senza tipologia/periodo.
\beginsong{Ti saluto, o croce santa}[by={Gazzera, Damilano}]
\beginchorus
Ti saluto, o croce santa, che portasti il Redentor:
gloria, lode, onor ti canta ogni lingua ed ogni cuor.
\endchorus
\beginverse
Sei vessillo glorioso di Cristo, sua vittoria e segno d'amor:
il suo sangue innocente fu visto come fiamma sgorgare dal cuor.
\endverse
\beginverse
Tu nascesti fra braccia amorose d'una Vergine madre, o Gesù;
tu moristi fra braccia pietose d'una croce che data ti fu.
\endverse
\beginverse
O Agnello divino, immolato sulla croce crudele, pietà!
Tu, che togli dal mondo il peccato, salva l'uomo che pace non ha.
\endverse
\beginverse
Dona a tutti speranza, Signore, crocifisso e risorto per noi:
tu che effondi la pace del cuore nel tuo Spirito di santità.
\endverse
\beginverse
Del giudizio nel giorno tremendo sulle nubi del cielo verrai.
Piangeranno le genti, vedendo qual trofeo di gloria sarai.
\endverse
\endsong

% Fonte: testo fornito dall'utente.
% Nell'elenco: Finale, Pasqua. CD: Il mistero pasquale (2011).
\beginsong{Tutto il mondo deve sapere}[by={Gen Verde (CD Il mistero pasquale)}]
\beginchorus
Alleluia, alleluia, alleluia. \rep{2}
\endchorus
\beginverse
Come potrò raccontare?
È una gioia che fa piangere e fa gridare.
Io l'ho visto con i miei occhi,
era vivo, era Lui.
E m'ha chiamata per nome,
era la sua voce, era il mio Signore.
Io l'ho visto con i miei occhi,
era lì davanti a me.
\endverse
\beginverse
Gli angeli ci hanno parlato
davanti a quel sepolcro spalancato:
«Donne, il Signore è risorto,
non cercatelo qui».
Che corsa senza respiro
per gridare a tutti di quel mattino.
Tutto il mondo deve sapere
che è rimasto qui tra noi.
\endverse
\beginverse
Alba di un tempo diverso,
è il mattino dei mattini per l'universo.
Tutto già profuma d'eterno,
c'è il Risorto fra noi.
E le sue piaghe e la gloria
sono vive dentro la nostra storia,
segni di un amore che resta
qui per sempre, qui con noi.
\endverse
\beginverse*
Alleluia.
\endverse
\endsong

% Fonte: LibrettoAccordi.pdf (Canti Coro), O9 (sezione Settimana Santa - Ceneri).
% Nell'elenco: senza tipologia/periodo.
\beginsong{Venite dal profondo}[by={R. J. Stamps}]
\beginverse
Venite dal profondo dei tempi, o cuori dalla sete piagati.
Aprite il cuore alla gioia più profonda.
\endverse
\beginchorus
Dio ci ha messo il suo Corpo tra le mani. \rep{3}
\endchorus
\beginverse
Stasera l'acqua si trasforma in vino, sul monte il pane è moltiplicato.
Precoce è il frutto della vigna in fiore.
\endverse
\beginverse
Come ferita aperta è la parola. Parola di un eterno senza fine,
il Verbo si fa carne e nutrimento.
\endverse
\endsong

\end{songs}

% !TEX root = ../libretto.tex
% Sezione: Altri canti
\momento{Altri canti}
\begin{songs}{titoli}

% Fonte: testo fornito dall'utente (tolti gli accordi; corretti i caratteri rovinati: tutti, restiamo, gettiamo, attimo).
% Nell'elenco: "Beato il cuore che perdona". Inno della GMG di Cracovia 2016 (versione italiana).
\beginsong{Beato il cuore che perdona}[by={J. Blycharz}]
\beginverse
Sei sceso dalla tua immensità
in nostro aiuto.
Misericordia scorre da te
sopra tutti noi.
Persi in un mondo d'oscurità
lì tu ci trovi.
Nelle tue braccia ci stringi e poi
dai la vita per noi.
\endverse
\beginchorus
Beato è il cuore che perdona!
Misericordia riceverà da Dio in cielo!
\endchorus
\beginverse
Solo il perdono riporterà
pace nel mondo.
Solo il perdono ci svelerà
come figli tuoi.
\endverse
\beginverse
Col sangue in croce hai pagato tu
le nostre povertà.
Se noi ci amiamo e restiamo in te
il mondo crederà!
\endverse
\beginverse*
Le nostre angosce ed ansietà
gettiamo ogni attimo in te.
Amore che non abbandona mai,
vivi in mezzo a noi!
\endverse
\endsong

% Fonte: testo fornito dall'utente (tolti i vocalizzi "oh oh"). Autore non indicato.
\beginsong{Cerco il tuo volto}
\beginverse
Contemplo in silenzio
un quadro scomposto,
ricerco il mio volto,
desidero te.
Mi incontri spandendo,
sprecando l'amore,
consegni la vita
umiliato per me.
Accetta, Signore, le lacrime amare
di chi ha rinnegato e venduto l'amore.
\endverse
\beginchorus
Io cerco il tuo volto,
il tuo sguardo d'amore
per chi ha rinnegato
e venduto l'amore.
Cerco il tuo volto, cerco il tuo volto,
cerco il tuo volto,
il tuo volto, il tuo volto, Signore.
\endchorus
\beginverse
Tra polvere e rabbia
che ingombra il mio cuore
riveli il tuo volto
di Figlio che muore.
Quando un gelido masso
serra forte l'amore,
io ti riconosco
Maestro e Signore.
Accetta, Signore, il grido stupito
di chi t'ha incontrato e ti vuole annunciare.
\endverse
\beginverse
Ho trovato il tuo volto,
il tuo sguardo soave,
la segreta dolcezza
che mi vuol perdonare.
Ho trovato il tuo volto,
il tuo sguardo soave:
ora vado nel mondo
ed insegno a pregare.
\endverse
\beginchorus
Cerco il tuo volto, cerco il tuo volto,
cerco il tuo volto,
il tuo volto, il tuo volto, Signore.
\endchorus
\endsong

% Fonte: "Con-Te-Camminero-Official-Lead-Sheet.pdf" (melodia con accordi).
% Nell'elenco: senza tipologia/periodo.
\beginsong{Con te camminerò}[by={Gen Verde}]
\beginverse
Mi hai chiamato, ora sono qui.
La mia vita ormai, Signore, trova senso solo in te.
Cerco nel profondo e tu ci sei.
Guardo intorno e poi ti trovo accanto a me.
\endverse
\beginchorus
Insieme a te camminerò.
In ogni volto, in ogni pianto, io ti riconoscerò.
Ti seguirò ovunque andrai.
Io con te camminerò.
\endchorus
\beginverse
Come servo sei venuto qui,
per chi è solo e chi è ferito, per chi vive in povertà.
Ora, o Signore, manda me,
segno del tuo amore per l'umanità.
\endverse
\beginverse*
Tu sei qui, sei con noi, sempre accanto a noi.
Nel dolore di chi non ha più una casa, tu sei lì.
Nel silenzio di chi ha perso ogni speranza, tu sei lì.
Nei rimpianti, nei miei dubbi
e nel buio delle mie fragilità, tu sei lì.
\endverse
\beginverse*
And I will walk on holy ground.
I'll hear your voice, I'll see your face, I will find you in the crowd.
Lord, I will go where you will lead me.
I will walk on holy ground.
\endverse
\beginverse*
Et avec toi, j'irai, Seigneur.
J'entends ta voix, et je te vois, dans les pauvres autour de moi.
Je te suivrai où tu me conduis.
Je te vois auprès de moi.
\endverse
\beginverse*
You are near, you are here.
Io con te camminerò.
\endverse
\endsong

% Fonte: LibrettoAccordi.pdf (Canti Coro), F1.
% Nell'elenco: Credo, Ordinario. In Altri canti su indicazione dell'utente (unico Credo, sezione Credo eliminata).
\beginsong{Credo in unum Deum}[by={M. Balduzzi (CD Verbum Panis)}]
\beginverse
Credo in unum Deum Patrem omnipotentem.
Credo in unum Deum factorem coeli et terrae,
visibilium omnium et invisibilium
Credo in unum Deum. Amen.
\endverse
\beginverse
Credo in unum Deum Dominum Jesum Christum
Credo in unum Deum Filium Dei unigenitum
et ex Patre natum ante omnia saecula
Credo in unum Deum. Amen.
\endverse
\beginverse
Credo in unum Deum Spiritum Sanctum
Credo in unum Deum Dominum et vivificantem
qui ex Patre Filioque procedit
Credo in unum Deum. Amen.
\endverse
\beginverse
Credo in unum Deum Patrem omnipotentem
Credo in unum Deum Dominum Jesum Christum
Credo in unum Deum Spiritum Sanctum
Credo in unum Deum. Amen. Amen.
\endverse
\endsong

% Fonte: testo fornito dall'utente.
% Nell'elenco: "Incendio", Fraternità & Serena Lambertucci, Adorazione; l'utente lo mette in Altri canti.
\beginsong{Incendio}[by={Fraternità, S. Lambertucci}]
\beginverse
Adesso come fuoco partiremo per il mondo
e canteremo la parola d'ogni uomo:
fraternità.
Eccoci, manda noi come luce nell'oscurità,
porteremo salvezza in ogni estremità
della terra.
\endverse
\beginchorus
Tu ci hai scelto dal mondo perché portiamo frutto,
come fosse un incendio:
tu sarai insieme a noi in ogni attimo per farci tuoi. \rep{2}
Per farci tuoi.
\endchorus
\beginverse
Adesso come membra del tuo corpo cureremo
ogni cuore che ha bisogno di conforto:
fraternità.
L'amore che tu ci hai donato ci guiderà,
non sarà più da sola ogni singola anima
in eterno.
\endverse
\beginverse*
Non temiamo perché tu ci hai chiamato
e lo Spirito arde in noi. \rep{4}
Arde in noi.
\endverse
\beginchorus
Tu ci hai scelto dal mondo perché portiamo frutto,
come fosse un incendio:
tu sarai insieme a noi in ogni attimo per farci tuoi. \rep{3}
Per farci tuoi.
\endchorus
\endsong

% Fonte: LibrettoAccordi.pdf (Canti Coro), J17.
% Nell'elenco: senza tipologia (nel libretto accordi tra i canti di comunione). L'utente lo mette in Altri canti.
\beginsong{Nascerà}[by={Gen Rosso}]
\beginverse
Non c'è al mondo chi mi ami, non c'è stato mai nessuno,
in fondo alla mia vita, come Te.
È con te la mia partita, come sabbia tra le dita
scorrono i miei giorni insieme a Te.
Inquietudine o malinconia:
non c'è posto per loro in casa mia.
Sempre nuovo è il tuo modo di
inventare il gioco del tempo per me.
\endverse
\beginchorus
Nascerà dentro me, sul silenzio che abita qui.
Fiorirà un canto che mai nessuno ha cantato per Te.
\endchorus
\beginverse
Se la strada si fa dura come posso aver paura?
Nel buio della notte ci sei Tu.
Se mi assale la fatica di cancellare la sconfitta,
dietro ogni ferita sei ancora Tu.
È una cosa che non mi spiego mai,
cosa ho fatto perché tu scegliessi me?
Cosa mai dirò quando mi vedrai,
quando dai confini del mondo verrai?
\endverse
\endsong

% Fonte: LibrettoAccordi.pdf (Canti Coro), I3 "Segno della pace".
% Nell'elenco: "Pace", Balduzzi, Pace, Ordinario. CD: Giorno senza tramonto (2013).
% Omessa la ripresa finale a due cori (le due parti sovrapposte).
\beginsong{Pace}[by={M. Balduzzi (CD Giorno senza tramonto)}]
\beginchorus
Pace! Pace! Segno della pace che hai portato agli uomini.
Pace! Pace! Pace che qui in terra hai dato a tutti gli uomini.
\endchorus
\beginverse*
Dove c'è odio porterò pace, dov'è tempesta il sole,
dov'è silenzio sarò una voce, dov'è deserto una fonte.
In mezzo al gelo porterò fuoco, in mezzo a spine una rosa,
dove c'è noia porterò vita, dov'è tristezza un sorriso.
\endverse
\endsong

\end{songs}


% ---- Indice alfabetico -------------------------------------------------
\clearpage
\phantomsection\addcontentsline{toc}{section}{Indice alfabetico}
\nellindicetrue
\showindex[2]{Indice alfabetico}{titoli}

\end{document}
